% The neUROn design manifest!
\documentclass{report}
\usepackage{color}
\usepackage{makeidx}
\usepackage{graphicx}
\usepackage[dvipdfm]{hyperref} % this should be the last package

% Colors
\definecolor{myLinkColor}{rgb}{0.5, 0.0, 0.7}
\def\beginpurple{{\special{pdf:bc \purple}}}
\def\purple{[0.2 0.0 0.45]}
\def\beginred{{\special{pdf:bc \red}}}
\def\red{[0.8 0.0 0.0]}
\def\begincolor#1{\special{pdf:bc #1}}
\def\endcolor{{\special{pdf:ec}}}

% hyperref configuration
\hypersetup{colorlinks=true,
linkcolor=myLinkColor}

\setlength{\unitlength}{1in}
\newcounter{tabcount}
\setcounter{tabcount}{1}
\newcounter{figcount}
\setcounter{figcount}{1}

% Math from maths.tex
\newcommand{\D}[2]{\frac{\partial #1}{\partial #2}}

\title{neUROn Design Manifest}
\author{C. Niederberger\\and\\the neUROn group}

\makeindex

\begin{document}

\maketitle

\chapter{Introduction and History}

The original neUROn\index{History!neUROn} was coded in C in 1992 when
I was at Baylor College of Medicine.  It was named
``\underline{ne}ural computational for \underline{URO}logical
\underline{n}umericals,'' a somewhat mercurial reference to its
application domain, the medical field of Urology.  In actuality,
neUROn is a general purpose neural computational modeling environment.
neUROn has been, and will continue to be used to model many problems
beyond the Urological domain.

Between 1992 and 2000, neUROn evolved in fundamental ways.  In 1996,
Young Hong migrated neUROn to C++, and the resulting environment was
named ``neUROn++''.\index{History!neUROn++}  Many have coded these
fundamental changes, including Vinod Kutty, Yuan Qin, Luke Cho, Joe
Jovero, and Kai Liu.  In general, beyond migrating neUROn to
object-oriented programming, these fundamental changes have taken 2
forms: programming novel neural computational training algorithms and
programming statistical analysis of training behavior.  Throughout
this project, Richard Golden Ph.D.\index{Golden, Richard} at the
University of Texas at Dallas has suggested algorithms which the team
has implemented.

The impetus for programming novel training algorithms should be
obvious.  Neural computational algorithms are iterative, and problems
modeled by the neural computational approach can easily become
intractable, even with modern computational hardware.  A single
training session may require several hours, and multiple training
sessions may require days or even weeks.  There are 3 basic ways in
which training sessions may become shorter:

\begin{enumerate}
\item Faster hardware may be used.  This is clearly the way in which
      the programmer has the least control.
\item More efficient code may be written.  The importance of this
      cannot be stessed enough, and it will serve as a basic design
      principle for the neUROn project.  In the original neUROn code,
      for example, arrays were expressed entirely as pointer
      operations.  It was not until the compiler benchmarks
      demonstrated similar outcomes for array constructs that arrays
      were used in the code.
\item Faster training algorithms may be encoded.  A substantial part
      of the neUROn project is devoted to investigating novel training
      algorithms, and this has paid handsomely.  Investigated
      algorithms have been demonstrated to speed training by 2 orders
      of magnitude or more.
\end{enumerate}

The impetus for programming statistical analysis of network behavior
was less obvious, but perhaps more significant.  When we began
reporting results to the Urological community, we were surprised to
find a universal response was, ``so the network accurately predicts
$X$, but what were the important variables?''  In retrospect, this
response is unsurprising.  The majority of our audience, our potential
users, were physicians or scientists involved in the medical field.
It is not enough to accurately model an outcome for these users, it is
critical to find the significant input features so that these features
may be altered by the user in the clinical environment.  For example,
if it is determined that medication $X$ is statistically significant
in the model $Y$ that predicts cancer recurrence, the user will put
that information to use in the clinical domain.

Beginning in 2000, neUROn was entirely redesigned using this document
as a specification, and taking full advantage of the object-oriented
nature of C++.  Whereas the migration from 1996 to 2000 was an
incorporation of C code into the C++ paradigm, the hybrid code which
resulted was brittle, with previous features breaking as new features
were added.  In 2000, the code was redesigned from the ``ground up''.
This document is intended to be a description of neUROn programming,
past, present, and future.  As such, the document is a plan that is
expected to evolve as neUROn is redesigned and recoded in the future.

\chapter{Methodology}\index{Golden, Richard|(}

The following chapter outlines algorithms currently implemented in
neUROn, as well as many intended to be implemented in the future. As
usual, we are indebted to Dr. Golden. The terminology ANN is used to
refer to ``Artificial Neural Network''.

\begin{enumerate}

\begin{item}\index{Forward propagation}
{\bf Error Gradient Calculation Algorithm.}  The gradient of the
objective function is computed by calling a {\em forward propagation
algorithm} followed by a {\em backward error propagation algorithm}.
\end{item}

\begin{item}
{\bf Parameter Estimation Algorithm.}  The gradient computed by the
{\em Error Gradient Calculation Algorithm} is then used by the
parameter estimation algorithm for the purposes of learning.
\end{item}

\begin{item}
{\bf Data Analysis Algorithm.}  If this algorithm decides that
reliable statistical inference is possible, this algorithm indicates
which input nodes in the network are making statistically significant
contributions to the network's performance.  Otherwise, this algorithm
makes recommendations to the user regarding possible manipulations
designed to increase statistical inference reliability. Future work
will involve computing confidence intervals on the predictions and
relaxing the ``Goodness-of-Fit'' assumption which is assumed
throughout this document.
\end{item}
\end{enumerate}

Note that math calculations for the following algorithms such as
matrix inverse, eigenvalues, ``line-search'', cumulative distribution
functions for the chi-square random variable for statistical test
result reporting can be done using the routines in {\em Numerical
Recipes in C}.\index{Numerical Recipes in C}  In fact,
{\em Numerical Recipes} has a conjugate gradient routine which could
be used to develop a learning algorithm for the ANN system as well.

\section{Error Gradient Calculation Algorithm}

\subsection{Data Records}
Note that the data is assumed to have the following format: $({\bf
s}_k, {\bf t}_k)$ where ${\bf s}_k$ is the $d$-dimensional input
vector for the $k$th data record and $d$-dimensional vector ${\bf
t}_k$ is a list of ``targets'' for the $d$ output units. If the
training data consists of $N$ data records, then we have $k = 1,
\ldots, N$.

\subsection{Objective Function}\label{objFunc}
\index{Error function|(}

The {\em error gradient} is defined as the gradient (i.e., vector
first derivative) of a sufficiently smooth {\em error function} with
respect to the ``weights'' (i.e., parameters) of the ANN. In this
document version, only feedforward $m$-layer ANN systems will be
considered where the number of ``layers'' is defined by
``connections''. Thus, a $1$-layer ANN is essentially a linear
feedforward ANN with one set of input units and one set of output
units and no hidden units.  Let ${\bf y}$ denote a vector containing
all weights of an $m$-layer ANN. The activation levels of the units in
the output layer of an $m$-layer ANN generated in response to training
stimulus $k$ are denoted by the {\em activation vector}, ${\bf
a}^{(m)}_k$.  Note that ${\bf a}^{(m)}_k$ is a function of both ${\bf
y}$ and the input vector ${\bf s}_k$.

The {\em error function}, $E_k$, for the $k$th training pattern is
defined by the formula:
\begin{equation}\label{deferrorfunk}
E_k({\bf y}) =
e(t^k, {\bf a}^{(m)}_k) +
\frac{1}{2\sigma_w^2} \sum_{i=1}^m |{\bf y}|^2.
\end{equation}

The second term on the right hand side of (\ref{deferrorfunk}) is
required to guarantee reliable statistical inference.
\index{Weight decay} It basically ``stabilizes'' the learning process
by preventing the weights in the network from becoming too large or
too small.  In practice, a connection weight is expected to lie
between $-1.96 \sigma_w$ and $+1.96 \sigma_w$ with probability equal
to about 95\%. Adjustment of this parameter can ``change results'' of
the statistical analysis so it is best to fix $\sigma_w$ to be some
constant number and try not to change its value. A reasonable choice
for $\sigma_w$ would be $\sigma_w = 100$. This basically constrains
the ``learning process'' to find connection strengths for the ANN such
that the connection strength weight ranges between $-200$ and $+200$.
The choice of $\sigma_w$ should be reported in conjunction with all
statistical analysis results. Note that as
$\sigma_w \rightarrow \infty$, the effect of the summation term
becomes negligible.

The first term on the right hand side of (\ref{deferrorfunk}) is
called the {\em prediction error function} term. The prediction error
function $e$ must have continuous second partial derivatives with
respect to the weights.  For example, the classic choice of the
prediction error function $e$ is to choose
\index{Error function!Least squares}
\begin{equation}
\label{leastsquaresfunk}
e({\bf t}^k, {\bf a}^{(m)}_k) = (1/2)|{\bf t}^k - {\bf a}^{(m)}_k|^2
\end{equation}
to be the {\em least squares prediction error function}.  This is a
good choice when the target vector ${\bf t}^k$ is a continuous (as
opposed to a discrete) variable.

When the target vector consists of a set of binary variables taking on
the values of zero or one, then a good choice is the
{\em cross-entropy prediction error function} defined by the
formula:\index{Error function!Cross-entropy}
\begin{equation}\label{xentropyfunk}
e({\bf t}^k, {\bf a}^{(m)}) =
-{\bf t}^k .* {\bf log} [{\bf a}^{(m)}_k] - ({\bf 1} - {\bf t}^k) .*
{\bf log} [1 - {\bf a}^{(m)}_k].
\end{equation}
Note that the notation $.*$ denotes element by element matrix
multiplication, the notation ${\bf 1}$ denotes a column vector of
ones, and the notation ${\bf log}$ denotes a vector-valued function
defined such that the ${\bf y} = {\bf log}({\bf x})$ whenever the
$i$th element of the vector ${\bf y}$, $y_i$ is equal to the natural
log (log base $e$) of the $i$th element of the vector ${\bf x}$. When
the cross-entropy prediction error function is used, the transfer
function for all output units should be the sigmoidal logistic
transfer function which is defined formally in the next section of
this report.\index{Transfer function!Sigmoidal}
\index{Error function|)}

\subsection{Forward Propagation of Activations}
\index{Forward propagation}

The response of the artificial neuron units in layer $j$ is computed
from the activation levels of the artificial neuron units in layer
$j-1$ using the {\em forward propagation algorithm} formulas. These
formulas are defined by an update formula for computing the {\em
netinput}, ${\bf u}^{(j)}_k$, to a unit in layer $j$ for the $k$th
training stimulus:
\begin{equation}
\label{fwd1}
{\bf u}^{(j)}_k = {\bf W}^{(j)} {\bf a}^{(j)}_k + {\bf b}^{(j)}
\end{equation}
and the {\em output activation level} for the $k$th training stimulus
in layer $j+1$:
\begin{equation}
\label{fwd2}
{\bf a}^{(j)}_k = {\bf f}^{(j)}({\bf u}^{(j)}_k)
\end{equation}
using the {\em layer transfer function}, ${\bf f}^{(j})$.
Note that ${\bf a}^{(0)}_k = {\bf s}_k$ by definition.

The layer transfer function, ${\bf f}^{(j)}$, is the chosen
vector-valued transfer function for layer $j$, ${\bf W}^{(j)}$ is the
connection matrix from the units in layer $j-1$ to the units in layer
$j$, and ${\bf b}^{(j)}$ is the bias vector for the units in layer
$j$. The subscript $k$ refers to the $k$th training stimulus. The
function ${\bf f}^{(j)}$ must have continuous second partial
derivatives.

For example, let ${\bf u} = [u_1, u_2, u_3, u_4, u_5]$.
A function form for ${\bf f}^{(j)}$ might be:
\begin{equation}
\label{transferexample}
{\bf f}^{(j)}({\bf u}) =
[ u_1, u_2, {\cal S}(u_3), exp(u_4^2), 2{\cal S}(u_5) - 1]
\end{equation}
where ${\cal S}(u_3) = 1/(1+exp(-u_3))$.  The layer transfer function
in (\ref{transferexample}) illustrates a situation where the first two
elements are {\em linear transfer functions}, the third element is a
{\em logistic sigmoidal} or just {\em sigmoidal} function,
\index{Transfer function!Sigmoidal} the fourth element is an example
of one type of {\em radial basis} transfer function, and the fifth
element is an example of one type of hyperbolic tangent sigmoidal
function (i.e., the asymptotes are -1 and +1 unlike the 0 and 1
asymptotes of the logistic sigmoid).

\subsection{Backward Propagation of Errors}\index{Backpropagation}
Let the vector $\delta^{(j)}_k = dE_k/d{\bf a}^{(j)}_k$ (i.e., the
derivative of the error with respect to the activation pattern vector
in layer $j$ evaluated at the current set of weights and the $k$th
training vector). The vector $\delta^{(j)}_k$ is called the {\em error
signal} for the $j$th layer of units whose activation levels are
defined by the vector ${\bf u}^{(j)}$ (note this notation may slightly
differ from the Rumelhart et al., 1986, paper).  Let the matrix ${\bf
Df}^{(j)}_k$ be the vector first derivative (i.e., the ``jacobian'')
of the vector transfer function ${\bf f}^{(j)}_k$ and $k$ refers to
the $k$th training simulus.

Consider a $m$-layer feedforward network consisting of layers $1,
\ldots, m$.  The {\em backward error propagation algorithm} formula
for propagating an error signal for the output units of layer $j$ to
the output units of layer $j-1$ (i.e., the input units of layer $j$)
using the formula:
% BEGIN RMG March 13, 2002 modification
%\begin{equation}
%\label{bwd}
%\delta^{(j-1)}_k = {\bf Df}^{(j)}_k [{\bf W}^{(j)}]^T \delta^{(j)}_k
%\end{equation}
\begin{equation}
\label{bwd}
\delta^{(j-1)}_k =  [{\bf W}^{(j)}]^T {\bf Df}^{(j)}_k \delta^{(j)}_k
\end{equation}
% END RMG March 13, 2002 modification
for $j = 1, \ldots, m-1$. The error signal for the $m$th layer
is the vector $\delta^{(m)}_k = de_k/d{\bf a}^{(m)}_k$
where $e_k$ is defined as in (\ref{deferrorfunk}).

For example, suppose ${\bf f}$ is defined such that ${\bf u} = {\bf
f}({\bf u})$ (i.e., a ``linear transfer function''), then ${\bf Df}$
is the identity matrix.  As another example, suppose that all output
units are logistic sigmoidal units\index{Transfer function!Sigmoidal}
so that the activation level of the $i$th output unit is given by
${\cal S}(u_i)$ where $u_i$ is the netinput to the $i$th output unit
and ${\cal S}(u_i) = 1/(1 + exp(-u_i))$.  Assume there are only three
units in a layer.  In this case, ${\bf f}$ will be a three-dimensional
vector whose $i$th on-diagonal element is defined by ${\cal S}(u_i)$
and ${\bf Df}$ will be a three-dimensional matrix whose off-diagonal
elements are zero and whose $i$th on-diagonal element is
${\cal S}(u_i) (1 - {\cal S}(u_i))$.

If the prediction error is a sum-squared error signal then
\begin{equation}
\delta^{(m)}_k = -({\bf t}_k - {\bf a}^{(m)})
\end{equation}
If the prediction error is a cross-entropy error function, then
\index{Error function!Cross-entropy}
\begin{equation}
\delta^{(m)}_k = -({\bf t}_k - {\bf a}^{(m)}) ./
[{\bf a}^{(m)} .* ({\bf 1} - {\bf a}^{(m)}]
\end{equation}
where it is assumed that the $m$ layer consists of logistic sigmoidal
functions. Note the notation ./ refers to element by element division.
\index{Transfer function!Sigmoidal}

\section{Parameter Estimation Formulas}

\subsection{Calculation of the gradient vector (and overall error).}

For each data record ($k = 1, \ldots, N$), compute the following:
\begin{itemize}
\begin{item}\index{Forward propagation}
{\em Step 1: Forward Activation Propagation.} Given stimulus ${\bf
s}_k$, propagate activation levels forward layer by layer through the
network using the forward propagation equations (equations
(\ref{fwd1}) and (\ref{fwd2})) in order to generate output activation
vector ${\bf a}^{(m)}_k$ for the last layer (layer $m$).
\end{item}
\begin{item}
{\em Step 2: Compute data record error.}  Compute the output error
$E_k$ using equation (\ref{deferrorfunk}).
\end{item}
\begin{item}
{\em Step 3: Backward error Propagation.} Compute the output error
signal $\delta^{(m)}_k$ for layer $m$ and propagate this backwards so
that an error signal is obtained for layers $m-1$ through layer 1
using equation (\ref{bwd}).
\end{item}
\begin{item}
{\em Step 4: Compute and store data record gradient.}
Let the column vector ${\bf f}^{(j-1)}_k$ be defined such that
${\bf f}^{(j-1)}_k = [{\bf a}^{(j-1)}_k \;,\; 1]^T$.
% BEGIN RMG March 13, 2002 modification
%\begin{displaymath}
%{\bf G}^{(j)}_k = \delta^{(j)}_k [{\bf f}^{(j-1)}_k]^T + (1/\sigma_w^2)
%[{\bf W} \;,\; {\bf b}]
%\end{displaymath}
\begin{equation}\label{gradcalc}
{\bf G}^{(j)}_k = \delta^{(j)}_k  {\bf Df}^{(j)}_k
[{\bf f}^{(j-1)}_k]^T + (1/\sigma_w^2)[{\bf W} \;,\; {\bf b}]
\end{equation}
% END RMG MARCH 13, 2002 modification
for $j = 1, \ldots, m$.

Let the notation $STACK[{\bf G}^{(j)}_k]$ refer to a vector formed by
``stacking'' the rows of the matrix ${\bf G}^{(j)}_k$ one on top of
another.  Thus, if ${\bf G}^{(j)}_k$ has $d$ rows, then the first $d$
elements of $STACK[{\bf G}^{(j)}_k]$ will be the first row of ${\bf
G}^{(j)}_k$, the second $d$ elements of $STACK[{\bf G}^{(j)}_k]$ will
be the second row of ${\bf G}^{(j)}_k$, and so on.  It will also be
convenient to have the function $UNSTACK$ which is the inverse of the
function $STACK$.  Now define the column vectors:
${\bf g}^{(j)}_k = STACK[{\bf G}^{(j)}_k]$ and
${\bf g}_k = [{\bf g}^{(1)}_k, \ldots, {\bf g}^{(m)}_k]$.
The vector ${\bf g}_k$ is called the {\em data record gradient}
for the $k$th data record in the training data and should be
temporarily stored once it is computed.
\end{item}
\end{itemize}

\subsection{Example learning algorithms.}\label{learningAlgs}

{\bf Important indexing issues and definitions.}
Define ${\bf y}^{(j)} = STACK [{\bf W}^{(j)} \;\;\; {\bf b}^{(j)}]$
and ${\bf y} = [{\bf y}^{(1)}, \ldots, {\bf y}^{(m)}]$. We will use
the index $t$ to denote the {\em iteration number} of the learning
algorithm.
\begin{em}
The index $t$ is not the same as the index $k$!! The $k$ index ranges
over the set of $N$ data records in the training set, while the $t$
index identifies the learning algorithm iteration!!
\end{em}
Using this notation, $E_k({\bf y}_t)$ will be used to refer to the
error for the $k$th data record at iteration number $t$ of the
learning algorithm. Similarly, ${\bf g}_k({\bf y}_t)$ refers to the
gradient of the error for the $k$th data record at iteration number
$t$ of the learning algorithm. Similarly, the variable $N$ {\bf is
NEVER} equal to the number of learning trials but is {\bf ALWAYS}
defined as the number of data records in the training set.

{\bf Initialization issues.}  In all of the following learning
algorithms the choice for the initial value ${\bf y}_0$ is absolutely
crucial. The elements of this vector should be chosen so that they are
independent and identically distributed with mean zero and variance
$\sigma_0^2$. The parameter $\sigma_0^2$ should be ``highly visible''
to the user. Learning then proceeds until the parameters of the
network converge. If $\sigma_0^2$ is too small, then the learning
process will be very long. If $\sigma_0^2$ is too large, then the
learning process will rapidly converge to a crummy solution. It is
best to err on the side of selecting $\sigma_0^2$ to be too small.

{\bf Classic on-line backprop.}\index{Backpropagation!On-line}
The {\em classic on-line}
backpropagation learning algorithm is given by the formula:
\begin{equation}
{\bf y}_{t+1} = {\bf y}_t - \eta {\bf g}_k({\bf y}_t)
\end{equation}
where $k$ (the particular training data record index) is chosen at
random on iteration $t$ of the learning algorithm.  The parameter
$\eta$ is a strictly positive number called the {\em learning rate} or
{\em stepsize} for the learning algorithm (e.g., $\eta = 1/N$ is a
good choice).

{\bf Stochastic approximation on-line backprop.}
\index{Backpropagation!On-line}
An improved version of the {\em classic on-line} algorithm called {\em
stochastic approximation on-line} is given by the formula:
\begin{equation}
{\bf y}_{t+1} = {\bf y}_t -
\left(\frac{\eta}{t_0+t} \right) {\bf g}_k({\bf y}_t)
\end{equation}
where $t_0$ is a positive integer whose magnitude is roughly equal to
the number of records, $N$, in the training data set.  The stepsize
$\eta$ is a positive number (e.g., $\eta = 1$ is a good choice).

{\bf Classic off-line (``batch'') backprop.}\label{batchepochdef}
\index{Backpropagation!Off-line}
The {\em classic off-line} or ``classic batch'' backpropagation
learning algorithm is given by the formula:
\begin{equation}
{\bf y}_{t+1} = {\bf y}_t - \eta {\bf g}({\bf y}_t)
\end{equation}
where the {\em average gradient} ${\bf g}$ is defined by the formula:
\begin{equation}
\label{gradavg}
{\bf g} = (1/N) \sum_{k=1}^N {\bf g}_k.
\end{equation}
A good choice for $\eta$ for this algorithm would be $\eta = 1$ but as
usual $\eta$ must be a strictly positive real number.

{\bf Offline backprop with automatic stepsize selection.}
\index{Backpropagation!Off-line}
Define the {\em average error} $E$ by the formula:
\begin{equation}
\label{erroravg}
E = (1/N) \sum_{k=1}^N E_k.
\end{equation}
An improved version of the {\em classic off-line} backpropagation
learning algorithm is given by the formula:
\index{Backpropagation!Off-line}
\begin{equation}
{\bf y}_{t+1} = {\bf y}_t - \eta_t {\bf g}({\bf y}_t)
\end{equation}
where ${\bf g}$ is defined as in (\ref{gradavg}) and $\eta_t$ is
chosen such that quantity $E({\bf y}_t - \eta {\bf g}({\bf y}_t))$ is
minimized with respect to $\eta$ (treating ${\bf y}_t$ as a known
constant) subject to the constraint that $0 \leq \eta_t \leq 1$.  A
``line search'' routine from {\em Numerical Recipes} may be used to
solve this one-dimensional optimization problem.
\index{Numerical Recipes in C}

\section{Data Analysis Algorithms}

\subsubsection{Conditions for Statistical Inference to be Valid.}
The four assumptions that we make here are as follows. All four of
these assumptions must be satisfied in order to guarantee reliable
statistical inferences.

{\bf Log-likelihood Objective Function Assumption.}
First, it is assumed that the prediction error objective function
$e_k$ may be interpreted as the negative natural logarithm of the
likelihood of a training data record given the network's parameters.
When the target ${\bf t}_k$ takes on a continuous range of values and
the least-squares error function defined in (\ref{leastsquaresfunk})
\index{Error function!Least squares} then this assumption is satisfied
by assuming a multivariate conditional Gaussian distribution for the
target vector given the input vector and network parameters. When the
target consists of elements taking on the values of either zero or one
and the output units are sigmoidal,\index{Transfer function!Sigmoidal}
then this assumption is satisfied by interpreting the activation level
of an output unit as the probability that the target for that output
unit will take on the value of one and using the cross-entropy
\index{Error function!Cross-entropy} function defined in
(\ref{xentropyfunk}).  The best way to
satisfy this assumption is to make sure that the user only gets to
choose from a set of valid log-likelihood objective functions.

{\bf Goodness-of-Fit Assumption.}
Second, it is assumed that the predictions of the model with respect
to the training data would be ``perfect'' if the model was exposed to
a sufficiently large set of training data and was provided with a
sufficiently large number of learning iterations. That is, we assume
the model provides a good fit to the observed training data. More
technically, it is assumed that the probabilistic modeling assumptions
introduced in the ``log-likelihood objective function assumption'' are
valid.

We will call the parameters which provide a good fit to the data ${\bf
y}^*$. Note that in many real-world situations, this assumption will
not be satisfied. For example, suppose we teach a linear ANN the
``exclusive-or'' problem. No matter how long we train the ANN, the ANN
will never ``fit the data''.

To check this assumption, provide the user with plots of training and
test set performance. Warn user that statistical inferences may not be
reliable if the ANN is not a good data model.

{\bf Convergence Assumption.}
Third, it is assumed that the gradient vector, ${\bf g}({\bf y}^*)$,
is sufficiently small in magnitude. Typically, we check this
assumption by seeing if the element in ${\bf g}_t$ with the largest
absolute value is less than $\epsilon$. The constant $\epsilon$ is
typically about $0.0001$.  If this assumption is not satisfied, the
user should be recommended to adjust the parameters of the learning
procedure or alter the initial conditions (initialization of the ANN
connection weights).

{\bf Convergence to a Strict Local Minimum Assumption.}
Fourth, it is required that this local uniqueness assumption be
satisfied. Typically, we check this assumption by computing the
matrix:
\begin{equation}
\label{hessavg}
{\bf H}^* = (1/N) \sum_{k=1}^N {\bf g}_k({\bf y}^*) [{\bf g}_k({\bf y}
^*)]^T.
\end{equation}
All eigenvalues of ${\bf H}^*$ can be shown to be always non-negative
(if you see any negative eigenvalues then there are numerical problems
with the eigenvalue calculation routine). Compute the condition number
of ${\bf H}^*$ which is defined as the largest eigenvalue of ${\bf
H}^*$ divided by the smallest eigenvalue of ${\bf H}^*$. If the
condition number is less than some constant $K$ then we say the local
uniqueness assumption is satisfied. A constant $K$ with a value of
$1000$ is acceptable.  Most statisticians use a constant $K$ of value
of $100$.  Usually, $K > 1000000$ is too large. If this assumption is
not satisfied, the user should consider collecting more data, changing
the network architecture, or decreasing the value of $\sigma_w$. Of
course the preferred course of action is collecting more data. Note
that altering the value of $\sigma_w$ is formally equivalent to
changing the network architecture.

\subsubsection{Wilks Likelihood Ratio Test}\label{Wilks}
Wilk's Likelihood Ratio Test works as follows. Estimate the parameters
of the network. Call that set of parameters ${\bf y}_{full}$. Now
constrain $r$ free parameters of the network equal to the value of
zero and re-estimate the network parameters using ${\bf y}_{full}$ as
the ``initialization'' for the learning algorithm. The resulting
parameter estimates will be called ${\bf y}_{reduced}$. It is assumed
that ${\bf y}_{full}$ satisfies all four assumptions and ${\bf
y}_{reduced}$ satisfies all assumptions except for the Goodness-of-Fit
assumption. It is assumed that ${\bf y}_{reduced}$ and ${\bf
y}_{full}$ are sufficiently close to one another in the sense that
under the null hypothesis that any difference between ${\bf
y}_{reduced}$ and ${\bf y}_{full}$ is due to the fact that we only
have a finite random sample.

Using equation (\ref{erroravg}), compute
\begin{equation}
G^2 = -2N [E({\bf y}_{full}) - E({\bf y}_{reduced})].
\end{equation}
If $G^2$ exceeds the critical value of a chi-square random variable
with $r$ degrees of freedom at the $\alpha$ significance level
($r$=number of parameters in full model minus number of parameters in
reduced model), then we reject the null hypothesis at the $\alpha$
significance level that the $r$ parameters could be set equal to zero
without hurting the fit of the model to the data.

\subsubsection{Another Version of Wilks Test called the Wald Test}
A generalization of the Z-test can be used to answer the same
questions as Wilks test. This statistical test is called the Wald
Test. It has the advantage that you can do the statistical test with
only one set of parameter estimates. First, you estimate the parameter
vector which we will call ${\bf y}^*$. Now suppose you want to test
the null hypothesis that $r$ specific elements of ${\bf y}^*$ are
equal to zero.  We will denote this subvector by the $r$-dimensional
vector ${\bf q}^*$. Obviously this is just as general as Wilks test
because you can select the $r$ elements to the connection strengths
leaving a particular input unit (for example).

Now you define a {\em selection matrix} ${\bf S}$ which has $r$ rows
and $h$ columns where $h$ is the dimension of ${\bf y}^*$. You want to
select ${\bf S}$ in just such a way so that ${\bf q}^* = {\bf S} {\bf
y}^*$. This can always be done and it might make sense to have a set
of ``canned'' selection matrix generation routines available so the
user doesn't have to figure this one out.

For example, suppose ${\bf y}^* = [1, 2, 11, 4, 6]$ and we want to
test the null hypothesis that the third and fourth weights are both
equal to zero. Then we would construct a selection matrix ${\bf S}$
with two rows whose first row is $[0,\; 0,\;1, \; 0, \; 0]$ and whose
second row is $[0, \; 0, \; 0, \; 1, \; 0]$. The only technical
restriction on the choice of the selection matrix ${\bf S}$ is that
the rank of the matrix ${\bf S}$ must be exactly equal to the number
of rows in ${\bf S}$ (i.e., the rows of ${\bf S}$ must be linearly
independent vectors).  clearly in this example the dot product of the
two row vectors in ${\bf S}$ is equal to zero implying that the two
row vectors are orthogonal and thus linearly independent.

Once you've got the selection matrix the remaining calculations are
straightforward.  Compute the formulas (using equation
(\ref{hessavg})):
\begin{equation}
{\bf C}^* = {\bf S}^* [{\bf H}^*]^{-1} ({\bf S}^*)^{T}
\end{equation}
and
\begin{equation}
{\cal W}_N = N ({\bf S}{\bf y}^*)^T [{\bf C}^*]^{-1} ({\bf S} {\bf y}^
*)
\end{equation}
where $N$ is the number of data records.

If the Wald statistic ${\cal W}_N$ exceeds the value of a chi-square
random variable with $r$ degrees of freedom (where $r$ is the number
of rows of ${\bf S}$) at the $\alpha$ significance level, then reject
the null hypothesis that the $r$ selected weights take on the value of
zero at the $\alpha$ significance level.
\index{Golden, Richard|)}

\chapter{Documentation}\label{documentation}

There is a reason that documentation is described even before the most
general features of neUROn design.  Students of computer science often
think of documentation as ancillary to code.  Nothing could be further
from the truth.  When one realizes that code is secondary to
documentation, one becomes a mature programmer.

The reason for this transcendence of documentation is actually quite
simple.  Code is merely an instance of an abstract idea, the
algorithm.  Different computer languages yield different instances,
but the algorithm is at the core.  In order to access the algorithm, a
method is coded with inputs and outputs.  These inputs and outputs are
specified by the programmer, and allow the method to become a ``black
box'' which can be used at will, and, if all goes well, as part of a
library of methods that can be used and reused often.

The job of a student programmer is to write code.  The job of a master
programmer is to write the specifications for the inputs and outputs
of methods, which when combined, solve a specific problem.  After the
specifications are written, then the algorithms within the methods are
coded.  Documentation ultimately becomes these specifications.

Documentation thus has several purposes.  It serves the programmer who
wants to reuse code written in the past, but that use is for minor
programming tasks.  It serves the programmer who wants the code used
by other programmers, a legacy for major programming tasks.
Ultimately, it serves the programmer who wants the specification to be
used by other programmers, and that is the stuff which drives, for
example, national economies.

\begin{quotation}
The neUROn programmer will document code in several ways, but in all,
{\em documentation will be written before the code.}  Not after, not
even while.  Documentation will be written {\em before} the code.
\end{quotation}

Each method (and feature corresponding to a set of methods) coded in
neUROn is to be documented in 3 specific places:

\begin{enumerate}
\item Here, in the Design Manifest.  Here a natural language
      description, algorithmic description if applicable, and
      specification is to be documented.
\item Immediately prior to the method in the code, as a specification.
\item Throughout the code.
\end{enumerate}

As an example, here is the code for the dot product method in
matrix.h:

\section{Documentation as it appears in the Design Manifest:}

Given a matrix object $M$, and vectors $\vec{v}$ and $\vec{o}$, where
$ M \cdot \vec{v} = \vec{o} $, e.g.
% NOTE: the formula must have NO blank lines!
\[
\left[ \begin{array}{cccc}
	M_{11} & M_{12} & \cdots & M_{1c} \\
	M_{21} & M_{22} & \cdots & M_{2c} \\
	\vdots & \vdots & \vdots & \vdots  \\

	M_{r1} & M_{r2} & \cdots & M_{rc}
\end{array} \right]
\cdot
\left[ \begin{array}{c}
	\vec{v_{1}} \\
	\vec{v_{2}} \\
	\vdots \\
	\vec{v_{c}}
\end{array} \right]
=
\left[ \begin{array}{c}
	\vec{o_{1}} \\
	\vec{o_{2}} \\
	\vdots \\
	\vec{o_{r}}
\end{array} \right]
\]
the following dot product operators are available in the Matrix class:

\begin{itemize}
\item $M$.dotprod( vector\& $v$ ) returns a new vector object which is
      the dot product of the matrix and the passed vector
\item $M$.dotprod( vector\& $v_{in}$, vector\& $v_{out}$ ) which
      populates the output vector $v_{out}$ with the dot product of
      the matrix and the passed vector $v_{in}$
\item $M$.dotprod( vector\& $v_{in}$, vector\& $v_{out}$, unsigned
      $begin$, unsigned $end$ ) which populates the output vector
      range
      $v_{out}[begin]$ to $v_{out}[end]$ with the dot product of the
      matrix and the passed vector $v_{in}$
\end{itemize}

\subsubsection{Examples}

\begin{verbatim}
// Matrix A has 4 rows and 3 columns, vector v1 has 3 elements
vector< double > v2;
v2 = A.dotprod( v1 ); // v1 is unchanged
                      // v2 now contains the dot product

// vector v3 has 4 elements
A.dotprod( v1, v3 ); // v1 is unchanged
                     // v3 now contains the dot product

// vector v4 has 5 elements
A.dotprod( v1, v4, 0, 3 ); // v4[0] through v4[3] now
                           // contains the dot product
\end{verbatim}

\subsubsection{Notes}
For $M$.dotprod( $v$ ), the number of columns in $M$ {\em must equal}
the number of elements in $v$.  For $M$.dotprod( $v_{in}$, $v_{out}$
), the number of columns in $M$ {\em must equal} the number of
elements in $v_{in}$, and the number of rows in $M$ {\em must equal}
the number of elements in $v_{out}$.  Because $M$.dotprod( $v$ ) must
construct a new vector object, it is less efficient than $M$.dotprod(
$v_{in}$, $v_{out}$ ).

\subsection{General principles of documentation in the Manifest:
emphasis on use of the method}

The emphasis of documentation in the Manifest is on how the method is
used.  The programmer calling the code should not need to know
anything about the inner workings of the method, only how it is
called, and what it returns.

\subsection{Emphasis on readability}

Documentation in the manifest is written in a conversational tone so
that it is easy to read and understand.

\subsection{Examples}

Good, simple examples are provided to demonstrate the use of the
method.

\subsection{Important notes}

Any specific behavior of the code is noted so that it may be used in
the most effective manner.

\section{Documentation as it would appear immediately before the
procedure in the code:}

\begin{verbatim}
// Dot product
// For vectors o and v and Matrix M, where M . v = o
//
//                         [ v1
// [ M11 M21 M31 M41    .    v2  =  [ o1
//   M12 M22 M32 M42 ]       v3       o2 ]
//                           v4 ]
//
// ( Note that the number of columns in M, 4, is the number of rows in
v )
//    by convention, ALL vectors are assumed to be vertical
// Calling sequence: M.dotprod( v, o );
\end{verbatim}

\subsection{General principles of documentation immediately before the
procedure in the code: clear argument specification}

Inputs and when applicable, returned values to the method are clearly
and succinctly stated so that casual inspection will allow the code to
be used correctly.

\subsection{Critical notes}

Critical usage issues (such as the assumption that all vectors are
assumed to be vertical in this example) are included in the comments
immediately prior to the procedure.

\section{Documentation in the code itself:}

Finally, the code itself is written and commented {\em while it is
being written.}  Debugging is the {\em last} step, and comments are
{\em modified} as the method is debugged.

\begin{verbatim}
template< class T >
vector< T >& Matrix< T >::dotprod( const vector< T >& iVec, vector< T
>& oVec ) const
{
   // Number of Matrix columns must equal number of input vector
   elements
   // Number of Matrix rows must equal number of output vector
   elements
   assert ( iVec.size() == ncols_ && oVec.size() == nrows_);

   // Set a pointer ( pData ) to Matrix's data array
   T* pData = data_;

   // Set a const iterator ( pi ) to the input vector
   vector< T >::const_iterator pi;

   // Set an iterator ( po ) to the output vector
   vector< T >::iterator po;

   // Clear output vector by setting all elements to zero
   std::fill( oVec.begin(), oVec.end(), 0 );

   // Calculate the dot product using vector iterators
   for ( po = oVec.begin(); po != oVec.end(); ++po )
      // Calculate an output vector element = Matrix row . input
      vector
      for ( pi = iVec.begin(); pi != iVec.end(); ++pi )
         *po += ( ( *pData++ ) * ( *pi ) );

   return oVec; // enables use in vector formulae
}
\end{verbatim}

\subsection{General principles of documentation within code:}

{\em Note how easy the code is to read.  These basic principles are to
be followed:}
\begin{itemize}
\item {\em Comments are meaningful.}  Nothing is more useless than a
      comment which simply restates the code, for example:

\begin{verbatim}
// Set the end of outputText to NULL
outputText[ 9 ] = '\0';
\end{verbatim}

is truly useless, whereas

\begin{verbatim}
// Truncate output string
outputText[ 9 ] = '\0';
\end{verbatim}

while less verbose, is actually more meaningful.

\item {\em Comments are capitalized correctly.}  Using all capitals is
      equivalent to email shouting, and equally offensive.  Use

\begin{verbatim}
	// Set the end of outputText to NULL
\end{verbatim}

instead of:

\begin{verbatim}
	// SET THE END OF OUTPUTTEXT TO NULL
\end{verbatim}

\item {\em Pretty printing makes the code easier to read.}  Consistent
      tabbing and bracket placement go a long way towards making clear
      code.

\item {\em Frequent meaningful comments make the code easier to read.}
      Every single code block is documented with a meaningful comment.

\end{itemize}

\subsection{Variable names:}\label{varnames}

Variables are to be meaningfully named.  It is difficult to sort
through previously written code when variables assume generic names
like `x' or `px'.  Verbose variable names are rarely a hindrance.
Often student programmers dislike verbose variable names because they
do not allow easy packing of variables into 1 line of code.  However,
this very restriction often makes the code much easier to read by
requiring the programmer to split the code into several lines.
Consider the example:

\begin{verbatim}
irpp = (EXP(LN((EXP(n*LN(1+(i/n)))-1)+1)/(ap*ppy))-1)*ap*ppy;
\end{verbatim}

While it does fit on 1 line, it is thoroughly unreadable.  As such, it
is difficult for the programmer to debug while programming, and nearly
impossible for succeeding programmers to intelligently use in future
code.  Consider the alternative:

\begin{verbatim}
// Use times_per_year_compounded = 2 for Canada, 12 for U.S.
incremental_rate = interest_rate / times_per_year_compounded;

// Calculate effective annual rate
effective_rate = EXP( times_per_year_compounded
   * LN( 1 + incremental_rate ) ) - 1;

// Finally, calculate incremental_interest = interest rate per period:
amortization = amortization_period * payments_per_year;
incremental_interest = (EXP( LN( effective_rate + 1 ) / amortization )
- 1 )
   * amortization;
\end{verbatim}

Not only is the algorithm that the programmer used to calculate the
interest rate per period manifestly clear, the variables are so well
named as to make the `//' comments complimentary rather than critical.
And when a future programmer sees the variable `incremental\_interest'
2000 lines of code later, he or she will have half a clue about what
that variable means, whereas the meaning of `irpp' would be completely
opaque.

\subsection{Space}

One of the simplest ways to make code readable is to insert space
around operators in code.  Consider the example in section
\ref{varnames},

\begin{verbatim}
irpp = (EXP(LN((EXP(n*LN(1+(i/n)))-1)+1)/(ap*ppy))-1)*ap*ppy;
\end{verbatim}

One of the difficulties in reading this code is the lack of spaces
between operators.

Rules for spaces in coding neUROn2++ are:

\begin{itemize}

\item Put spaces after binary operators (like = ):
\begin{verbatim}
a = b + c;
\end{verbatim}
rather than:
\begin{verbatim}
a=b+c;
\end{verbatim}

\item After control keywords like 'for' and 'while', put spaces around
      the ()s, e.g.:
\begin{verbatim}
for ( i = 1; i <= 10; i++ ) {
\end{verbatim}
rather than:
\begin{verbatim}
for(i = 1; i <= 10; i++){
\end{verbatim}

\item Inside of brackets, braces or parentheses, put spaces right
      after the first bracket, brace, or parenthesis, and right before
      the second, e.g.:
\begin{verbatim}
a[ i ] = 10;
\end{verbatim}
rather than:
\begin{verbatim}
a[i] = 10;
\end{verbatim}

\end{itemize}

\chapter{C++ coding}

neUROn was written in C, and took advantage of the high- and low-level
compilation features available in that language.  While neUROn2++ was
written in C++, coding began shortly after the first standardized C++
implementation became available.  Because C++ is a superset of C,
neUROn++ became a hodgepodge of C and C++ coding styles, and
ultimately, became too cumbersome to maintain, let alone extend.
neUROn2++ is designed to take advantage of the extensibility offered
by the object-oriented nature of C++, but it is also designed to take
advantage of other significant advances in C++ such as operator and
method overloading.  The following C++ style rules are expected to be
followed closely, but are not intended to be an exhaustive list.  For
those interested in an excellent source of good C++ coding style
rules, see Deitel and Deitel's {\em C++ How to
Program}.\footnote{Deitel and Deitel, C++ How to Program, 3rd edition,
Prentice Hall, 2001, ISBN 0-13-089571-7}

In general, while coding, think of the 3 cardinal rules of good
programming: code must be
\begin{itemize}
\item {\em Readable} (see Chapter \ref{documentation} on rules for
      documentation.)
\item {\em Efficient}
\item {\em Bulletproof}
\end{itemize}

Two excellent resources for C++ coding questions are the C++ Usenet
group {\em comp.lang.c++} and Marshall Cline's companion FAQ at \\
\href{http://www.parashift.com/c++-faq-lite}
{http://www.parashift.com/c++-faq-lite}.  If you have {\em any}
question regarding the most efficient way to code a particular
problem, post it on comp.lang.c++: you'll likely get a response within
a couple of hours if not minutes, and the responses often contain
excellent strategies that you would not have considered.  However,
look first in the FAQ, because if it's contained there, your post will
likely be met with the terse and somewhat embarrasing reply, ``look in
the FAQ.''

\section{Cross compilation}

All code must be compiled successfully on both GNU C++ (g++) and
Microsoft Visual C++ (MSVC++).  Code will not be considered complete
until it can be successfully compiled and executed on both platforms.

\section{Use friends sparingly}

Friend methods are to be used {\em only} when absolutely necessary.
Consider using a friend a failure to find a good solution.  Obviously,
friends are used in overloading the ostream operator $<<$ and istream
operator $>>$,
e.g.:
\begin{verbatim}
friend ostream& operator << ( ostream&, const ... );
\end{verbatim}
because there is simply no other good way to overload the ostream
operator for a new object.  However, friends should {\em not} be used
to overload mathematical operators such as $+=$ and $+$.  {\em Do not
use a friend unless you discuss it with the neUROn2++ project group
first.}

\section{Use const correctness}

Const correctness provides one way to write efficient and bulletproof
code right from the start.  By being const-correct, you think about
which objects can be mutated in a method and which cannot.  By
not allowing certain methods to mutate objects, you prevent
particularly intractable runtime errors, and by judiciously allowing
certain methods to mutate objects (e.g. passed array references), you
know exactly where your code is behaving efficiently, and exactly
where
runtime errors may be introduced.  {\em The label const is to be
applied while code is written, not after.  Constness applied as a
patch to already written code is uniformly considered to be very poor
programming practice.}

\section{Use asserts generously}

Asserts allow you to easily track down runtime errors that would be
otherwise very difficult to find, {\em but you have to include the
asserts in the first place.}  For example, in the Matrix class, here
is an assert that will catch a programmer trying to add 2 Matrices of
different sizes:
\begin{verbatim}
// Catch different size matrices
assert ( nrows_ == rhs.nrows_ && ncols_ == rhs.ncols_ );
\end{verbatim}
If the program encounters an add operation with 2 Matrices of
different sizes, it will cease immediately, and print out a meaningful
error statement with an appropriate program line, pointing you
directly to the problem.  Without the assert, the program may bomb out
with the ubiquitous stackdump or core runtime error, giving you
absolutely no debugging information, and sending you on a debugging
task that may take many hours to complete.

\chapter{vector and Matrix}\label{matrix_chapter}

Establishment of numerical vector and 2-dimensional matrix objects in
neUROn2++ serves as the basis by which clear, concise, efficient
neural computational code may be written.  The header file
``vector\_ops.h'' extends the STL $<$vector$>$ by adding simple unary
and binary mathematical operations, and overloaded ostream $<<$ and
istream $>>$ operators for simplified output and input.  The header
file ``matrix.h'' implements a matrix class for simplified
2-dimensional matrix operations.

\section{vector\_ops.h}

Inclusion of the header file ``vector\_ops.h'' overloads operators and
adds methods to the $<$vector$>$ STL.  Please see $<$vector$>$ STL
documentation for standard $<$vector$>$ STL methods.

\subsection{Unary mathematical operators}

The following operators perform element by element addition,
subtraction, multiplication and division on the right hand side (rhs)
vector and left hand side (lhs) vector, returning the result in the
left hand side (lhs) vector:

\begin{itemize}
\item $+=$: unary addition
\item $-=$: unary subtraction
\item $*=$: unary multiplication
\item $/=$: unary division
\end{itemize}

\subsubsection{Example}

\begin{verbatim}
// Standard C-style arrays
double data1[ 4 ] = { 1.0, 2.0, 3.0, 4.0 },
       data2[ 4 ] = { 2, 3, 4, 5 },
       data3[ 3 ] = { 8, 9, 10 };

// vector's std array ctor
vector< double > v1( data1, data1 + 4 ), v2( data2, data2 + 4 );

// Unary addition element by element
v1 += v2;

// Different size vectors
v3 += v1; // v1[ 3 ] will be ignored
\end{verbatim}

\subsubsection{Notes}
The vectors may have different numbers of elements, but the right hand
side (rhs) vector size {\em must be} greater than the left hand side
(lhs) vector.  If the vectors have different sizes, the remaining
elements at the end of the rhs vector that are not represented in the
lhs vector will be ignored.

\subsection{Unary mathematical operators with scalars}

The following operators perform scalar addition, subtraction,
multiplication and division on a right hand side (rhs) scalar and left
hand side (lhs) vector, returning the result in
 the left hand side
(lhs) vector:

\begin{itemize}
\item $+=$: unary scalar addition
\item $-=$: unary scalar subtraction
\item $*=$: unary scalar multiplication
\item $/=$: unary scalar division
\end{itemize}

\subsubsection{Example}

\begin{verbatim}
// Standard C-style array
double data1[ 4 ] = { 1.0, 2.0, 3.0, 4.0 };

// vector's std array ctor
vector< double > v1( data1, data1 + 4 );

// Unary scalar addition
v1 += 2.0;
\end{verbatim}

\subsubsection{Notes}
Type resolution may be a problem if the right hand side argument is
not specified, e.g. ``v1 $+=$ 2;'' may give an error.  To fix, either
cast appropriately, or include decimal points in double values,
e.g. ``v1 $+=$ 2.0;''.

\subsection{Binary mathematical operators}

The following operators perform element by element addition,
subtraction, multiplication and division on the right hand side (rhs)
vectors, returning the result for allocation in a new left hand side
(lhs) vector:

\begin{itemize}
\item $+$: binary addition
\item $-$: binary subtraction
\item $*$: binary multiplication
\item $/$: binary division
\end{itemize}

\subsubsection{Example}

\begin{verbatim}
// Standard C-style arrays
double data1[ 4 ] = { 1.0, 2.0, 3.0, 4.0 },
       data2[ 4 ] = { 2, 3, 4, 5 },
       data3[ 3 ] = { 8, 9, 10 };

// vector's std array ctor
vector< double > v1( data1, data1 + 4 ),
   v2( data2, data2 + 4 ), v3( data3, data3 + 3 ),
   v4, v5;

// Binary addition element by element
v4 = v1 + v2;

// Binary addition with different size vectors
v5 = v3 + v2; // v5 will have 3 elements
\end{verbatim}

\subsubsection{Notes}
In the passed arguments to the overloaded operator, the left hand side
(lhs) is {\em not} a reference, but a local copy of the lhs vector,
and a local copy is returned.  There is thus some inefficiency
introduced in the use of binary vector operators.  The vectors may
have different numbers of elements, but like their unary counterparts,
the right hand side (rhs) size {\em must be} greater than the size of
the lhs vector.  The final size of the returned vector will be the
same as the lhs vector size.

\subsection{Binary mathematical operators with scalars}

The following operators perform scalar addition,
subtraction, multiplication and division on a vector and a scalar, and
return the result for allocation in a new vector:

\begin{itemize}
\item $+$: binary scalar addition
\item $-$: binary scalar subtraction
\item $*$: binary scalar multiplication
\item $/$: binary scalar division
\end{itemize}

\subsubsection{Example}

\begin{verbatim}
// Standard C-style arrays
double data1[ 4 ] = { 1.0, 2.0, 3.0, 4.0 };

// vector's std array ctor
vector< double > v1( data1, data1 + 4 ), v2;

// Binary scalar addition
v2 = v1 + 3.0;
\end{verbatim}

\subsubsection{Notes}
Type resolution may be a problem if the right hand side argument is
not specified, e.g. ``v1 $+=$ 2;'' may give an error.  To fix, either
cast appropriately, or include decimal points in double values,
e.g. ``v1 $+=$ 2.0;''.
In the passed arguments to the overloaded operator, the left hand side
(lhs) is {\em not} a reference, but a local copy of the lhs vector,
and a local copy is returned.  There is thus some inefficiency
introduced in the use of binary vector operators.

\subsection{Overloaded $==$ and $!=$}

The boolean operators $==$ and $!=$ perform element by element
comparisons between vectors, $==$ returning {\em true} if they are
equal, and $!=$ returning {\em false} if they are not.

\subsubsection{Example}

\begin{verbatim}
vector< double > v1, v2;

// Code to load the vectors would follow...

if ( v1 == v2 )
   cout << "They're equal" << endl;
else
   cout << "They're not equal" << endl;
\end{verbatim}

\subsection{Overloaded $<<$}

The ostream operator $<<$ is overloaded for vector output.

\subsubsection{Example}

\begin{verbatim}
cout << "v1 = " << v1 << endl;
\end{verbatim}

\subsubsection{Notes}

This overloaded operator may also be used for vector of vector (etc.),
and will insert double (triple, etc.) spaces between the output
vectors.

\subsection{Overloaded $>>$}

The istream operator $>>$ is overloaded for vector input.

\subsubsection{Example}

\begin{verbatim}
vector< double > v1( 5 );
cout << "Enter a 5 element vector: ";
cin >> v1;
cout << "Your 5 element vector is " << v1 << endl;
\end{verbatim}

\subsection{Dot product}

The {\em dotprod} method may be used to compute a scalar dotproduct
from 2 vectors.  Its signature is:

dotprod( vector\& $v_1$, vector\& $v_2$ )

\subsubsection{Example}

\begin{verbatim}
x = dotprod( v1, v2 );
\end{verbatim}

\subsubsection{Notes}
$v_1$ and $v_2$ {\em must have} the same number of elements.

\subsection{Function passing}

The following methods allow functions to be applied to a vector
element by element:

\begin{itemize}
\item func( vector\& $v_{in}$, function $f$, vector\& $v_{out}$ ):
	applies function $f$ element by element to vector $v_{in}$,
	populating vector $v_{out}$ with the result.
\item func( vector\& $v_{in}$, function $f$, vector\& $v_{out}$,
	unsigned $begin$, unsigned $end$ ): applies function $f$
	element by element to vector $v_{in}$ in the range
	$v_{in}[begin]$ to $v_{in}[end]$, populating vector $v_{out}$
	in the range $v_{out}[begin]$ to $v_{out}[end]$ with the
	result.
\item func( vector\& $v$, function $f$ ): applies function $f$ element
	by element to vector $v$, and returns a new vector object.
\item func( vector\& $v$, function $f$, unsigned $begin$, unsigned
	$end$ ): applies function $f$ element by element to vector $v$
	in the range $v[begin]$ to $v[end]$, and returns a new
	vector object of size $v[begin] - v[end] + 1$.
\end{itemize}

\subsubsection{Example}

\begin{verbatim}\index{Transfer function!Sigmoidal}

#include <vector>
#include <math.h>

using namespace std;

#include "vector_ops.h"

// Define sigmoidal function as a class which extends
// unary_function.  This is necessary for algorithm::transform.
class sigmoidal : public unary_function< double, double >
{
   public:
      inline double operator()( const double& x )
      {
         return 1.0 / ( 1.0 + exp( -x ) );
      }
};

int main() {
   // Standard C-style arrays
   double data1[ 4 ] = { 1.0, 2.0, 3.0, 4.0 },
          data2[ 4 ] = { 2, 3, 4, 5 };

   // vector's std array ctor
   vector< double > v1( data1, data1 + 4 ),
      v2( data1, data1 + 4 ), v3( 4 ), v4;

   // Apply sigmoidal function to each element of v1
   // Populate v3 with the result
   func( v1, sigmoidal(), v3 );

   // Apply sigmoidal function to v1 with itself as result
   func( v1, sigmoidal(), v1 );

   // Apply sigmoidal function to each element of v2
   // After this function call, v4 will contain the result
   v4 = func( v2, sigmoidal() );

   // Copy vector v1 into a new vector v5
   vector< double > v5( v1 );

   // Apply the sigmoidal function to elements 1 through 2
   // in vector v1, with the result placed in vector v5
   // elements 1 through 2
   func( v1, sigmoidal(), v5, 1, 2 );

   // The receiving vector may be smaller than the sending
   // vector
   vector< double > v6( 3 );
   func( v1, sigmoidal(), v6, 0, 2 ); // v1 has 4 elements
                                      // v6 has 3 elements

   // Make a new vector
   vector< double > v7;

   // And put the sigmoidal function applied to vector v1
   // elements 1 through 2 into it
   v7 = func( v1, sigmoidal, 1, 2 ); // v7 has 2 elements

   return 0;
}
\end{verbatim}

\subsubsection{Notes}
The passed function {\em must} be defined as a class extension of
unary\_function as demonstrated in the example.  Common functions
useful for neural computation (e.g. sigmoidal) are packaged with
neUROn2++.\index{Transfer function!Sigmoidal}

If a vector to be populated is passed without bounds arguments, it
{\em must have} the same number of elements as the input vector.  If a
vector to be populated is passed with bounds arguments, the size of
the vector {\em must be} greater than the upper bound.  As func(
vector\& $v$, function $f$ ) and func( vector\& $v$, function $f$,
unsigned $begin$, unsigned $end$ ) must construct new vectors, they
are less efficient than func( vector\& $v_{in}$, function $f$,
vector\& $v_{out}$ ) and func( vector\& $v_{in}$, function $f$,
vector\& $v_{out}$, unsigned $begin$, unsigned $end$ ).  func(
vector\& $v_{in}$, function $f$, vector\& $v_{out}$ ) and func(
vector\& $v_{in}$, function $f$, vector\& $v_{out}$, unsigned $begin$,
unsigned $end$ ) are preferred.

To use a class extension of unary\_function with a scalar, the
following may be written (using sigmoidal as an example):
\index{Transfer function!Sigmoidal}

\begin{verbatim}
double x, y;
y = sigmoidal()( x );
\end{verbatim}

is valid if sigmoidal is defined as described above.

\subsection{Calculating the sum of squares of a vector elements}

The method squared( vector\& $v$ ) returns the value of the sum of
squares of vector $v$ elements.

\subsubsection{Example}
\begin{verbatim}
vector< double > v; // v is then populated with double values...
double sum;
sum = squared( v );
\end{verbatim}

\subsection{Calculating the maximum absolute value of vector elements}

The method maxabs( vector\& $v$ ) returns the maximum absolute value
of vector $v$ elements.

\subsubsection{Example}
\begin{verbatim}
vector< double > v; // v is then populated with double values...
double max;
max = maxabs( v );
\end{verbatim}

\subsection{Calculating the minimum absolute value of vector elements}

The method minabs( vector\& $v$ ) returns the minimum absolute value
of vector $v$ elements.

\subsubsection{Example}
\begin{verbatim}
vector< double > v; // v is then populated with double values...
double min;
min = minabs( v );
\end{verbatim}

\subsection{Randomization}

{\bf namespace nvec}

The method random( vector\& $v$, double $n$ ) populates a double
precision vector $v$ with double precision random numbers between $-n$
and $+n$.

\subsubsection{Example}

\begin{verbatim}
vector< double > v( 6 ); // Make a 6 element vector
// Populate it with random numbers between -0.5 and +0.5,
//    and print it out
cout << "The random vector is " << nvec::random( v, 0.5 );
\end{verbatim}

\subsubsection{Notes}
This method is only coded for vectors of type double, e.g. vector$<$
double $>$.

\subsection{Random positions}

{\bf namespace nvec}

The method random\_positions( vector\& $v$ ) populates an unsigned
vector $v$ with a series of unique integers from 0 to $n - 1$, where
$n$ is the number of elements in the vector.

\subsubsection{Example}

The example code:

\begin{verbatim}
vector< unsigned > v( 6 );
cout << nvec::random_positions( v ) << endl;
\end{verbatim}

May yield:

\begin{verbatim}
5 2 1 3 4 0
\end{verbatim}

\subsubsection{Notes}
For obvious reasons, this method is only coded for vectors of type
unsigned, e.g. vector$<$ unsigned $>$.

\subsection{Converting a vector of vector into a vector}

The following methods allow conversion of a vector of vector into a
vector, thus "flattening" the vector of vector:

\begin{itemize}
\item flatten( vector$<$ vector $>$\& $container$, vector\& $v$ )
      converts the data structure vector of vector $container$ to
      passed vector $v$, returns a reference to vector $v$.
\item flatten( vector$<$ vector $>$\& $container$ ) converts the data
      structure vector of vector $container$ to a new vector, returns
      a reference to the new vector.
\end{itemize}

\subsubsection{Example}

See section \ref{rm_var} for method {\em variable\_parse}.

\begin{verbatim}
// Create a vector of vector using util::variable_parse
string parse_test = "0-3,5; 4,8; 6,7";
vector< vector< unsigned > > parse_numbers;
parse_numbers = util::variable_parse( parse_test );
cout << "Here it is: " << parse_numbers << endl;

// Now flatten it
cout << endl << "Flattening:" << endl;
vector< unsigned > v_flat;
flatten( parse_numbers, v_flat );
cout << v_flat << endl;
cout << flatten( parse_numbers ) << endl;
\end{verbatim}

\subsubsection{Notes}
If you're using MSVC, you {\em must} include in your code the
following, as MSVC will give you error warnings for each vector$<$
vector$<$ unsigned $>$ $>$ reference:
\begin{verbatim}
#ifdef _MSC_VER	// MSVC
#pragma warning (disable: 4786)
#endif
\end{verbatim}

\subsection{Binning a vector into a vector of vector}

\begin{itemize}
\item bin( const vector\& $v$, unsigned $b$, bool $binFlag$, vector$<$
      vector $>$\& $vv$ ): partitions vector $v$ into the data
      structure vector of vector $vv$ according to specified bin size,
      with flag $binFlag$ indicating number in the last bin:
\begin{itemize}
\item {\bf false} if the last bin contains $<=$ bin size (last bin $=$
      remaining elements,
\item {\bf true} if the last bin contains $>=$ bin size (remaining
      elements are added to the last bin.
\end{itemize}
\item bin( const vector\& $v$, unsigned $b$, bool $binFlag$ ):
      partitions vector $v$ into a new data structure vector of vector
      according to specified bin size, with flag $binFlag$ indicating
      number in the last bin:
\begin{itemize}
\item {\bf false} if the last bin contains $<=$ bin size (last bin $=$
      remaining elements,
\item {\bf true} if the last bin contains $>=$ bin size (remaining
      elements are added to the last bin.
\end{itemize}
and returns a reference to the new data structure vector of vector.
\end{itemize}

\subsubsection{Example}

\begin{verbatim}
double vpart_[ 10 ] = { 0.0, 1.0, 2.0, 3.0, 4.0, 5.0, 6.0, 7.0, 8.0,
9.0 };
vector< double > vpart( vpart_, vpart_ + 10 );
vector< vector< double > > vbin;

bin( vpart, 3, false, vbin );
cout << vpart << endl << "partitioned with small last bin =" << endl <
< vbin;
cout << endl << "partitioned with large last bin =" << endl <<
   bin( vpart, 3, true ) << endl;
\end{verbatim}

Produces the following output:

\begin{verbatim}
0 1 2 3 4 5 6 7 8 9
partitioned with small last bin =
0 1 2  3 4 5  6 7 8  9
partitioned with large last bin =
0 1 2  3 4 5  6 7 8 9
\end{verbatim}

\subsubsection{Notes}
As with the previous method {\em flatten}, if you're using MSVC, you
{\em must} include in your code the following:
\begin{verbatim}
#ifdef _MSC_VER	// MSVC
#pragma warning (disable: 4786)
#endif
\end{verbatim}

\section{matrix.h and the Matrix class}

Including the header file matrix.h allows implementation of the Matrix
class, and allows simplified numerical operations on 2 dimensional
arrays.

\subsection{Constructors}

The following constructors are implemented in the Matrix class:
\begin{itemize}
\item Matrix(): Constructs a Matrix object of zero size, with zero
      rows and zero columns.
\item Matrix( unsigned $r$, unsigned $c$ ): Constructs a Matrix
      object with $r$ rows and $c$ columns.
\item Matrix( unsigned $r$, unsigned $c$, $value$ ): Constructs a
Matrix
      object with $r$ rows, $c$ columns, and fills it with
      $value$.
\item Matrix( vector\& $q$, vector\& $p^t$ ): Constructs a Matrix
      object which is the outer product of vectors $q$ and the
      transpose of $p$ (see \ref{outprod}).
\item Matrix( string\& $filename$, unsigned $c$ ): Constructs a Matrix
      object, and loads it from a file named ``$filename$'' with $c$
      columns.
\item Matrix( bool $report$, string\& $filename$ ): Constructs a
      Matrix object, and loads it from a file named ``$filename$''.
      Delimeters are any whitespace character, or any character that
      is non-numeric (+, -, e and E are considered numeric, as is a
      decimal point.)  Rows are delimited by carriage returns in the
      file, and the number of columns is the maximum number of
      elements in any one row.  Elements in unpopulated rows will be
      filled with 0, beginning from the {\em right} hand side of the
      row.  If the bool flag $report$ is true, a report will be
      printed indicating which rows were found to be unpopulated.
\end{itemize}

\subsubsection{Examples}

\begin{verbatim}
// Make some Matrices
Matrix< double > A( 4, 3 ), B, C( 5, 4, 0.01 );

// Get a file name
string filename;
cout << "The file name is ";
cin >> filename;

// Make a Matrix from the file
Matrix< double > C( filename, 3 ), // It is known to have 3 columns
                 D( true, filename ); // Number of columns unknown,
                                      // report will be printed

// Make a matrix from the outer product of vectors v1 and v2
Matrix< double > E( v1, v2 );
\end{verbatim}

\subsubsection{Notes}

A Matrix constructed by {\em Matrix( $r$, $c$ )} is {\bf not}
initially populated.  Consider it to be filled with garbage.

\subsection{Accessors}

Given a Matrix object $M$, the following accessors are implemented in
the Matrix class:

\begin{itemize}
\item $M$( unsigned $r$, unsigned $c$ ): overloaded () to
      access the element at row $r$ and column $c$ in Matrix
      {\em matrix\_object}
\item $M$.rows(): returns number of rows in Matrix
\item $M$.cols(): returns number of columns in Matrix
\end{itemize}

\subsubsection{Examples}
\begin{verbatim}
// Make a Matrix with 4 rows and 3 columns
Matrix< double > A( 4, 3 );

// Fill Matrix A with numbers
for ( i = 0; i < 4; i++ )
   for ( j = 0; j < 3; j++ )
      A( i, j ) = 3 * i + j + 1;

// Print out number of rows and number of columns in Matrix
cout << "Matrix has " << A.rows() << " rows and ";
cout << A.cols() << " columns." << endl;

\end{verbatim}

\subsection{Submatrix accessors}

The following accessors return a submatrix from a Matrix:

\begin{itemize}
\item $M$.submatrix( unsigned $row_1$, unsigned $row_N$, unsigned
      $col_1$, unsigned $col_N$, Matrix\& $M_{out}$):
      populates passed matrix $M_{out}$ with Matrix $M$ rows $row_1$
      to $row_N$, and columns $col_1$ to $col_N$.
\item $M$.submatrix( unsigned $row_1$, unsigned $row_N$, unsigned
      $col_1$, unsigned $col_N$ ): returns a new Matrix object with
      Matrix $M$ rows $row_1$ to $row_N$, and columns $col_1$ to
      $col_N$.

\end{itemize}

\subsubsection{Examples}

\begin{verbatim}
// Make some matrices
Matrix< double > P( 5, 6, 2.5 ), R( 2, 4 ), S;

// Populate matrix R with the resulting submatrix
P.submatrix( 3, 4, 2, 5, R );

// Construct a new matrix for result
S = P.submatrix( 3, 4, 2, 5 );
\end{verbatim}

\subsubsection{Notes}
Row and column dimensions must be within the bounds of the originating
matrix.  If a matrix is passed as an argument, its number of columns
{\em must equal} $col_N$ - $col_1$ + 1, and its number of rows {\em
must equal} $row_N$ - $row_1$ + 1.  A passed matrix to be populated is
more efficient than constructing a new matrix for the result.

\subsection{Row and column accessors}

The following accessors return a vector from a row or column of a
Matrix:

\begin{itemize}
\item $M$.row( unsigned $r$, vector\& $v$ ): populates vector $v$ with
      row $r$ from Matrix $M$.
\item $M$.row( unsigned $r$ ): returns a new vector with row $r$ from
      Matrix $M$.
\item $M$.col( unsigned $c$, vector\& $v$ ): populates vector $v$ with
      column $c$ from Matrix $M$.
\item $M$.col( unsigned $c$ ): returns a new vector with column $c$
      from Matrix $M$.
\end{itemize}

\subsubsection{Examples}

\begin{verbatim}
// Make a matrix with 5 rows and 6 columns, filled with Pi
Matrix< double > P( 5, 6, 3.14159 );

// Make some vectors
vector< double > v5( 5 ), v6( 6 ), v7, v8;

// Put row 2 of Matrix P in vector v6
P.row( 2, v6 );

// Construct a new vector v7 for row 2 of Matrix P
v7 = P.row( 2 );

// Put column 2 of Matrix P in vector v5
P.col( 2, v5 );

// Construct a new vector v8 for column 2 of Matrix P
v8 = P.col( 2 );
\end{verbatim}

\subsubsection{Notes}
For the methods which take passed vectors as arguments, the size of
the vector $v$ {\em must equal} the number of Matrix columns for
$M$.row( $r$, $v$ ), and the size of the vector $v$ {\em must equal}
the number of Matrix rows for $M$.col( $c$, $v$ ).  The methods which
take passed vectors as arguments are more efficient than the methods
which construct vectors.

\subsection{Methods which append vectors to Matrices}

The following methods append a vector to a Matrix:

\begin{itemize}
\item $M$.addrow( vector\& $v$, Matrix $N$ ): populates Matrix $N$ by
      appending vector $v$ to the {\em end} of Matrix $M$ as a {\em
      row}.
\item $M$.addrow( vector\& $v$ ): returns a new Matrix with vector $v$
      appended to the {\em end} of Matrix $M$.
\item $M$.addcol( vector\& $v$, Matrix $N$ ): populates Matrix $N$ by
      appending vector $v$ to the {\em right} of Matrix $M$ as a {\em
      column}.
\item $M$.addrow( vector\& $v$ ): returns a new Matrix with vector $v$
      appended to the {\em right} of Matrix $M$.
\end{itemize}

\subsubsection{Examples}

\begin{verbatim}
// Matrix A has 3 rows and 4 columns, vector v1 has 4 elements
// Matrix B has 4 rows and 4 columns
A.addrow( v1, B );
C = A.addrow( v1 );

// Matrix D has 3 rows and 4 columns, vector v2 has 3 elements
// Matrix E has 3 rows and 5 columns
D.addrow( v2, E );
F = D.addrow( v2 );
\end{verbatim}

\subsubsection{Notes}

The number of columns in $M$ and $N$ {\em must} be the same as the
number of elements of vector $v$ for $M$.addrow, and the number of
rows in $M$ and $N$ {\em must} be the same as the number of elements
of vector $v$ for $M$.addcol.

Because they construct new Matrices, $M$.addrow( vector\& $v$ ) and
$M$.addcol( vector\& $v$ ) are less efficient than $M$.addrow(
vector\& $v$, Matrix $N$ ) and $M$.addcol( vector\& $v$, Matrix $N$ ).
$M$.addrow( vector\& $v$, Matrix $N$ ) and $M$.addcol( vector\& $v$,
Matrix $N$ ) are to be preferred.

$M$.addcol is computationally expensive, and its use should be
minimized.

\subsection{Methods which replace a row or column of a Matrix with a
vector}

The following methods replace a row or column of a Matrix with a
vector:

\begin{itemize}
\item $M$.replacerow( unsigned $n$, vector\& $v$ ): replaces row $n$
      in Matrix $M$ with vector $v$, returning a reference to Matrix
      $M$ with the result.
\item $M$.replacecol( unsigned $n$, vector\& $v$ ): replaces column
      $n$ in Matrix $M$ with vector $v$, returning a reference to
      Matrix $M$ with the result.
\end{itemize}

\subsubsection{Examples}

\begin{verbatim}
// Replace row 3 of Matrix A with vector v1, and print the result
cout << A.replacerow( 3, v1 );

// Replace column 2 of Matrix A with vector v2, and print the result
cout << A.replacecol( 2, v2 );
\end{verbatim}

\subsubsection{Notes}

The number of columns in $M$ {\em must} be the same as the number of
elements of vector $v$ for $M$.replacerow, and the number of rows in
$M$ {\em must} be the same as the number of elements of vector $v$ for
$M$.replacecol.

\subsection{Column inclusion and exclusion}

The following operators include or exclude those columns indicated by
the passed vector, which includes the (unsigned) column values:

\begin{itemize}
\item $M$.includecols( Matrix\& $M_{out}$, vector$<$ unsigned $>$\&
      $v$ ): includes in the returned Matrix $M_{out}$ only those
      columns from Matrix $M$ indicated by $v$, returning a reference
      to Matrix $M_{out}$ with the result.
\item $M$.includecols( vector$<$ unsigned $>$\& $v$ ): returns a new
      Matrix {\em with} only those columns from Matrix $M$ indicated
      by $v$
\item $M$.excludecols( Matrix\& $M_{out}$, vector$<$ unsigned $>$\&
      $v$ ): excludes in the returned Matrix $M_{out}$ those columns
      from Matrix $M$ indicated by $v$, returning a reference to
      Matrix $M_{out}$ with the result.
\item $M$.excludecols( vector$<$ unsigned $>$\& $v$ ): returns a new
      Matrix {\em without} those columns from Matrix $M$ indicated by
      $v$
\end{itemize}

\subsubsection{Examples}

\begin{verbatim}
// Make a 5x4 Matrix, fill with some numbers
Matrix< double > M( 5, 4 );
for ( unsigned row = 0; row < 5; row++ )
   for ( unsigned column = 0; column < 4; column++ )
      M( row, column ) = 4 * column + row;
cout << "Here is the original Matrix:" << M;

// Make some smaller Matrices
Matrix< double > M1( 5, 2 ), M2, M3( 5, 2 ), M4;

// Make a vector to indicate columns 1 & 3
vector< unsigned > v;
v.push_back( 1 );
v.push_back( 3 );

// Passing a Matrix to include columns
M.includecols( M1, v );
cout << "From the matrix above, including columns "
   << v << ": " << M1;

// A new Matrix to include columns
M2 = M.includecols( v );
cout << "Returning a new Matrix: " << M2;

// Passing a Matrix to exclude columns
M.excludecols( M3, v );
cout << "From the matrix above, excluding columns "
   << v << ": " << M3;

// A new Matrix to exclude columns
M4 = M.excludecols( v );
cout << "Returning a new Matrix: " << M4;
\end{verbatim}

\subsubsection{Notes}
If the Matrix $M_{out}$ to be populated is passed as an argument, its
rows {\em must} be the same as the number of rows in $M$, and its
number of columns {\em must} be the same as the difference between the
number of columns in $M$ and the number of elements in the passed
vector $v$.  The elements in the passed vector $v$ {\em must} be {\em
unique} (no duplicated values), and no element can be $>=$ to the
number of columns in Matrix $M$.

These methods are useful to include or exclude variables or nodes
represented by columns in a Matrix.

\subsection{Boolean equality operators}

The following operators perform boolean equality and inequality
comparisons, element by element:

\begin{itemize}
\item $==$: equality
\item $!=$: inequality
\end{itemize}

\subsubsection{Example}

\begin{verbatim}
// Make 2 Matrices with 4 rows and 3 columns
Matrix< double > A( 4, 3 ), B( 4, 3 );

// Fill Matrices A and B with numbers

for ( i = 0; i < 4; i++ )
{
   for ( j = 0; j < 3; j++ )
   {
      A( i, j ) = 3 * i + j;
      B( i, j ) = 3 * i + j + 1;
   }
}

// Compare Matrix A to B:
if ( A == B )
   cout << "They're equal" << endl;
if ( A != B )
   cout << "They're not equal" << endl;
\end{verbatim}

\subsection{Unary mathematical operators}

The following operators perform element by element addition,
subtraction, multiplication and division on the right hand side (rhs)
Matrix and left hand side (lhs) Matrix, returning the result in the
left hand side (lhs) Matrix:

\begin{itemize}
\item $+=$: unary addition
\item $-=$: unary subtraction
\item $*=$: unary multiplication
\item $/=$: unary division
\end{itemize}

\subsubsection{Example}

\begin{verbatim}
// Make 2 Matrices with 4 rows and 3 columns
Matrix< double > A( 4, 3 ), B( 4, 3 );

// Fill Matrices A and B with numbers

for ( i = 0; i < 4; i++ )
{
   for ( j = 0; j < 3; j++ )
   {
      A( i, j ) = 3 * i + j;
      B( i, j ) = 3 * i + j + 1;
   }
}

// Add matrix B to A element by element
A += B;
\end{verbatim}

\subsection{Unary mathematical operators with scalars}

The following operators perform scalar addition, subtraction,
multiplication and division on a right hand side (rhs) scalar and left
hand side (lhs) Matrix, returning the result in the left hand side
(lhs) Matrix:

\begin{itemize}
\item $+=$: unary scalar addition
\item $-=$: unary scalar subtraction
\item $*=$: unary scalar multiplication
\item $/=$: unary scalar division
\end{itemize}

\subsubsection{Example}

\begin{verbatim}
// Make a Matrix with 4 rows and 3 columns, filled with 1.0
Matrix< double > A( 4, 3, 1.0 );

// Add 3.0 to each element of A
A += 3.0;

\end{verbatim}

\subsection{Binary mathematical operators}

The following operators perform element by element addition,
subtraction, multiplication and division on the right hand side (rhs)
Matrices, returning a new Matrix object:

\begin{itemize}
\item $+$: binary addition
\item $-$: binary subtraction
\item $*$: binary multiplication
\item $/$: binary division
\end{itemize}

\subsubsection{Example}

\begin{verbatim}
// Make 2 Matrices with 4 rows and 3 columns, and a 3rd empty Matrix
Matrix< double > A( 4, 3 ), B( 4, 3 ), C;

// Fill Matrices A and B with numbers
for ( i = 0; i < 4; i++ )
{
   for ( j = 0; j < 3; j++ )
   {
      A( i, j ) = 3 * i + j;
      B( i, j ) = 3 * i + j + 1;
   }
}

// Add matrix A and B element by element, the result is matrix C
C = A + B;
\end{verbatim}

\subsubsection{Notes}
The Matrices {\em must have} the same number of rows and columns.
As a new matrix object must be constructed, unary operations are
preferred for efficiency.

\subsection{Binary mathematical operators with scalars}

The following operators perform scalar addition, subtraction,
multiplication and division on a Matrix, returning a new Matrix
object:

\begin{itemize}
\item $+$: binary scalar addition
\item $-$: binary scalar subtraction
\item $*$: binary scalar multiplication
\item $/$: binary scalar division
\end{itemize}

\subsubsection{Example}

\begin{verbatim}
// Make some Matrices
Matrix< double > A( 4, 3, 1.0 ), B;

// Add 3.0 to each element of A, B is the result
B = A + 3.0;
\end{verbatim}

\subsubsection{Notes}
As a new matrix object must be constructed, unary operations are
preferred for efficiency.  {\em \bf Known bug: the scalar must
be on the right side of the Matrix for all operators.}

\subsection{Transpose}

The following operators transpose a Matrix:

\begin{itemize}
\item $M$.t() transposes matrix $M$, and returns a matrix object
\item $M$.t( $O$ ) populates passed matrix $O$ with the transpose
      of $M$
\end{itemize}

\subsubsection{Example}

\begin{verbatim}
// A is a 3 x 4 matrix
Matrix< double > B( 4, 3 ), C;

A.t( B );  // B now contains the transpose of A
C = A.t(); // C now contains the transpose of A
\end{verbatim}

\subsubsection{Notes}

For $M$.t( $O$ ), the number of columns in $M$ {\em must equal} the
number of rows in $O$ and the number of rows in $M$ {\em must equal}
the number of columns in $O$.

Because $M$.t() must construct a new matrix object, it is less
efficient than $M$.t( $O$ ).

\subsection{Overloaded $<<$}

The ostream operator $<<$ is overloaded for Matrix output.

\subsubsection{Example}

\begin{verbatim}
cout << "Matrix A = " << A << endl;
\end{verbatim}

\subsubsection{setHeader}

The setHeader method:

\begin{itemize}
\item $M$.setHeader( bool $flag$ )
\end{itemize}

is provided for ostream output.  If $flag$ is true, a header line
containing the format:
\begin{center}
$r$ rows and $c$ columns:
\end{center}
is output to the stream prior to the Matrix, whereas if $flag$ is
false, the header line is not output, only the Matrix.  {\em By
default, this flag is set to} {\bf true}.

\subsubsection{Examples}

\begin{verbatim}
cout << "Matrix A alone = " << A.setHeader( false ) << endl;
A.setHeader( true );
cout << "Matrix A with header = " << A << endl;
\end{verbatim}

\subsection{Overloaded $>>$}

The istream operator $>>$ is overloaded for Matrix input.

\subsubsection{Example}

\begin{verbatim}
Matrix< double > A( 3, 2 );
cout << "Enter Matrix A: ";
cin >> A;
\end{verbatim}

\subsection{Dot product}

Given a matrix object $M$, and vectors $\vec{v}$ and $\vec{o}$, where
$ M \cdot \vec{v} = \vec{o} $, e.g.
% NOTE: the formula must have NO blank lines!
\[
\left[ \begin{array}{cccc}
	M_{11} & M_{12} & \cdots & M_{1c} \\
	M_{21} & M_{22} & \cdots & M_{2c} \\
	\vdots & \vdots & \vdots & \vdots  \\

	M_{r1} & M_{r2} & \cdots & M_{rc}
\end{array} \right]
\cdot
\left[ \begin{array}{c}
	\vec{v_{1}} \\
	\vec{v_{2}} \\
	\vdots \\
	\vec{v_{c}}
\end{array} \right]
=
\left[ \begin{array}{c}
	\vec{o_{1}} \\
	\vec{o_{2}} \\
	\vdots \\
	\vec{o_{r}}
\end{array} \right]
\]
or other matrix objects $N$ and $O$, where $M \cdot N = O$,
the following dot product operators are available in the Matrix class:

\begin{itemize}
\item $M$.dotprod( vector\& $v$ ) returns a new vector object which is
      the dot product of the matrix and the passed vector
\item $M$.dotprod( vector\& $v_{in}$, vector\& $v_{out}$ ) which
      populates the output vector $v_{out}$ with the dot product of
      the matrix and the passed vector $v_{in}$
\item $M$.dotprod( vector\& $v_{in}$, vector\& $v_{out}$, unsigned
      $begin$, unsigned $end$ ) which populates the output vector
      range $v_{out}[begin]$ to $v_{out}[end]$ with the dot product of
      the matrix and the passed vector $v_{in}$
\item $M$.dotprodt( vector\& $v$ ) returns a new vector object which
      is the dot product of the {\em transpose} of the matrix and the
      passed vector
\item $M$.dotprodt( vector\& $v_{in}$, vector\& $v_{out}$ ) which
      populates the output vector $v_{out}$ with the dot product of
      the {\em transpose} of the matrix and the passed vector $v_{in}$
\item $M$.dotprodt( vector\& $v_{in}$, vector\& $v_{out}$, unsigned
      $begin$, unsigned $end$ ) which populates the output vector
      range $v_{out}[begin]$ to $v_{out}[end]$ with the dot product of
      the {\em transpose} of the matrix and the passed vector $v_{in}$
\item $M$.dotprod( $N$ ) returns a new Matrix object which is the dot
      product of the matrix $M$ and the passed matrix $N$
\item $M$.dotprod( $N$, $O$ ) which populates the passed matrix $O$
      with the dot product of the matrix $M$ and the passed matrix $N$

\end{itemize}

\subsubsection{Examples}

(dotprodt($\dots$) follows the same form as dotprod($\dots$).)

\begin{verbatim}
// Matrix A has 4 rows and 3 columns, vector v1 has 3 elements
vector< double > v2;
v2 = A.dotprod( v1 ); // v1 is unchanged
                      // v2 now contains the dot product

// vector v3 has 4 elements
A.dotprod( v1, v3 ); // v1 is unchanged
                     // v3 now contains the dot product

// vector v4 has 5 elements
A.dotprod( v1, v4, 0, 3 ); // v4[0] through v4[3] now
                           // contains the dot product

// Matrix B has 3 rows and 5 columns
Matrix< double >C ( 4, 5 ), D;
A.dotprod( B, C );  // C now contains the dot product
D = A.dotprod( B ); // D now contains the dot product
\end{verbatim}

\subsubsection{Notes}
For $M$.dotprod( $v$ ), the number of columns in $M$ {\em must equal}
the number of elements in $v$.  For $M$.dotprod( $v_{in}$, $v_{out}$
), the number of columns in $M$ {\em must equal} the number of
elements in $v_{in}$, and the number of rows in $M$ {\em must be equal
to or greater than} the number of elements in $v_{out}$.  If the
number of rows in $M$ is greater than the number of elements in
$v_{out}$, the resulting vector will be truncated to the size of
$v_{out}$, and the remainder ignored.

For $M$.dotprodt( $v$ ), the number of rows in $M$ {\em must equal}
the number of elements in $v$.  For $M$.dotprodt( $v_{in}$,
$v_{out}$),
the number of rows in $M$ {\em must equal} the number of elements
in $v_{in}$, and the number of columns in $M$ {\em must be equal to or
greater than} the number of elements in $v_{out}$.  If the number of
columns in $M$ is greater than the number of elements in $v_{out}$,
the resulting vector will be truncated to the size of $v_{out}$, and
the remainder ignored.

For $M$.dotprod( $N$ ), the number of columns in $M$ {\em must equal}
the number of rows in $N$.  For $M$.dotprod( $N$, $O$ ), the number of
rows in $M$ {\em must equal} the number of rows in $O$, the number of
columns in $N$ {\em must equal} the number of columns in $O$, and the
number of columns in $M$ {\em must equal} the number of rows in $N$.

Because $M$.dotprod( $v$ ) and $M$.dotprodt( $v$ ) must construct new
vector objects, and $M$.dotprod( $N$ ) must construct a new matrix
object, these methods are less efficient than $M$. dotprod( $v_{in}$,
$v_{out}$ ), $M$.dotprodt( $v_{in}$, $v_{out}$ ) and $M$.dotprod( $N$,
$O$ ).

\subsubsection{dotprod\_row}

The {\em dotprod\_row} method takes a single row from a Matrix, and
transposes it to be used as the right hand side argument for a dot
product.  It is useful for extracting a single row from a dataset and
passing it to the first layer of a neural network.  It is functionally
equivalent to $M \cdot {\vec{d}}^{t} = \vec{o}$, where $M$ is the
matrix, ${\vec{d}}^{t}$ is the transpose of the vector $\vec{d}$
derived
from a single row of the dataset matrix, and $\vec{o}$ is the output
vector.

Suppose that $d$ is a single row of the dataset matrix $D$, e.g.:
\[
\left[ \begin{array}{cccc}
	\cdots & \cdots & \cdots & \cdots \\
	d_{r1} & d_{r2} & \cdots & d_{rc} \\
	\vdots & \vdots & \vdots & \vdots \\
	\cdots & \cdots & \cdots & \cdots
\end{array} \right]
\]

Then $M$.dotprod\_row( $D$, $r$, $o$ ) is equivalent to:

\[
\left[ \begin{array}{cccc}
	M_{11} & M_{12} & \cdots & M_{1c} \\
	M_{21} & M_{22} & \cdots & M_{2c} \\
	\vdots & \vdots & \vdots & \vdots  \\
	M_{r1} & M_{r2} & \cdots & M_{rc}
\end{array} \right]
\cdot
\left[ \begin{array}{c}
	\vec{d_{r1}} \\
	\vec{d_{r2}} \\
	\vdots \\
	\vec{d_{rc}}
\end{array} \right]
=
\left[ \begin{array}{c}
	\vec{o_{1}} \\
	\vec{o_{2}} \\

	\vdots \\
	\vec{o_{r}}
\end{array} \right]
\]

where $r$ is the row number.  The following dotprod\_row operators are
available in the Matrix class:

\begin{itemize}
\item $M$.dotprod\_row( Matrix\& $D$, unsigned $r$, vector\& $v$ )
which
      populates the passed vector $v$ with the dot product of Matrix
      $M$
      and the transposed row $r$ in matrix $D$
\item $M$.dotprod\_row( Matrix\& $D$, unsigned $r$ ) which returns a
new
      Matrix object with the dot product of Matrix $M$
      and the transposed row $r$ in matrix $D$
\item $M$.dotprod\_row( Matrix\& $D$, unsigned $r$, vector\& $v$,
      unsigned $begin$, unsigned $end$ ) which populates the range in
      the passed vector $v[begin]$ to $v[end]$ with the dot product of
      Matrix $M$ and the transposed row $r$ in matrix $D$
\end{itemize}

\subsubsection{Examples}

\begin{verbatim}
// Matrix A and D both have 3 columns, vector v1 has 3 elements
A.dotprod_row( D, 3, v1 ); // v1 now contains the dot product
                           // of A and the transpose of D's 3rd row

vector< double > v2;
v2 = A.dotprod_row( D, 3 ); // v2 now contains the dot product
                            // of A and the transpose of D's 3rd row

// Matrix A has 4 rows, vector v3 has 5 elements
A.dotprod_row( D, 3, v3, 0, 3 ); // v3[0] through v3[3] now
                                 // contains the dot product of A
                                 // and the transpose of D's 3rd row
\end{verbatim}

\subsubsection{Notes}
Both matrices {\em must have} the same number of columns.  If a vector
is passed, it {\em must have} the same number of elements as the
matrices have columns.  Because $M$.dotprod\_row( $D$, $r$ ) must
construct a new vector object, it is less efficient than
$M$.dotprod\_row( $D$, $r$, $v$ ).

\subsection{Outer product}\label{outprod}

The {\em outprod} method is available in the Matrix class to populate
a Matrix with the outer product of 2 vectors.  For example, given 2
vectors $\vec{p}$ and $\vec{q}$, e.g.:

\[
\vec{q}
=
\left[ \begin{array}{c}
	q_1 \\
	q_2 \\
	\vdots \\
	q_m
\end{array} \right]
, \vec{p} =
\left[ \begin{array}{c}
	p_1 \\
	p_2 \\
	\vdots \\
	p_n
\end{array} \right]
\]

The outer product $\vec{q}\vec{p}^t$ is:

\[
\vec{q}\vec{p}^t =
\left[ \begin{array}{cccc}
	q_1p_1 & q_1p_2 & \cdots & q_1p_n \\
	q_2p_1 & q_2p_2 & \cdots & q_2p_n \\
	\vdots & \vdots & \vdots & \vdots  \\
	q_mp_1 & q_mp_2 & \cdots & q_mp_n
\end{array} \right]
\]

The outer product method is called as in:

$M$.outprod( vector\& $q$, vector\& $p$ )

\subsubsection{Example}

\begin{verbatim}
// Make some vectors
double array1[ 4 ] = { 0, 1, 2, 3 },
   array2[ 5 ] = { 0, 1, 2, 3, 4 };
vector< double > vec1( array1, array1 + 4 ),
   vec2( array2, array2 + 5 );

// Make the right size Matrix
Matrix< double > U( 4, 5 );

// Populate the Matrix with the outer product of the vectors
U.outprod( vec1, vec2 );
\end{verbatim}

\subsubsection{Notes}
The number of elements in the first vector argument {\em must equal}
the number of Matrix rows, and the number of elements in the second
vector argument {\em must equal} the number of Matrix columns.

\subsection{Function passing}

The following methods allow functions to be applied to a Matrix
element by element:

\begin{itemize}
\item func( Matrix\& $M_{in}$, function $f$, Matrix\& $M_{out}$ ):
      applies function $f$ element by element to Matrix $M_{in}$,
      populating Matrix $M_{out}$ with the result.
\item func( Matrix\& $M$, function $f$ ): applies function $f$ element
      by element to Matrix $M$, and returns a new Matrix object.
\end{itemize}

\subsubsection{Example}

\begin{verbatim}\index{Transfer function!Sigmoidal}

#include <vector>
#include <math.h>

using namespace std;

#include "vector_ops.h"

// Define sigmoidal function as a class which extends
// unary_function.  This is necessary for algorithm::transform.
class sigmoidal : public unary_function< double, double >
{
   public:
      inline double operator()( const double& x )
      {
         return 1.0 / ( 1.0 + exp( -x ) );
      }
};

int main() {
   // Make some Matrices
   Matrix < double > M1( 4, 3, 0.5 ), M2( 4, 3, 0.0 ),
      M3( 4, 3, 2.0 ), M4;

   // Apply sigmoidal function to each element of M1
   // Populate M2 with the result
   func( M1, sigmoidal(), M2 );

   // Apply sigmoidal function to a Matrix, with itself as result
   func( M3, sigmoidal(), M3 );

   // Apply sigmoidal function to each element of M1
   // After this function call, M4 will contain the result
   M4 = func( M1, sigmoidal() );

   return 0;
}
\end{verbatim}

\subsubsection{Notes}
The passed function {\em must} be defined as a class extension of
unary\_function as demonstrated in the example.  Common functions
useful for neural computation (e.g. sigmoidal) are packaged with
neUROn2++.\index{Transfer function!Sigmoidal}

If a Matrix to be populated is passed, it {\em must have} the same
number of rows and columns as the input Matrix.  As func( Matrix\&
$M$, function $f$ ) must construct a new vector, it is less efficient
than func( Matrix\& $M_{in}$, function $f$, Matrix\& $M_{out}$ ).

\subsection{Randomization}

The method $M$.random( double $n$ ) populates a double precision
Matrix $M$ with double precision random numbers between $-n$ and $+n$.

\subsubsection{Example}

\begin{verbatim}
Matrix< double > Z( 6, 5 ); // Make a 6 x 5 Matrix
// Populate it with random numbers between -0.5 and +0.5,
//    and print it out
cout << "The random matrix is " << Z.random( 0.5 );
\end{verbatim}

\subsubsection{Notes}
This method is only coded for Matrices of type double, e.g. Matrix$<$
double $>$.

\subsection{Loading a Matrix from a file}

The following methods will load a matrix from a file:
\begin{itemize}
\item $M$.loadfile( string\& $filename$, unsigned $c$ ): loads Matrix
      $M$ from a file named ``$filename$'' with $c$ columns.  The
      Matrix will be resized if necessary, and if the file does not
      exist, the user will be queried until a good filename is
      entered.
\item $M$.loadfile( bool $report$, string\& $filename$ ): loads Matrix
      $M$ from a file named ``$filename$''. Delimeters are any
      whitespace character, or any character that is non-numeric (+,
      -, e and E are considered numeric, as is a decimal point.)  Rows
      are delimited by carriage returns in the file, and the number of
      columns is the maximum number of elements in any one row.
      Elements in unpopulated rows will be filled with 0, beginning
      from the {\em right} hand side of the row.  If the bool flag
      $report$ is true, a report will be printed indicating which rows
      were found to be unpopulated.  The Matrix will be resized if
      necessary, and if the file does not exist, the user will be
      queried until a good filename is entered.
\end{itemize}

\subsubsection{Examples}

\begin{verbatim}
Matrix< double > A, B;
string filename;
unsigned cols;

cout << "What is the name of the data file? ";
cin >> filename;
cout << "How many columns does it have? ";
cin >> cols;

A.loadfile( filename, cols );

cout << "What is the name of another data file? ";
cin >> filename;

B.loadfile( true, filename ); // a report will be printed
\end{verbatim}

\subsection{Saving a Matrix to a file}

The following method will save a matrix to a file:
\begin{itemize}
\item bool $M$.savefile( string\& $filename$ ): saves Matrix $M$ to a
      file named ``$filename$'', and returns a Boolean flag indicating
      whether the save operation was successful.
\end{itemize}

\subsubsection{Example}

\begin{verbatim}
Matrix< double > A( 5, 4, 3.14159 ); // fill a 5 x 4 Matrix with pi
A.savefile( "pi.txt" ); // and save it to a file named pi.txt
\end{verbatim}

\subsection{Resizing a Matrix}

The method $M$.resize( unsigned $r$, unsigned $c$ ) resizes Matrix $M$
to $r$ rows and $c$ columns.

\subsubsection{Example}
\begin{verbatim}
A.resize( 5, 4 ); // resize Matrix A to 5 rows and 4 columns
\end{verbatim}

\subsubsection{Notes}
The following Matrix so resized is filled with {\em \bf garbage}.

\subsection{Clearing a Matrix}

The method $M$.clear() clears Matrix $M$, setting its dimensions and
data to zero.

\subsubsection{Example}
\begin{verbatim}
A.clear(); // clear Matrix A
\end{verbatim}

\subsection{Populating a Matrix with a value}

The method $M$.fill( $v$ ) populates Matrix $M$ with value $v$ of its
type.

\subsubsection{Example}
\begin{verbatim}
Matrix< double > A( 3, 2 ); // Make a Matrix of type double with
                            // 3 rows and 2 columns

A.fill( 3.14159 ); // and fill it with pi
\end{verbatim}

\subsection{Calculating the sums of columns of a Matrix}

The method $M$.colsums() returns a new vector with each element the
sum of the respective column in $M$.  The method $M$.colsums( vector\&
$v$ ) returns the result in the passed vector $v$.

\subsubsection{Example}
\begin{verbatim}
#include <iostream>
#include <string>
#include "matrix.h"
#include "vector_ops.h"

string filename;
cout << "The file name is ";
cin >> filename;
Matrix< double > M( true, filename ); // load Matrix from file
vector< double > v;
v = M.colsums(); // v will contain the sum of the columns in M
cout << "Sum of columns = " << v << endl;
\end{verbatim}

\subsubsection{Note}

If the vector to be returned is passed as an argument, its size {\em
must} be the same as the number of columns in the Matrix.

\subsection{Calculating the sum of squares of a Matrix elements}

The method $M$.squared() returns a the value of the sum of squares
of Matrix elements.

\subsubsection{Example}
\begin{verbatim}
Matrix< double > M; // M is then populated with double values...
double sum;
sum = M.squared();
\end{verbatim}

\subsection{Calculating the maximum absolute value of Matrix elements}

The method $M$.maxabs() returns the maximum absolute value of Matrix
elements.

\subsubsection{Example}
\begin{verbatim}
Matrix< double > M; // M is then populated with double values...
double max;
max = M.maxabs();
\end{verbatim}

\subsection{Finding the row elements with value 1 in a Matrix}

The method (bool) $M$.rowindex( vector$<$ unsigned $>$\& $v$ )
populates a passed (unsigned) vector $v$ with each vector element
representing in which Matrix column each of the Matrix row elements
was 1, takes passed vector as argument, returns true if successful.

\subsubsection{Example}
\begin{verbatim}
vector< unsigned > v( M.rows() );
if ( M.rowindex( v ) )
   cout << v << endl;
else
   cout << "can't do it: goofy Matrix." << endl;
\end{verbatim}

\subsubsection{Note}
The passed Matrix must consist of rows that have a single row element
with a value of 1, and the remainder with values of 0.  For example,
the Matrix

% NOTE: the formula must have NO blank lines!
\[ \left[ \begin{array}{cccc}
1 & 0 & 0 & 0 \\
0 & 0 & 1 & 0 \\
0 & 1 & 0 & 0 \\
\end{array} \right] \]

would return true and populate the passed vector with

% NOTE: the formula must have NO blank lines!
\[ \left[ \begin{array}{c}
0 \\
2 \\
1 \\
\end{array} \right] \]

whereas Matrices

\[ \left[ \begin{array}{cccc}
1 & 1 & 0 & 0 \\
0 & 0 & 1 & 0 \\
0 & 1 & 0 & 0 \\
\end{array} \right] \]

\[ \left[ \begin{array}{cccc}
0 & 0 & 0 & 0 \\
0 & 0 & 1 & 0 \\
0 & 1 & 0 & 0 \\
\end{array} \right] \]

\[ \left[ \begin{array}{cccc}
3 & 0 & 0 & 0 \\
0 & 0 & 1 & 0 \\
0 & 1 & 0 & 0 \\
\end{array} \right] \]

would all return false.

The size of the vector {\em must} be the same as the number of rows in
the Matrix.

\subsection{Converting a Matrix to a vector, and a vector to a Matrix}

A Matrix may be coverted row by row to a vector, e.g. Matrix

\[ \left[ \begin{array}{cccc}
1 & 2 & 3 & 4 \\
5 & 6 & 7 & 8 \\
9 & 10 & 11 & 12 \\
\end{array} \right] \]

would be converted to vector

\[ \left[ \begin{array}{cccccccccccc}
1 & 2 & 3 & 4 & 5 & 6 & 7 & 8 & 9 & 10 & 11 & 12
\end{array} \right] \]

A vector may also be converted to a Matrix.  Methods to perform the
conversion include:

\begin{itemize}
\item $M$.toVector( vector\& $v$ ): converts Matrix $M$ to passed
      vector $v$ which is populated with the result, and returns a
      reference to vector $v$
\item $M$.toVector(): converts Matrix $M$ to a new vector, and returns
      a reference to the new vector
\item toMatrix( Matrix\& $M$, vector\& $v$ ): converts vector $v$ to
      passed Matrix $M$ which is populated with the result, and
      returns a reference to Matrix $M$
\item toMatrix( vector\& $v$, unsigned $r$, unsigned $c$ ): converts
      vector $v$ to a new Matrix $M$ with $r$ rows and $c$ columns,
      and returns a reference to the new Matrix
\end{itemize}

\subsubsection{Example}
\begin{verbatim}
// Make a Matrix as an example, and print it out
Matrix< double > A( 4, 3 );
for ( unsigned i = 0; i < 4; i++ )
    for ( unsigned j = 0; j < 3; j++ )
        A( i, j ) = 3 * i + j + 1;
cout << "Matrix A: " << A << endl;

// Convert the Matrix to vectors
vector< double > v1( 12 ), v2;
A.toVector( v1 );
cout << "toVector passed existing vector: " << v1 << endl;
v2 = A.toVector();
cout << "toVector passed new vector: " << v2 << endl;

// Convert the vectors back to Matrices
Matrix< double > B( 4, 3 ), C;
toMatrix( B, v2 );
cout << "back to an existing Matrix: " << B << endl;
C = toMatrix( v2, 4, 3 );
cout << "back to a new Matrix: " << C << endl;
\end{verbatim}

\subsubsection{Notes}
If a vector is passed as an argument to toVector, its size {\em must
be} the same as the product of the rows and the columns of the Matrix.
If a Matrix is passed as an argument to toMatrix, the product of its
rows and columns {\em must be} the same as the size of the vector
which is also passed.  If a Matrix is not passed as an argument to
toMatrix, the product of the $2^{nd}$ and $3^{rd}$ arguments {\em must
be} the same as the size of the passed vector.

\subsection{First order covariance of a Matrix}

The method $M$.covariance() returns a new Matrix with the first-order
covariance Matrix of $M$.  The method $M$.covariance( Matrix$<$ double
$>$\& $N$ ) populates the passed Matrix $N$ with the result, and
returns a reference to Matrix $N$.

\subsubsection{Example}
\begin{verbatim}
#include <iostream>
#include <string>
#include "matrix.h"

string filename;
cout << "The file name is ";
cin >> filename;
Matrix< double > Data( true, filename ), // load data Matrix from file
   C; // for the covariance Matrix

C = Data.covariance();
cout << "Covariance Matrix = " << C << endl;
\end{verbatim}

\subsubsection{Notes}

The covariance method is only coded for Matrices of type $<$ double
$>$.  If the Matrix to be populated with the result is passed as an
argument, it's dimensions must be $Col \times Col$, where $Col$ is the
number of columns in the data Matrix.

\subsection{Matrix inverse and determinant}

The following methods calculate the inverse and determinant of a
Matrix:\footnote{Adapted from W. H. Press et. al, Numerical
Recipes in C, 2nd edition, Cambridge University Press, Cambridge,
1992, pages 36-48}\index{Numerical Recipes in C}

\begin{itemize}
\item $M$.inverse( Matrix$<$ double $>$\& $I$ ): calculates the
      inverse of Matrix $M$ by LU decomposition, populating Matrix $I$
      with the result, and returns a reference to Matrix $I$
\item $M$.inverse(): calculates the inverse of Matrix $M$ by LU
      decomposition, and returns a new Matrix with the result
\item $M$.inverse( Matrix$<$ double $>$\& $I$, double\& $d$ ):
      calculates the inverse of Matrix $M$ by LU decomposition,
      populating Matrix $I$ with the result, returns the determinant
      in the passed argument $d$, and returns a reference to Matrix
      $I$
\item $M$.inverse( double\& $d$ ): calculates the inverse of Matrix
      $M$ by LU decomposition, returns the determinant in the passed
      argument $d$, and returns a new Matrix with the result
\item $M$.determinant(): returns the determinant of Matrix $M$ by LU
      decomposition
\item $M$.inverse( Matrix$<$ double $>$\& $I$, unsigned $choice$ ):
      allows the user to force type of inverse calculation by setting
      $choice$ to ``GaussJordan,'' (default is LU decomposition)
      populating Matrix $I$ with the result, and returns a reference
      to Matrix $I$
\item $M$.inverse( unsigned $choice$ ): allows the user to force type
      of inverse calculation by setting $choice$ to ``GaussJordan''
      (default is LU decomposition), and returns a reference to the
      new Matrix
\end{itemize}

In the case of a singular Matrix, these methods throw exception
Matrix$<$ double $>$::Singular\& $e$, where $e$.what() is ``singular
Matrix''.

\subsubsection{Example}
\begin{verbatim}
#include <iostream>
#include <string>
#include "matrix.h"

string inverse_filename;
cout << "The file name for Matrix inverse is ";
cin >> inverse_filename;
Matrix< double > I( true, inverse_filename ), J, K;
double determinant;

try
{
   // Test populating a passed Matrix
   J.resize( I.rows(), I.cols() ); // resize receiving Matrix to input
   Matrix
   cout << "The inverse of " << I << "is " << I.inverse( J ) << endl;

   // Test returning a new Matrix
   K = I.inverse();
   cout << "The inverse of " << I << "is " << K << endl;

   // Test populating a passed Matrix, and returning a determinant
   cout << "The inverse of " << I << "is " << I.inverse( J,
   determinant )
      << endl;
   // Note g++ requires this to be on a 2nd line, but MSVC++ does not:
   cout << "The determinant is " << determinant << endl;

   // Test returning a new Matrix with a determinant
   K.clear();
   K = I.inverse( determinant );
   cout << "The inverse of " << I << "is " << K
      << endl << "The determinant is " << determinant << endl;

   // Test returning just the determinant
   cout << "That determinant again is " << I.determinant() << endl;

   // Test populating a passed Matrix, and forcing type of inverse
   procedure
   cout << "The Gauss Jordan inverse of " << I << "is "
      << I.inverse( J, GaussJordan ) << endl;

   // Test returning a new Matrix and forcing type of inverse
   procedure
   K.clear();
   K = I.inverse( GaussJordan );
   cout << "The Gauss Jordan inverse of " << I << "is " << K << endl;

   // Examine the difference between LU decomposition and Gauss Jordan
   cout << "The difference between LU decomposition and Gauss Jordan
   inverse of "
      << I << "is " << I.inverse() - I.inverse( GaussJordan ) << endl;
}
catch ( Matrix< double >::Singular& e )
{
   cout << "Can't do it: " << e.what() << endl;
} \end{verbatim}

\subsubsection{Notes}

{\bf Use a try block with these methods to catch the Singular
exception.}  If you don't, your program may bomb out if a singular
Matrix is encountered.

These methods are only coded for Matrices of type $<$ double $>$.
Matrices to invert must be square, e.g. dimensions must be $Col \times
Col$, where $Col$ is the number of columns in the Matrix. If the
Matrix to be populated with the result is passed as an argument, it's
dimensions must be the same as the Matrix object.

\subsection{Eigenvalues}

The calculation of the eigenvalues of a symmetric matrix is the only
current instance of using pre-existing code and libraries, as these
calculations can be quite intensive, and are currently beyond the
scope of the current program.  Gaurav Bansal ported the eigenvalue
routines for a symmetric matrix from the Gnu Scientific Library (GSL)
in 2003, recoding those portions of GSL to integrate with the matrix
specification in neUROn2++.  {\em \bf \beginred As a result, do not
expect GSL functionality beyond eigenvalue routines in the neUROn2++
environment. \endcolor \bf \em}

The best way to understand calling eigenvalue routines in neUROn2++ is
to study the given example.  A GSL matrix is first created from a
neUROn2++ matrix using $gsl\_matrix\_view$, as is a GSL vector using
$gsl\_vector$. A GSL workspace is then created using
$gsl\_eigen\_symm\_workspace$ for symmetric matrix eigenvalue
calculations, and the eigenvalues are calculated using
$gsl\_eigen\_symm$.  Finally, the GSL vector is returned to an STL
vector containing the eigenvalues using $gsl\_vector\_ptr$, and memory
is freed using $gsl\_eigen\_symm\_free$.

For those interested in GSL, it may be found at
\href{http://sources.redhat.com/gsl/}{http://sources.redhat.com/gsl/}
 as of the time of this writing (August 11, 2003).

\subsubsection{Example}
\begin{verbatim}
#include "gsl_eigen.h"

// Create a symmetric matrix using our matrix class, and fill with
random values
Matrix< double > m1( 4, 4 );
m1.random( 5 );

// Check to see that we perform eigenvalue routines to symmetric
routines only
assert ( m1.rows() == m1.cols() );

// Create a gsl matrix from our matrix--begin function is cautioned to
use but
//    here we need the value as double*
gsl_matrix_view m = gsl_matrix_view_array( m1.begin(), 4, 4 );

// Create a gsl vector to store all the eigenvalues
gsl_vector *eval = gsl_vector_alloc( 4 );

// Create gsl workspace for eigenvalue calculations
gsl_eigen_symm_workspace *w = gsl_eigen_symm_alloc( 4 );

// Calculate eigenvalues
gsl_eigen_symm( &m.matrix, eval, w );

// Create vector< double > from gsl vector
vector< double > eigenv( gsl_vector_ptr( eval, 0 ), gsl_vector_ptr(
eval, 4 ) );

// Free workspace
gsl_eigen_symm_free( w );

// Output results
cout << "Matrix: " << m1 << endl << "Eigenvalues are: " << eigenv <<
endl;
cout << "Maximum value = " << maxabs( eigenv ) << endl;
cout << "Minimum value = " << minabs( eigenv ) << endl;
\end{verbatim}

\subsubsection{Notes}

All of the files needed are contained in the directory ``gsl''.  If
you are working in a pre-existing environment, you do not need to be
concerned about searching for directory paths, they should be there
already.  However, if you are building a new environment in g++,
you'll need to add search functionality in your Makefile.  An example
Makefile which searches for code inside of an ``msvc'' directory, and
the gsl subdirectory within it is:

\begin{verbatim}
vpath %.cpp msvc msvc/gsl
vpath %.h msvc msvc/gsl
vpath %.c msvc msvc/gsl

PROJ = neuron
CPP = g++
CPPFLAGS = -g -O2 -Wall
LDFLAGS = -lm

OBJS = $(PROJ).o ...

default: $(PROJ)

$(PROJ): $(OBJS)
    $(CPP) $(CPPFLAGS) $(OBJS) -Imsvc -Imsvc/gsl -o $(PROJ) $(LDFLAGS)
        -Lmsvc -Lmsvc/gsl -llibgsl -llibgslcblas

$(PROJ).o: $(PROJ).cpp
    $(CPP) $(CPPFLAGS) -Imsvc/gsl -o $@ -c $^
...
\end{verbatim}
%$

If you're using a pre-existing MSVC project, the following settings
should already be present in your .dsw workspace file.  However, for a
new MSVC project, you must instruct it to look for the gsl include
path, the link path, and the linked libraries libgsl.lib
libgslcblas.lib.  Choose the ``Project$\rightarrow$Settings'' menu
choice, then set the following:

\begin{itemize}
\item C/C++ tab
\begin{itemize}
\item Category: Preprocessor
\begin{itemize}
\item Additional Include Directories:
\begin{itemize}
\item gsl
\end{itemize}
\end{itemize}
\end{itemize}
\end{itemize}

\begin{itemize}
\item Link tab
\begin{itemize}
\item Category: Input
\begin{itemize}
\item Additional Library Path:
\begin{itemize}
\item gsl
\end{itemize}
\end{itemize}
\end{itemize}
\end{itemize}

\begin{itemize}
\item Link tab
\begin{itemize}
\item Category: Input
\begin{itemize}
\item Object/Library Modules: add
\begin{itemize}
\item libgsl.lib libgslcblas.lib
\end{itemize}
\end{itemize}
\end{itemize}
\end{itemize}

\section{An example of the Matrix and vector classes}

The following example constructs a simple neural network with 3 hidden
nodes on one layer to train the exclusive-or problem.  The file
``xor.set'' contains the exclusive-or dataset:

\begin{verbatim}
-0.9 -0.9 0.1
-0.9 0.9 0.9
0.9 -0.9 0.9
0.9 0.9 0.1
\end{verbatim}

The following is the program:\footnote{References are to James
A. Freeman and David M. Skapura, Neural Networks, Algorithms,
Applications and Programming Techniques, Addison-Wesley, 1991, ISBN
0-201-51376-5}

\begin{verbatim}
#include <iostream>
#include <iomanip>
#include <string>
#include <math.h>

#include "matrix.h"
#include "vector_ops.h"
#include "function_defs.h"

using namespace std;

int main()
{
   // Assumes a file named xor.set
   string filename = "xor.set";

   // Load the XOR dataset from the file xor.set
   Matrix< double > dataset( filename, 3 ); // it has 3 columns

   // Get the input Matrix In from the dataset as a submatrix
   Matrix< double > In = dataset.submatrix( 0, 3, 0, 1 );

   // Get the output vector y from the dataset's last column
   vector< double > y = dataset.col( 2 );

   // Make the 3 x 2 hidden weight Matrix hW and the hW update Matrix
   hWup
   Matrix< double > hW( 3, 2 ), hWup( 3, 2 );
   hW.random( 0.5 ); // randomize hW

   // Make the 2 element vector for the inputs I
   // and the 3 element vector hidden output hO
   // and the 3 element vector output weights oW
   vector< double > I( 2 ), hO( 3 ), oW( 3 );
   nvec::random( oW, 0.5 ); // randomize oW

   // o is the neural net's output, o_err the output error term,
   // eta is the learning rate
   double o, o_err, eta = 0.05;

   // h_err is the vector with hidden error terms
   vector< double > h_err( 3 );

   for ( unsigned iteration = 0; iteration < 1000000; iteration++ )
   {
      for ( unsigned example = 0; example < 4; example++ )
      {
         // Get a single row from the input Matrix
         // (Freeman & Skapura p. 102, #1)
         In.row( example, I );

         // Take the dot product of the hidden weights Matrix hW and
         the
         // transpose of a row of the input Matrix I, and apply the
         // sigmoidal function to the resulting vector
         // (Freeman & Skapura p. 102, #2 & #3)
         func( hW.dotprod( I, hO ), sigmoidal(), hO );

         // Take the dotproduct of the hidden output vector and the
         output
         // weights vector, and apply the sigmoidal function to the
         result
         // (Freeman & Skapura p. 102, #4 & #5)
         o = sigmoidal()( dotprod( hO, oW ) );

         // Calculate the error term for the output unit
         // (Freeman & Skapura p. 102, #6)
         o_err = ( y[ example ] - o ) * d_sigmoidal()( o );

         // Calculate the error terms for the hidden units
         // (Freeman & Skapura p. 102, #7)
         // Note that using func which takes the output vector as the
         // last argument, and *= instead of *, is the most efficient
         way
         ( func( hO, d_sigmoidal(), h_err ) *= o_err ) *= oW;

         // Update the output weights, note again the use of *= for
         efficiency
         // (Freeman & Skapura p. 102, #8)
         oW += ( hO *= ( eta * o_err ) );

         // Update the hidden weights
         // (Freeman & Skapura p. 102, #9)
         hW += ( hWup.outprod( h_err, I ) *= eta );
      }

      // Every 1000 iterations, print out the error
      if ( iteration % 1000 == 0 )
         cout << setw( 8 ) << iteration << ": LMS error = " <<
         0.5 * ( y[ 3 ] - o ) * ( y[ 3 ] - o ) << endl;
   }

   return 0;
}
\end{verbatim}

\chapter{Utility methods}

The following are contained in the header file ``utility.h'', and
contain global utility methods.  These methods throw exception
util::utilErr\& e, where e.what() specifies the module from which the
error originated, and the nature of the error.

\section{Random number generator}

{\bf namespace util}

The following utilities seed or generate random numbers of double
precision:

\begin{itemize}
\item d\_random(): seeds the random number generator with an internal
      timestamp
\item d\_random( double $n$ ): returns a double precision random
      number between $-n$ and $+n$
\item d\_random( double $lower$, double $upper$ ): returns a double
      precision random number between $lower$ and $upper$
\end{itemize}

\subsection{Examples}

\begin{verbatim}
util::d_random(); // seed the generator

double x, y;

x = util::d_random( 0.5 ); // Generates a random number between -0.5
and +0.5

y = util::d_random( 1.0, 10.0 ); // Between 1.0 and 10.0
\end{verbatim}

\section{Timestamp}

The timestamp( unsigned $x$ ) utility inserts a timestamp of the form
HHHH:MM:SS (where H is hour, M is minute, S is second) into the output
stream, given $x$ seconds.

\subsection{Example}
\begin{verbatim}
unsigned start = time( 0 ); // gets the current time in seconds

// Do something for a while
for ( i = 0; i < 100000; i++ )
   cout << ".";

// Calculate the elapsed time
unsigned increment = time( 0 ) - start;

// Print out how long it took
cout << endl << "That took " << timestamp( increment ) << endl;
\end{verbatim}

\section{Round}

{\bf namespace util}

The following method will round a number to a specified number of
significant digits:

\begin{itemize}
\item round( double $x$, unsigned $d$ ): rounds the passed value $x$
      to $d$ significant digits, returns the rounded value
\end{itemize}

\subsection{Example}
\begin{verbatim}
double x;
unsigned d;

cout << "x = ";
cin >> x;
cout << "d = ";
cin >> d;
cout << "round( x, d ) = " << util::round( x, d ) << endl;
\end{verbatim}

\subsection{Notes}
Zero significant digits is not allowed.  Rounds positive numbers up,
and negative numbers down, e.g. util::round( 30.9999, 2 ) = 31, and
util::round( -30.9999, 2 ) = -31.


\section{Queries with bounds checking}

{\bf namespace util}

The {\em askI and askD} methods allow user queries with bounds
checking.  Specifications include:

\begin{itemize}
\item askI( string\& $query$, int $lower$ ) which queries the user
      with the string $query$ repeatedly if necessary until the user
      enters an integer value $>= lower$, and returns the value
\item askI( string\& $query$, int $lower$, int $upper$ ) which queries
      the user with the string $query$ repeatedly if necessary until
      the user enters an integer value $>= lower$ and $<= upper$, and
      returns the value
\item askD( string\& $query$, double $lower$ ) which queries the user
      with the string $query$ repeatedly if necessary until the user
      enters a value $>= lower$, and returns the value
\item askD( string\& $query$, double $lower$, double $upper$ ) which
      queries the user with the string $query$ repeatedly if necessary
      until the user enters a value $>= lower$ and $<= upper$, and
      returns the value
\end{itemize}

\subsection{Example}
\begin{verbatim}
int i, h;
double d;

i = util::askI( "How many input nodes? ", 1 );
h = util::askI( "How many hidden nodes? ", 1, 10 );
d = util::askD( "What is the learning rate? ", 0, 1 );
\end{verbatim}

\section{Query for yes/no answer}

{\bf namespace util}

The {\em askYesNo} method allows user queries with yes/no answers.
Specifications include:

\begin{itemize}
\item askYesNo( string\& $query$ ) queries the user with the string
      $query$ repeatedly if necessary until the user enters ``yes'',
      ``y'', ``no'' or ``n'', and returns boolean true for yes, false
      for no
\end{itemize}

\subsection{Example}
\begin{verbatim}
bool yesNoAnswer;
yesNoAnswer = util::askYesNo( "Shall we clean the closet? " );
\end{verbatim}

\section{Query for existing file name}

{\bf namespace util}

The {\em getGoodFile} method tests for a file name to see if it
actually exists, then queries for an existing file name repeatedly
until a file that actually exists is entered, and returns that file
name as a {\em string}:

\begin{itemize}
\item getGoodFile( string\& $filename$ ) tests $filename$ to see if it
      actually exists, then if it doesn't, queries the user for a file
      name repeatedly if necessary until the user enters an existing
      file name, and returns it as a string.
\end{itemize}

\subsection{Example}
\begin{verbatim}
#include <iostream>
#include <fstream>
#include <string>
#include "utility.h"

string filename; // for the name of a loaded file

cout << "What is the name of your data file (it must exist)? ";
cin >> filename;

// Open the file as read-only, and associate it with the 'label'
inputFile
ifstream inputFile( ( util::getGoodFile( filename ) ).c_str(), ios::in
);
\end{verbatim}

\section{Remove a carriage return}

{\bf namespace util}

The {\em chopEndl} method removes ASCII 13, generally a carriage
return, from the end of a string if it exists, and returns a reference
to the string.

\begin{itemize}
\item chopEndl( string\& $stringObj$ ) removes the last character from
      $stringObj$ if it is ASCII 13, and returns a reference to
      $stringObj$.
\end{itemize}

\subsection{Example}
\begin{verbatim}
#include <iostream>
#include <fstream>
#include <string>
#include "utility.h"

string lineString; // string to hold the 1st line of file

// Open the file "test.txt" for input
ifstream loadfile( "test.txt", ios::in );

getline( loadfile, lineString ); // get the 1st line of file
util::chopEndl( lineString ); // remove <cr> from end of string if
exists

loadfile.close(); // close the input file
\end{verbatim}

\section{Parse a variable definition string}\label{parsefx}

{\bf namespace util}

\begin{itemize}
\item variable\_parse( string\& $stringObj$ ) parses $stringObj$ for
      variable representations from nodes (see below), and returns a
      vector of vector of unsigned representing the nodes.
\item variable\_parse( ifstream\& $fileObj$ ) parses the file
      $fileObj$ for variable representations from nodes (see below),
      and returns a vector of vector of unsigned representing the
      nodes.
\end{itemize}

The {\em variable\_parse} method parses a string passed as the
argument to determine variable representation from nodes, and returns
a {\em new} vector of vector of unsigned.  In the string,
\begin{itemize}
\item {\bf ;} delimits variables
\item {\bf ,} delimits nodes
\item {\bf -} specifies a range of nodes
\item spaces are ignored
\end{itemize}

For example, passed string ``0-3,5; 4,8; 6,7'' returns vector of
vector of unsigned (each line is a vector):
\begin{verbatim}
0 1 2 3 5
4 8
6 7
\end{verbatim}

util::utilErr exceptions thrown by this method include e.what()
\begin{itemize}
\item ``util variable\_parse error: bad character found at string
      position {\em position} ( {\em character} )''
\item ``util variable\_parse error: too many hyphens''
\item ``util variable\_parse error: hyphen without a number on both
      ends ''
\item ``util variable\_parse error: 1st number in range greater than
      2nd ''
\item ``util variable\_parse error: numbers in range are the same''
\item ``util variable\_parse error: missing zero''
\item ``util variable\_parse error: duplicated positions''
\item ``util variable\_parse error: missing positions''
\end{itemize}

\subsection{Example}
\begin{verbatim}
#include "stdafx.h" // if you're using MSVC

// Because, apparently, vector< vector< unsigned > > symbol names
exceeds 255
//    characters in length, and MSCV will error with that (ugly!)
#ifdef _MSC_VER	// MSVC
#pragma warning (disable: 4786)
#endif

#include <iostream>
#include <fstream>
#include <vector>
#include <string>

#include "vector_ops.h"
#include "utility.h"

using namespace std;

int main()
{
   // The returned structures
   vector< vector< unsigned > > parse_numbers, parse2;

   string parse_test = "0-3,5; 4,8; 6,7"; // a test string

   // Parse the string
   try
   {
      parse_numbers = util::variable_parse( parse_test );

      // Then print out the returned structure, a vector at a time
      cout << "Here it is: " << endl;
      unsigned pCount;
      for ( pCount = 0; pCount < parse_numbers.size(); pCount++ )
         cout << parse_numbers[ pCount ] << endl;
   }
   catch ( util::utilErr& e )
   {
      cout << e.what() << endl;
   }

   // Do the same from a file
   string reply;
   cout << "What is the name of the file? ";
   cin >> reply;
   ifstream inputFile( ( util::getGoodFile( reply ) ).c_str(), ios::in
   );

   try
   {
      parse2 = util::variable_parse( inputFile );
      cout << "Here it is: " << endl;
      for ( pCount = 0; pCount < parse2.size(); pCount++ )
         cout << parse2[ pCount ] << endl;
   }
   catch ( util::utilErr& e )
   {
      cout << e.what() << endl;
   }

   return 0;
}
\end{verbatim}

\subsubsection{Notes}
If you're using MSVC, you {\em must} include in your code the
following, as MSVC will give you error warnings for each vector$<$
vector$<$ unsigned $>$ $>$ reference:
\begin{verbatim}
#ifdef _MSC_VER	// MSVC
#pragma warning (disable: 4786)
#endif
\end{verbatim}

{\bf variable\_parse( ifstream\& $fileObj$ ) will close and clear the
file for you.}  Parsing a file will support multiple line definitions
of variables by concatenating the lines prior to parsing the resulting
string.  However, MSVC places a limit of 2,048 characters per string,
and if compiling with MSVC, definitions containing more than 2,048
characters will not be supported.

\section{Remove a variable from a vector of vector of unsigned data
structure}\label{rm_var}

{\bf namespace util}

\begin{itemize}
\item remove\_variable( vector$<$ vector$<$ unsigned $>$ $>$\&
	$variables$, unsigned $n$, vector$<$ vector$<$ unsigned $>$
	$>$\& $result$ ) removes the subvector representing variable
	$n$ from the vector of vector of unsigned data structure
	$variables$, populates $result$ as a vector of vector of
	unsigned data structure representing the result, returns a
	reference to the resulting data structure.
\item remove\_variable( vector$<$ vector$<$ unsigned $>$ $>$\&
	$variables$, unsigned $n$ ) as above, except returns a new
	vector of vector of unsigned data structure with the result.
\end{itemize}

The {\em remove\_variable} method removes a single subvector,
representing a single variable, from a vector of vector of unsigned
data structure which represents a set of variables.  As neUROn2++
variables may contain more than 1 input node (e.g. a variable may
contain 2 nodes, the value of the variable itself, and whether or not
that variable was present in the dataset), this method is used to
remove a single variable from a set of variables.

\subsection{Example}
\begin{verbatim}
// Create a vector of vector of unsigned with some variables
string parse_test = "0-3,5; 4,8; 6,7";
vector< vector< unsigned > > parse_numbers, parse_new, parse_new2;
parse_numbers = util::variable_parse( parse_test );

// Test remove_variable
parse_new = parse_numbers;
unsigned position;
position = util::askI( "What is the position to remove? ", 0 );
util::remove_variable( parse_numbers, position, parse_new );
cout << "After I've removed variable " << position << ": " << endl;
cout << "Here is the original: " << parse_numbers << endl;
cout << "Here is the result: " << parse_new << endl;

parse_new2 = util::remove_variable( parse_numbers, position );
cout << "Making a new data structure : " << endl;
cout << "Here is the original: " << parse_numbers << endl;
cout << "Here is the result: " << parse_new2 << endl;
\end{verbatim}

\subsubsection{Notes}
If you're using MSVC, you {\em must} include in your code the
following, as MSVC will give you error warnings for each vector$<$
vector$<$ unsigned $>$ $>$ reference:
\begin{verbatim}
#ifdef _MSC_VER	// MSVC
#pragma warning (disable: 4786)
#endif
\end{verbatim}

\chapter{Statistical methods}

The following are contained in the header file ``stats.h'', and
implemented in ``stats.cpp'', and contain global statistical methods.
Those that are not classes are contained within namespace ``stats'',
and throw exception stats::statsErr\& e, where e.what() specifies the
module from which the error originated, and the nature of the error.

\section{Log gamma}

{\bf namespace stats}

The following method returns the natural logarithm of the gamma
function $\log \Gamma(a)$\footnote{Adapted from W. H. Press et. al,
Numerical Recipes in C, 2nd edition, Cambridge University Press,
Cambridge, 1992, p. 214} where \index{Numerical Recipes in C}

\begin{equation}
\Gamma(z) = \int_0^\infty t^{z-1}e^{-t}dt
\end{equation}

\begin{itemize}
\item gammln( double $a$ ): returns the double value $\log \Gamma(a)$
\end{itemize}

\subsection{Example}
\begin{verbatim}
double x;
cout << "For log gamma, x = ";
cin >> x;
cout << "Log gamma " << x << " = " << stats::gammln( x ) << endl;
\end{verbatim}

\section{Incomplete gamma by fraction}

{\bf namespace stats}

The following method returns the incomplete gamma function
\begin{equation}\label{incmpgammfx}
P(a,x) \equiv \frac{\gamma(a,x)}{\Gamma(a)}
\equiv \frac{1}{\Gamma(a)}\int_0^x e^{-t}t^{a-1}dt \quad (a>0)
\end{equation}

by continued fraction development:\footnote{Adapted from W. H. Press
et. al, Numerical Recipes in C, 2nd edition, Cambridge University
Press, Cambridge, 1992, p. 219} \index{Numerical Recipes in C}

\begin{equation}
\Gamma(a,x) = e^{-x}x^a\left(\frac{1}{x+\frac{1-a}{1+\frac{1}{x+
\frac{2-a}{1+\frac{2}{x+\cdots}}}}}\right) \quad (x>0)
\end{equation}

\begin{itemize}
\item gcf( double $a$, double $x$ ): returns the double value
$\Gamma(a,x)$
\end{itemize}

\subsection{Example}
\begin{verbatim}
double a, x;
cout << "For gcf, a = ";
cin >> a;
cout << "x = ";
cin >> x;
try {
     cout << "gcf( " << a << " , " << x << " ) = " << stats::gcf( a, x
     );
} catch ( stats::statsErr& e ) {
     cout << e.what();
}
cout << endl;
\end{verbatim}

\section{Incomplete gamma by series}

{\bf namespace stats}

The following method returns the incomplete gamma function
(\hyperref[incmpgammfx]{equation \ref{incmpgammfx}}) by series:
\footnote{Adapted from W. H. Press et. al, Numerical Recipes in C, 2nd
edition, Cambridge University Press, Cambridge, 1992, pp. 218--219}
\index{Numerical Recipes in C}

\begin{equation}
\gamma(a,x) = e^{-x}x^a \sum_{n=0}^\infty \frac{\Gamma(a)}{\Gamma(a+1+
n)}x^n
\end{equation}

\begin{itemize}
\item gser( double $a$, double $x$ ): returns the double value
$\gamma(a,x)$
\end{itemize}

\subsection{Example}
\begin{verbatim}
double a, x;
cout << "For gser, a = ";
cin >> a;
cout << "x = ";
cin >> x;
try {
     cout << "gser( " << a << " , " << x << " ) = " << stats::gser( a,
     x );
} catch ( stats::statsErr& e ) {
     cout << e.what();
}
cout << endl;
\end{verbatim}

\section{Incomplete gamma composite}

{\bf namespace stats}

The following method returns the incomplete gamma function
(\hyperref[incmpgammfx]{equation \ref{incmpgammfx}}) by either
continued fraction or series, based on speed of convergence:\footnote{
Adapted from W. H. Press et. al, Numerical Recipes in C, 2nd edition,
Cambridge University Press, Cambridge, 1992, p. 218}
\index{Numerical Recipes in C}

\begin{itemize}
\item gammp( double $a$, double $x$ ): returns the double value
$P(a,x)$
\end{itemize}

\subsection{Example}
\begin{verbatim}
double a, x;
cout << "For gammp, a = ";
cin >> a;
cout << "x = ";
cin >> x;
try {
     cout << "gammp( " << a << " , " << x << " ) = " << stats::gammp(
     a, x );
} catch ( stats::statsErr& e ) {
     cout << e.what();
}
cout << endl;
\end{verbatim}

\section{Incomplete gamma complement}

{\bf namespace stats}

The following method returns the complement of the incomplete gamma
function
\begin{equation}\label{gammQlbl}
Q(a,x) \equiv 1-P(a,x)
\end{equation}

(see \hyperref[incmpgammfx]{equation \ref{incmpgammfx}} for the
definition of $P(a,x)$): \footnote{Adapted from W. H. Press et. al,
Numerical Recipes in C, 2nd edition, Cambridge University Press,
Cambridge, 1992, p. 218} \index{Numerical Recipes in C}

\begin{itemize}
\item gammq( double $a$, double $x$ ): returns the double value
$Q(a,x)$
\end{itemize}

\subsection{Example}
\begin{verbatim}
double a, x;
cout << "For gammq, a = ";
cin >> a;
cout << "x = ";
cin >> x;
try {
     cout << "gammq( " << a << " , " << x << " ) = " << stats::gammq(
     a, x );
} catch ( stats::statsErr& e ) {
     cout << e.what();
}
cout << endl;
\end{verbatim}

\section{Beta function}

{\bf namespace stats}

The following method returns the beta function $B(z,w)$
\footnote{Adapted from W. H. Press et. al, Numerical Recipes in C, 2nd
edition, Cambridge University Press, Cambridge, 1992, pp. 215--216}
where \index{Numerical Recipes in C}

\begin{equation}\label{Blbl}
B(z,w) = B(w,z) = \int_0^1t^{z-1}(1-t)^{w-1}dt = \frac{\Gamma(z)\Gamma
(w)}{\Gamma(z+w)}
\end{equation}

\begin{itemize}
\item beta( double $z$, double $w$ ): returns the double value
$B(z,w)$
\end{itemize}

\subsection{Example}
\begin{verbatim}
double z, w;
cout << "For beta function, z = ";
cin >> z;
cout << "w = ";
cin >> w;
cout << "Beta = " << stats::beta( z, w ) << endl;
\end{verbatim}

\section{Incomplete Beta function}

{\bf namespace stats}

The following method returns the incomplete beta function $I_x(a,b)$
\footnote{Adapted from W. H. Press et. al, Numerical Recipes in C, 2nd
edition, Cambridge University Press, Cambridge, 1992, pp. 226--228}
where \index{Numerical Recipes in C}

\begin{equation}\label{incompBlbl}
I_x(a,b) \equiv \frac{B_x(a,b)}{B(a,b)} \equiv \frac{1}{B(a,b)}\int_0^
xt^{a-1}(1-t)^{b-1}dt
\end{equation}

Note that the continued fraction representation is also available, but
as a utility function should not be necessary to call directly:

\begin{equation}
I_x(a,b) = \frac{x^a(1-x)^b}{aB(a,b)} = \left(\frac{1}{1+\frac{d_1}{1+
\frac{d_2}{1+\cdots}}}\right)
\end{equation}

where

\begin{equation}
d_{2m+1} = -\frac{(a+m)(a+b+m)x}{(a+2m)(a+2m+1)}
\end{equation}

and

\begin{equation}
d_{2m} = \frac{m(b-m)x}{(a+2m-1)(a+2m)}
\end{equation}

\begin{itemize}
\item betai( double $a$, double $b$, double $x$ ): returns the double
value $I_x(a,b)$
\item betacf( double $a$, double $b$, double $x$ ): returns the double
value $I_x(a,b)$
      by continued fraction evaluation, used by betai
\end{itemize}

\subsection{Example}
\begin{verbatim}
double a, b, x;
cout << "For incomplete beta function, a = ";
cin >> a;
cout << "b = ";
cin >> b;
cout << "x = ";
cin >> x;
cout << "Incomplete beta = " << stats::betai( a, b, x ) << endl;
\end{verbatim}

\section{Error functions}

{\bf namespace stats}

The following methods returns the error function (commonly referred to
as {\em erf})
\footnote{Adapted from W. H. Press et. al, Numerical Recipes in C, 2nd
edition, Cambridge University Press, Cambridge, 1992, pp. 220--221}
and its complement
\index{Numerical Recipes in C}

\begin{equation}
erf(x) = \frac{2}{\sqrt{\pi}}\int_0^x e^{-t^2}dt = P(\frac{1}{2},x^2)
\end{equation}
\begin{equation}
erfc(x) \equiv 1 - erf(x) = \frac{2}{\sqrt{\pi}}\int_x^\infty e^{-t^2}
dt = Q(\frac{1}{2},x^2)
\end{equation}

\begin{itemize}
\item erff( double $x$ ): returns the double value $erf(x)$
\item erfc( double $x$ ): returns the double value $erfc(x)$
\item erffc( double $x$ ): returns the double value $erfc(x)$ by
Chebyshev fitting
\end{itemize}

\subsection{Example}
\begin{verbatim}
double x;
cout << "x = ";
cin >> x;
try {
     cout << "erf( " << x << " ) = " << stats::erff( x );
} catch ( stats::statsErr& e ) {
     cout << e.what();
}
cout << endl;
try {
     cout << "erfc( " << x << " ) = " << stats::erfc( x );
} catch ( stats::statsErr& e ) {
     cout << e.what();
}
cout << endl;
\end{verbatim}

\subsection{Notes}
Note the 2 {\em f}~s in erff, so that this is not confused with the
innate compiler function erf.

\section{Error functions inverses}

{\bf namespace stats}

The following methods return the inverse of the error function and the
inverse of the complement of the error function:\footnote{Algorithm
for $erfc^{-1}(y)$ by Takuya Ooura $\copyright$ 1996,
$erf^{-1}(y) = erfc^{-1}(1-y)$}

\begin{itemize}
\item invErff( double $y$ ): returns the double value $erf^{-1}(y)$
\item invErfc( double $y$ ): returns the double value $erfc^{-1}(y)$
\end{itemize}

\subsection{Example}
\begin{verbatim}
double y;
cout << "y = ";
cin >> y;
try {
     cout << "inverse erf( " << y << " ) = " << stats::invErff( y );
} catch ( stats::statsErr& e ) {
     cout << e.what();
}
cout << endl;
try {
     cout << "inverse erfc( " << y << " ) = " << stats::invErfc( y );
} catch ( stats::statsErr& e ) {
     cout << e.what();
}
cout << endl;
\end{verbatim}

\section{Gaussian distribution}

{\bf namespace stats}

The following method returns the unit Gaussian function

\begin{equation}
\varphi(z) = \frac{1}{\sqrt{2\pi}} e^{-z^2/2}
\end{equation}

where given mean $\mu$ and standard deviation $\sigma$ the following
transformation $x \rightarrow z$ is made:

\begin{displaymath}
z \equiv \frac{x-\mu}{\sigma}
\end{displaymath}

$x$, $\mu$ and $\sigma$ may also be passed, where the Gaussian
function
is then

\begin{equation}
\varphi(x) = \frac{1}{\sigma\sqrt{2\pi}} e^{-(x-\mu)^2/(2\sigma^2)}
\end{equation}

\begin{itemize}
\item gaussD( double $z$ ): returns the double value $\varphi(z)$
\item gaussD( double $x$, double $\mu$, double $\sigma$): returns
the double value $\varphi(x, \mu, \sigma)$
\end{itemize}

\subsection{Example}
\begin{verbatim}
double z, x, u, s;
cout << "for gaussD, z = ";
cin >> z;
cout << "gaussD = " << stats::gaussD( z ) << endl;
cout << "for gaussD, x = ";
cin >> x;
cout << "u = ";
cin >> u;
cout << "s = ";
cin >> s;
cout << "gaussD = " << stats::gaussD( x, u, s ) << endl;
\end{verbatim}

\section{z Area}

{\bf namespace stats}

The following method returns the area under a unit Gaussian
distribution

\begin{equation}
\Phi(z) \equiv \frac{1}{\sqrt{2\pi}} \int_0^z e^{-z^2/2} dz =
\frac{1}{2} erf(\frac{z}{\sqrt{2}})
\end{equation}

\begin{itemize}
\item Zarea( double $z$ ): returns the double value $\Phi(z)$
\end{itemize}

\subsection{Example}
\begin{verbatim}
double z;
cout << "for z Area, z = ";
cin >> z;
try {
     cout << "z Area = " << stats::Zarea( z );
} catch ( stats::statsErr& e ) {
     cout << e.what();
}
cout << endl;
\end{verbatim}

\section{z Area inverse}\label{invZarea}

{\bf namespace stats}

The following method returns the inverse of the area under a unit
Gaussian distribution

\begin{equation}
z = Z(A) = \Phi^{-1}(A) = \sqrt{2}erf^{-1}(2A - 1)
\end{equation}

\begin{itemize}
\item invZarea( double $A$ ): returns the double value
$Z(A) = \Phi^{-1}(A)$
\end{itemize}

\subsection{Example}
\begin{verbatim}
double A;
cout << "for invZarea, A = ";
cin >> A;
try {
     cout << "invZarea ( " << A << " ) = " << stats::invZarea( A ) <<
     endl;
} catch ( stats::statsErr& e ) {
     cout << e.what();
}
\end{verbatim}

\section{Population metrics}

The {\em Population} class provides statistical metrics such as mean,
variance, standard deviation, skew and kurtosis.  It's constructor
takes vector$<$ double $>$ as its argument, representing the
population from which to generate metrics.  Errors are caught by {\em
PopulationErr} (see below for examples).

\begin{itemize}
\item Population( vector$<$ double $>$ $x$ ) defines a population from
      vector$<$ double $>$ $x$ from which to generate statistical
      metrics
\end{itemize}

Accessors include:
\begin{itemize}
\item $P$.mean(): returns the mean
\item $P$.var(): returns the variance
\item $P$.std(): returns the standard deviation
\item $P$.skew(): returns the skew
\item $P$.kurtosis(): returns the kurtosis
\end{itemize}

\subsection{Example}
\begin{verbatim}
double numbers[ 4 ] = { 1.0, 2.0, 3.0, 4.0 }; // define the dataset
vector< double > dset( numbers, numbers + 4 );
Population P( dset ); // define the population from the dataset
cout << "average = " << P.mean() << endl;
cout << "variance = " << P.var() << endl;
cout << "standard deviation = " << P.std() << endl;
try {
   cout << "skew = " << P.skew() << endl;
   cout << "kurtosis = " << P.kurtosis() << endl;
} catch ( Population::PopulationErr& e ) {
   cout << e.what() << endl;
}
\end{verbatim}

\subsection{Notes}
The only allowable constructor takes vector$<$ double $>$ as an
argument, and there is no copy constructor or overloaded $=$ operator.
The vector$<$ double $>$ argument {\em must} contain 2 or more
elements.

\section{F and Student's t}

{\bf namespace stats}

The following methods return probabilities for the F-test that 2
sample
variances are different ($p$=1.0 if they are equal), and Student's t-
tests
that 2 samples are different.  All take 2 vectors of double as the
samples
$S_1$ and $S_2$.

\begin{itemize}
\item ftest( vector$<$ double $>$ $S_1$, vector$<$ double $>$ $S_2$ ):
      returns a double value with the probability for an F-test.
\item ttest( vector$<$ double $>$ $S_1$, vector$<$ double $>$ $S_2$ ):
      returns a double value with the probability for a Student's t-
      test
      with equal sample variance.
\item tutest( vector$<$ double $>$ $S_1$, vector$<$ double $>$ $S_2$
):
      returns a double value with the probability for a Student's t-
      test
      with unequal sample variance.
\item tptest( vector$<$ double $>$ $S_1$, vector$<$ double $>$ $S_2$
):
      returns a double value with the probability for a paired
      Student's
      t-test.
\end{itemize}

Note that\footnote{Adapted from W. H. Press et. al, Numerical
Recipes in C, 2nd edition, Cambridge University Press, Cambridge,
1992, pp. 228--229, 615--619}\index{Numerical Recipes in C}
\begin{itemize}
\item For the F-test,
\begin{equation}
\alpha = I_{\frac{\nu_2}{\nu_2+\nu_1 F}}
(\frac{\nu_2}{2},\frac{\nu_1}{2})
\end{equation}
      (see \hyperref[incompBlbl]{equation \ref{incompBlbl}}), where
      $F = \frac{\sigma_1^2}{\sigma_2^2}$ ($\sigma_1^2 > \sigma_2^2$),
      and $\nu_1 = N_1 - 1$ and $\nu_2 = N_2 - 1$
\item For Student's distribution, $\alpha=1-A(t|\nu)$, where
\begin{equation}
A(t|\nu)=\frac{1}{\sqrt{\nu}B(\frac{1}{2},\frac{\nu}{2})}
\int_{-t}^t(1+\frac{x^2}{\nu})^{\frac{\nu+1}{2}}dx
\end{equation}
      (see \hyperref[Blbl]{equation \ref{Blbl}}) and is related to the
      incomplete beta function by
 \begin{equation}
A(t|\nu)=1-I_{\frac{\nu}{\nu+t^2}} (\frac{\nu}{2},\frac{1}{2})
\end{equation}
      (see \hyperref[incompBlbl]{equation \ref{incompBlbl}})
\begin{itemize}
\item For Student's t with equal sample variance,
\begin{equation}
t = \frac{\mu_1-\mu_2}{s_D}
\end{equation}
where
\begin{equation}
s_D=\sqrt{\frac{\sum_{i}^{N_1}(x_i-\mu_1)^2+\sum_{i}^{N_2}(x_i-\mu_2)^
2}
{N_1+N_2-2}\left(\frac{1}{N_1}+\frac{1}{N_2}\right)}
\end{equation}
      and $\nu=N_1+N_2-2$
\item For Student's t with unequal sample variance,
\begin{equation}
t=\frac{\mu_1-\mu_2}{\sqrt{\frac{\sigma_1^2}{N_1}+\frac{\sigma_2^2}{
N_2}}}
\end{equation}
and
\begin{equation}
\nu=\frac{\left(\frac{\sigma_1^2}{N_1}+\frac{\sigma_2^2}{N_2}\right)^2
}
{\frac{\left(\frac{\sigma_1^2}{N_1}\right)^2}{N_1-1}+
\frac{\left(\frac{\sigma_2^2}{N_2}\right)^2}{N_2-1}}
\end{equation}
\item For paired Student's t,
\begin{equation}
t = \frac{\mu_1-\mu_2}{s_D}
\end{equation}
where
\begin{equation}
s_D=\sqrt{\frac{\sigma_1^2+\sigma_2^2-2Cov(x1,x2)}{N}}
\end{equation}
where
\begin{equation}
Cov(x1,x2) \equiv \frac{1}{N-1}\sum_{i=1}^{N}(x_{i1}-\mu_1)(x_{i2}-
\mu_2)
\end{equation}
      and $\nu=N-1$, where $N$ is the {\em number of pairs}.
\end{itemize}
\end{itemize}

\subsection{Example}
\begin{verbatim}
double d1[ 9 ] = { 3, 4, 5, 8, 9, 1, 2, 4, 5 };
double d2[ 9 ] = { 6, 19, 3, 2, 14, 4, 5, 17, 1 };
vector< double > v1( d1, d1 + 9 ), v2( d2, d2 + 9 );

cout << "F-test = " << stats::ftest( v1, v2 ) << endl;
cout << "Student's t for same variance = " << stats::ttest( v1, v2 ) <
< endl;
cout << "Student's t for unequal variance = " << stats::tutest( v1, v2
) << endl;
cout << "Paired Student's t = " << stats::tptest( v1, v2 ) << endl;
\end{verbatim}

\subsection{Notes}
Note that in all cases, sample (passed vector) sizes {\em must be}
$>1$, and
for the paired Student's t-test, they {\em must be} equal.

\section{Chi-square}

{\bf namespace stats}

The following method returns the probability $p$ for $df$ degrees of
freedom and Chi-square value $\chi^2$:

\begin{itemize}
\item pX2( unsigned $df$, double $\chi^2$ ): returns a double value
      with the probability.
\end{itemize}

Note that the probability $p$ for the Chi-square value $\chi^2$ with
$df$ degrees of freedom is $Q(\frac{df}{2},\frac{\chi^2}{2})$ (see
\hyperref[gammQlbl]{equation \ref{gammQlbl}}).

\subsection{Example}
\begin{verbatim}
double x; // chi-square value
unsigned df; // degrees of freedom

// Query the user for the value and degrees of freedom
cout << "For chi-square, df = ";
cin >> df;
cout << "X2 = ";
cin >> x;

// Then calculate the probability p
try {
     cout << "p = " << stats::pX2( df, x );
} catch ( stats::statsErr& e ) { // catch a stats error
     cout << e.what();
}
cout << endl;
\end{verbatim}

\section{Kolmogorov-Smirnov statistic}\label{ks_p}

{\bf namespace stats}

The following method returns the Kolmogorov-Smirnov
probability\index{Kolmogorov-Smirnov!Probability}
$Q_{KS}$\footnote{Adapted from W. H. Press et. al, Numerical
Recipes in C, 2nd edition, Cambridge University Press, Cambridge,
1992, pages 623-626}\index{Numerical Recipes in C}

\begin{equation}
Q_{KS}(\lambda) = 2\sum_{j=1}^{\infty}(-)^{j-1}e^{-2j^2\lambda^2}
\end{equation}

\begin{itemize}
\item double probks( double $\lambda$ ): returns a double value with
      $Q_{KS}(\lambda)$.
\end{itemize}

\subsection{Example}
\begin{verbatim}
double alam;
cout << "For Kolmogorov-Smirnov, lambda = ";
cin >> alam;
cout << "Q_KS " << alam << " = " << stats::probks( alam ) << endl;
\end{verbatim}

\section{Functions}

The {\em Funct} class provides functional manipulation such as root
finding, etc.  The {\em Funct} constructor takes as its arguments an
instance of the function (generally {\em *this}) followed by a
reference to the {\em member function of a class}, for example:

\begin{verbatim}
class TestIt {
public:
   TestIt(){}
   ~TestIt(){}
   Funct< TestIt > squarer(){ return Funct< TestIt >( *this, &
   TestIt::xsqrd ); }

private:
   double xsqrd( double x ){ return x * x - 4; }
};
\end{verbatim}

Note that {\em Funct} is a template class, and takes as its template
argument the name of the class containing the member function.  Errors
are caught by {\em functErr} (see below for examples).

\subsection{Bracketing}

The following methods bracket a function:\footnote{Adapted from W. H.
Press et.\ al, Numerical Recipes in C, 2nd edition, Cambridge
University Press, Cambridge, 1992, pp. 400--401}\index{Numerical
Recipes in C}
\begin{itemize}
\item F.mnbrak( double $x1$, double $x2$ ) takes boundary points $x1$
and $x2$ to bracket function $F$.  Accessors below return the
bracketed result, where $ax$, $bx$ and $cx$ are 3 sequential $x$-axis
points, with $fb$ = $f(bx)$ less than $fa$ = $f(ax)$ and $fc$ =
$f(cx)$.  All returned values are of type double.
\item F.mnbrak\_ax(): accessor returns $ax$
\item F.mnbrak\_bx(): accessor returns $bx$
\item F.mnbrak\_cx(): accessor returns $cx$
\item F.mnbrak\_fa(): accessor returns $fa$
\item F.mnbrak\_fb(): accessor returns $fb$
\item F.mnbrak\_fc(): accessor returns $fc$
\end{itemize}

\subsubsection{Example}
\begin{verbatim}
class TestIt {
public:
   TestIt(){}
   ~TestIt(){}
   void mnbrak( double x1, double x2 )
   {
      try
      {
         Funct< TestIt > f( *this, &TestIt::xsqrd );
         f.mnbrak( x1, x2 );
         cout << "ax = " << f.mnbrak_ax() << endl;
         cout << "bx = " << f.mnbrak_bx() << endl;
         cout << "cx = " << f.mnbrak_cx() << endl;
         cout << "fa = " << f.mnbrak_fa() << endl;
         cout << "fb = " << f.mnbrak_fb() << endl;
         cout << "fc = " << f.mnbrak_fc() << endl;
      }
      catch ( Funct< TestIt >::functErr& e ) { cout << e.what() <<
      endl; }
   }

private:
   double xsqrd( double x ){ return x * x - 4; }
};

int main()
{
   TestIt tester;
   double x1, x2;
   cout << "For testing mnbrak, x1 = ";
   cin >> x1;
   cout << "x2 = ";
   cin >> x2;
   tester.mnbrak( x1, x2 );

   return 0;
}
\end{verbatim}

\subsection{Find Minimum}

The following methods isolate the minimum of function $F$ using
Brent's method:\footnote{Adapted from W. H. Press et.\ al, Numerical
Recipes in C, 2nd edition, Cambridge University Press, Cambridge,
1992, pp. 402--405}\index{Numerical Recipes in C}
\begin{itemize}
\item F.brent( double $ax$, double $bx$, double $cx$, double $tol$ ):
takes initial bracketing triplet of abscissas $ax$, $bx$ and $cx$ and
tolerance $tol$ to isolate minimum of function $F$.
\item F.brent\_xmin(): returns the isolated minimum (type double) of
function $F$.
\end{itemize}

\subsubsection{Example}
\begin{verbatim}
class TestIt {
public:
   TestIt(){}
   ~TestIt(){}
   void brent( double a1, double b1, double c1, double tol )
   {
      try
      {
         Funct< TestIt > f( *this, &TestIt::xsqrd );
         cout << "brent = " << f.brent( a1, b1, c1, tol ) << endl;
         cout << "xmin = " << f.brent_xmin() << endl;
      }
      catch ( Funct< TestIt >::functErr& e ) { cout << e.what() <<
      endl; }
   }

private:
   double xsqrd( double x ){ return x * x - 4; }
};

int main()
{
   TestIt tester;
   double a, b, c, tol;
   cout << "For testing brent, a = ";
   cin >> a;
   cout << "b = ";
   cin >> b;
   cout << "c = ";
   cin >> c;
   cout << "tolerance = ";
   cin >> tol;
   tester.brent( a, b, c, tol );

   return 0;
}
\end{verbatim}

\subsection{Find Roots}

The following method finds roots of function $F$ using Brent's
method:\footnote{Adapted from W. H. Press et.\ al, Numerical Recipes
in C, 2nd edition, Cambridge University Press, Cambridge, 1992, pp.
359--362}\index{Numerical Recipes in C}
\begin{itemize}
\item F.zbrent( double $x1$, double $x2$, double $tol$ ): takes 2
$x$-axis bounds $x1$ and $x2$, and tolerance $tol$, and returns the
root of function $F$ between those bounds.
\end{itemize}

\subsubsection{Example}
\begin{verbatim}
class TestIt {
public:
   TestIt(){}
   ~TestIt(){}
   void zbrent( double x1, double x2, double tol )
   {
      try
      {
         Funct< TestIt > f( *this, &TestIt::xsqrd );
         cout << "zbrent = " << f.zbrent( x1, x2, tol ) << endl;
      }
      catch ( Funct< TestIt >::functErr& e ) { cout << e.what() <<
      endl; }
   }

private:
   double xsqrd( double x ){ return x * x - 4; }
};

int main()
{
   TestIt tester;
   double x1, x2, tol;
   cout << "For testing zbrent, x1 = ";
   cin >> x1;
   cout << "x2 = ";
   cin >> x2;
   cout << "tolerance = ";
   cin >> tol;
   tester.zbrent( x1, x2, tol );

   return 0;
}
\end{verbatim}

\section{Linear Regression}

The {\em XY} class performs least squares linear regression. The {\em
XY} constructors takes as their arguments vector$<$ double $>$s
containing the x- and y-coordinates, and if present, the x- and
y-standard deviations:\footnote{Adapted from W. H. Press et.\ al,
Numerical Recipes in C, 2nd edition, Cambridge University Press,
Cambridge, 1992, pp. 661--670}\index{Numerical Recipes in C}
\begin{itemize}
\item XY( vector$<$ double $>$\& $x$, vector$<$ double $>$\& $y$ ):
loads an XY dataset with x-coordinates $x$ and y-coordinates $y$
without standard deviations
\item XY( vector$<$ double $>$\& $x$, vector$<$ double $>$\& $y$,
vector $<$ double $>$\& $\sigma_x$): loads an XY dataset with
x-coordinates $x$, y-coordinates $y$, and $x$ standard deviations
$\sigma_x$
\item XY( vector$<$ double $>$\& $x$, vector$<$ double $>$\& $y$,
vector $<$ double $>$\& $\sigma_x$, vector $<$ double $>$\&
$\sigma_y$): loads an XY dataset with x-coordinates $x$, y-coordinates
$y$, $x$ standard deviations $\sigma_x$, and $y$ standard deviations
$\sigma_y$
\end{itemize}

The following accessors return results of least squares linear
regression on $XY$ dataset $D$ (all of type double):
\begin{itemize}
\item $D$.a(): returns coefficient $a$ (y-intercept) in least squares
linear fit $y = a + bx$
\item $D$.b(): returns coefficient $b$ (slope) in least squares linear
fit $y = a + bx$
\item $D$.x0(): returns x-intercept $x_0$ in least squares linear fit
$y = a + bx$
\item $D$.siga(): returns standard errors on $a$
\item $D$.sigb(): returns standard errors on $b$
\item $D$.chi2(): returns minimized chi-squared value of least squares
fit
\item $D$.q(): returns chi-squared probability of least squares fit
(small is a poor fit)
\item $D$.r(): returns correlation coefficient $r$
\end{itemize}

\subsection{Notes}
Errors are caught by {\em XYErr} (see below for examples).

\subsection{Examples}
\begin{verbatim}
double xs[ 5 ] = { 1.0, 2.0, 3.0, 4.0, 5.0 };
double ys[ 5 ] = { 1.0, 2.2, 3.0, 4.0, 6.0 };
double stdxs[ 5 ] = { 0.2, 0.4, 0.2, 0.2, 1 };
double stdys[ 5 ] = { 0.2, 0.3, 0.2, 0.2, 0.8 };

vector< double > xv( xs, xs + 5 ), yv( ys, ys + 5 ), stdx( stdxs,
stdxs + 5 ),
   stdy( stdys, stdys + 5 );

try {
   XY xyset( xv, yv );
   cout << "For dataset without errors, a = " << xyset.a() << endl;
   cout << "b = " << xyset.b() << endl;
   cout << "x0 = " << xyset.x0() << endl;
   cout << "siga = " << xyset.siga() << endl;
   cout << "sigb = " << xyset.sigb() << endl;
   cout << "chi2 = " << xyset.chi2() << endl;
   cout << "q = " << xyset.q() << endl;
   cout << "r = " << xyset.r() << endl << endl;
}
catch ( XY::XYErr& e ) { cout << e.what() << endl;

try {
   XY xyset2( xv, yv, stdx );
   cout << "For dataset with x errors only, a = " << xyset2.a() <<
   endl;
   cout << "b = " << xyset2.b() << endl;
   cout << "x0 = " << xyset2.x0() << endl;
   cout << "siga = " << xyset2.siga() << endl;
   cout << "sigb = " << xyset2.sigb() << endl;
   cout << "chi2 = " << xyset2.chi2() << endl;
   cout << "q = " << xyset2.q() << endl;
   cout << "r = " << xyset2.r() << endl << endl;
}
catch ( XY::XYErr& e ) { cout << e.what() << endl;

try {
   XY xyset3( xv, yv, stdx, stdy );
   cout << "For dataset with x & y errors, a = " << xyset3.a() <<
   endl;
   cout << "b = " << xyset3.b() << endl;
   cout << "x0 = " << xyset3.x0() << endl;
   cout << "siga = " << xyset3.siga() << endl;
   cout << "sigb = " << xyset3.sigb() << endl;
   cout << "chi2 = " << xyset3.chi2() << endl;
   cout << "q = " << xyset3.q() << endl;
   cout << "r = " << xyset3.r() << endl << endl;
}
catch ( XY::XYErr& e ) { cout << e.what() << endl; }
\end{verbatim}

\chapter{Object Design}\label{obj_chap}
\setcounter{figcount}{1}

The design of the network object pursues 2 objectives.  The first is a
well organized, object oriented representation of a generic network
object that is flexible and easily extended to specific network types.
In fact, the current design allows for all models, including
statistical ones.

The second objective of the network object design is that the code be
readable.  This was the logic of constructing efficient yet readable
matrix and vector operations, as described in Chapter
\ref{matrix_chapter}.

\hyperref[hierarchy]{Figure~\ref{obj_chap}.\arabic{figcount}}
demonstrates the object hierarchy.  Objects that are to be built are
enclosed in dashed lines. Objects that are built (although features
will be added) are shown in solid lines.

\begin{figure}[htb] \addtocounter{figcount}{1}
\begin{center}
\begin{picture}(4.2,4.2)
\put(0,0){\framebox(4.2,4.2)}
% for dvi
\put(0.1,0.1){\includegraphics[width=4in,height=4in]{objects.bmp}}
% for dvipdfm
\put(0.1,0.1){\special{pdf: epdf width 4in height 4in (objects.pdf)}}
\end{picture}
\end{center}
\caption{Object design}\label{hierarchy}
\end{figure}

\begin{itemize}
\item {\bf DataSet} contains and manipulates the data to be modeled.
\item {\bf TwoSet} is a specialized object designed to calculate
      metrics for {\em one} output {\em Models}, including
      sensitivity, specificity, predictive value positive, predictive
      value negative, and ROC area.
\item {\bf Model} is the base object. {\em Model} provides access to
      the dataset, and contains the virtual methods common to all
      modeling techniques.
\item {\bf Iterative} extends {\em Model}.  Neural computational
      algorithms are iterative, and thus use {\em Iterative} as a
      base.  {\em Iterative} manages methods and parameters related to
      iterating through a dataset.
\item {\bf Network} extends {\em Iterative}.  Algorithms with hidden
      layers, weights, learning rates, transfer functions and other
      general properties of neural networks are encapsulated in {\em
      Network}.
\item {\bf RegressNet} is a compound object including a {\em Network}
      object which implements stepwise regression.  See Chapter
      \ref{regresschap}.
\item {\bf BackProp} extends {\em Network}.  {\em BackProp} is a
      multilayer multiple output back-propagation neural network.
      \index{Backpropagation}
\item {\bf SimpleProp} also extends {\em Network}.  {\em SimpleProp}
      is a backpropagation neural network which has 1 output node, 1
      hidden node layer, and biases.\index{Backpropagation}
\item {\bf BareProp} also extends {\em Network}.  {\em BareProp} is an
      example {\em Network} object for illustrative purposes which has
      1 output node, 1 hidden node layer, and no biases.
      \index{Backpropagation}
\item {\bf Logistic} also extends {\em Network}.  {\em Logistic}
      implements binary logistic regression, which is mathematically
      equivalent to a single output node neural network with a bias,
      the sigmoidal transfer function, no hidden nodes, and the
      cross-entropy error function.\index{Logistic regression}
\item {\bf DFA} extends {\em Model}, and encodes discriminant function
      analysis.  {\bf QDFA} extends {\em DFA}, and encodes quadratic
      discriminant function analysis.  {\bf LDFA} extends {\em DFA},
      and encodes linear discriminant function analysis.
\end{itemize}

\section{Function definitions}

For clarity and ease of use, function definitions are placed in the
file ``function\_defs.h''. The implementation of function definitions
in this file may take a number of forms, depending on how the function
is to be used in the code.  The reader is strongly encouraged to study
``function\_defs.h'' to see the various forms functions may take.

\subsection{In-line functions}

The simplest form is the in-line function.  An example is that which
calculates the least mean squared error for a single output node:
\index{Error function!Least squares}

\begin{verbatim}
// LMS error for a single output, takes known value y and model's
//    guess o as arguments, returns LMS error
inline double singleLMSerror( double y, double o )
{
     return 0.5 * ( y - o ) * ( y - o );
}
\end{verbatim}

\subsection{Function classes}

For functions that are to be passed to {\em vector} or {\em Matrix}
objects using the {\bf func} method, classes must be constructed.
An example is that which implements the sigmoidal transfer function:
\index{Transfer function!Sigmoidal}

\begin{verbatim}
// The sigmoidal function
class sigmoidal : public unary_function< double, double >
{
     public:
          inline double operator()( const double& x )
          {
               return 1.0 / ( 1.0 + exp( -x ) );
          }
};
\end{verbatim}

\subsection{The errorFunction object}\label{errFunc}

A very important object defined in {\bf function\_defs.h} is {\em
errorFunction}, which handles the calculation of model error.  As
neural network sigmoidal output is defined as
\index{Transfer function!Sigmoidal}

\begin{displaymath}
\vec{\bf{o}} = \frac{1}{1+e^{-\vec{\bf{x}}}}
\end{displaymath}

calculation of the error function may depend on $\vec{\bf{x}}$ as well
as $\vec{\bf{o}}$ (see the {\bf special note about numerical precision
and the cross-entropy error} below),
\index{Error function!Cross-entropy} and thus {\em errorFunction}
methods may be passed $\vec{\bf{x}}$ as well as $\vec{\bf{o}}$.
{\em errorFunction} consists of

\begin{itemize}

\item {\em Public Methods}

\item errorFunction(double $y$, double $o$, double $x$, bool
      $errorType$): Constructor for a {\em single} output, which takes
      known value $y$, model output $o$, argument $x$ in sigmoidal
      transfer function $1/(1+e^{-x})$ and bool flag $errorType$,
      where\index{Transfer function!Sigmoidal}

\begin{itemize}

\item $errorType = 0$ denotes the least mean squared error
\index{Error function!Least squares}

\begin{displaymath}
E = \frac{1}{2} (y_i - o_i)^2
\end{displaymath}

\item $errorType = 1$ denotes the cross-entropy error
\index{Error function!Cross-entropy}

\begin{displaymath}
E = -y\log(o)-(1-y)\log(1-o)
\end{displaymath}

\end{itemize}

\item errorFunction(vector$<$ double $>$ $\vec{\bf y}$, vector$<$
      double $>$ $\vec{\bf o}$, vector$<$ double $>$ $\vec{\bf x}$,
      bool $errorType$): Constructor for {\em multiple} outputs, which
      takes known values $\vec{\bf y}$, model outputs $\vec{\bf o}$,
      arguments $\vec{\bf x}$ in sigmoidal transfer function
      $1/(1+e^{-\vec{\bf x}})$ and bool flag $errorType$, where
      \index{Transfer function!Sigmoidal}

\begin{itemize}

\item $errorType = 0$ denotes the least mean squared error for $m$
output nodes \index{Error function!Least squares}

\begin{displaymath}
E = \frac{1}{2} \sum_{i=1}^{m} (y_i - o_i)^2
\end{displaymath}

\item $errorType = 1$ denotes the cross-entropy error for $m$ output
nodes\index{Error function!Cross-entropy}

\begin{displaymath}
E = \sum_{i=1}^{m} -y\log(o)-(1-y)\log(1-o)
\end{displaymath}

\end{itemize}

\item double value(): Accessor returns numerical error result

\item bool boundsErr() Accessor returns true if numerical bounds error
      (e.g. $\log(0)$) (see the {\bf special note about numerical
      precision and the cross-entropy error} below)
      \index{Error function!Cross-entropy}

\end{itemize}

\subsubsection{Example}

\begin{verbatim}
// Calculate x-entropy error for single output and add to set error
//    which is previously defined as double setError (double y, double
o
//    and double x have also previously been calculated)
errorFunction E( y, o, x, 1 );
setError += E.value();

if ( E.boundsErr() ) // check for out of bounds error in log(o)
    cout << "ERROR: Numerical bounds error" << endl;
\end{verbatim}

\subsubsection{A special note about numerical precision and the
cross-entropy error}\index{Error function!Cross-entropy}

As the \hyperref[xentropyfunk]{cross-entropy error function} is
defined for network output $o$ and target value $y$:

\begin{displaymath}\index{Error function!Cross-entropy}
E = - y\log(o)-(1-y)\log(1-o)
\end{displaymath}

And for the sigmoidal transfer function,
\index{Transfer function!Sigmoidal}

\begin{displaymath}
o = \frac{1}{1+e^{-x}}
\end{displaymath}

Then for large positive ($>> 0$) and large negative ($<< 0$) $x$,
$o \rightarrow 1$, and $o \rightarrow 0$ respectively, and thus
$log(1-o) \rightarrow \infty$ and $log(o) \rightarrow \infty$. In
these cases, the following transformations are made (which require
those methods to be passed the sigmoidal transfer function argument
$x$ for the cross-entropy error.)
\index{Transfer function!Sigmoidal}
\index{Error function!Cross-entropy}

\subsubsection{Large negative $x$}

If $x << 0$,

\begin{displaymath}
o = \frac{1}{1+e^{-x}} \approx \frac{1}{e^{-x}} = e^x
\end{displaymath}

and thus

\begin{displaymath}
\log(o) = \log(e^x) = x
\end{displaymath}

since $o \rightarrow 0$,

\begin{displaymath}
log(1-o) \rightarrow log(1) = 0
\end{displaymath}

thus

\begin{displaymath}
E = - y\log(o)-(1-y)\log(1-o) \approx -yx - 0 = -yx
\end{displaymath}

\subsubsection{Large positive $x$}

\begin{displaymath}
1-o = 1-\frac{1}{1+e^{-x}} = \frac{1+e^{-x}}{1+e^{-x}}
- \frac{1}{1+e^{-x}} = \frac{e^{-x}}{1+e^{-x}}
\end{displaymath}

If $x >> 0$, since $1+e^{-x} \rightarrow 1$

\begin{displaymath}
1-o=\frac{e^{-x}}{1+e^{-x}} \approx e^{-x}
\end{displaymath}

thus

\begin{displaymath}
\log(1-o) \approx log(e^{-x}) = -x
\end{displaymath}

if $x >> 0$, $o \rightarrow 1$ and $-y\log(o) \rightarrow 0$,
thus

\begin{displaymath}
E = - y\log(o)-(1-y)\log(1-o) \approx 0 - (1-y)(-x) = (1-y)x
\end{displaymath}

\section{What you really need to know about DataSet and TwoSet}

All models use training and test sets.  The model builds itself on the
training set, and is tested on the test set.  {\em DataSet} is
designed for generating and manipulating training and test sets, and
{\em TwoSet} is a special object that is designed to calculate
accuracies for one-output models {\em only}.

All training and test sets are expressed by Matrices in the following
format, where $i$ is an input to the model, and $o$ is an output:

\begin{center}
\begin{tabular}{cccccccc}
$i_{11}$ & $i_{12}$ & $\ldots$ & $i_{1I}$ & $o_{11}$ & $o_{12}$ &
$\ldots$ & $o_{1O}$ \\
$i_{21}$ & $i_{22}$ & $\ldots$ & $i_{2I}$ & $o_{21}$ & $o_{22}$ &
$\ldots$ & $o_{2O}$ \\
$i_{31}$ & $i_{32}$ & $\ldots$ & $i_{3I}$ & $o_{31}$ & $o_{32}$ &
$\ldots$ & $o_{3O}$ \\
$\ldots$ & $\ldots$ & $\ldots$ & $\ldots$ & $\ldots$ & $\ldots$ &
$\ldots$ & $\ldots$ \\
$i_{N1}$ & $i_{N2}$ & $\ldots$ & $i_{NI}$ & $o_{N1}$ & $o_{N2}$ &
$\ldots$ & $o_{NO}$ \\
\end{tabular}
\end{center}

for $N$ examples (``exemplars'') in the set, $I$ inputs and $O$
outputs.  Although there are many highly useful features in {\em
DataSet}, a programmer building a {\em Model} need only know how to
access the training and test Matrices, whether they've been loaded
into the {\em DataSet} object, how many inputs and outputs there are,
whether the output is discrete (held to 0 or 1), and the number of
examples in each set.  These methods are shown in \beginred red
\endcolor in the {\em DataSet} specification.

In addition, the programmer of a {\em Model} needs only to understand
that the methods setTrainTwoSet and setTestTwoSet must be called in
order for the method metricsReport, which reports accuracy metrics of
a one-output {\em Model}, to be called.  These methods are also shown
in \beginred red \endcolor in the {\em DataSet} specification.

\section{DataSet}

{\em DataSet} is a standalone object that contains and manipulates the
dataset that {\em Model} will access.  All {\em DataSets} have raw,
training and test datasets, and know the number of inputs and
outputs. {\em DataSet} consists of:

\begin{itemize}

\item {\em Public} methods

\begin{itemize}

\item Raw dataset methods:
\begin{itemize}
\item void loadRaw( string\& $filename$ ): Takes $filename$ as
      argument (string), and loads a file into the $Raw$ dataset
      Matrix, returns nothing
\item bool rawLoaded(): Returns a flag to indicate if $Raw$ dataset is
      loaded
\item Matrix$<$ double $>$\& getRawMatrix(): returns the $Raw$ dataset
\item void setRawMatrix( Matrix$<$ double $>$\& $inMatrix$ ): loads
      Matrix $inMatrix$ into the $Raw$ dataset Matrix
\end{itemize}

\item Training set methods:
\begin{itemize}
\item void loadTrain( string\& $filename$ ): Takes $filename$ as
      argument (string), and loads a file into the $TrainSetData$
      Matrix, returns nothing
\item \beginred bool trainLoaded(): Returns a flag to indicate if
      $TrainSetData$ Matrix is loaded \endcolor
\item bool saveTrain( string\& $filename$ ): Takes $filename$ as
      argument (string), and saves the $TrainSetData$ Matrix into a
      file, returns flag indicating if save operation was successful
\item bool saveScales( string\& $filename$ ): Takes $filename$ as
      argument (string), and saves the scaling factors for the
      training set into the file, returns flag indicating if save
      operation was successful

{\bf Note:} see \hyperref[normEqn]{equation \ref{normEqn}} below for
an explanation of scaling.  Saved to the file will be $lbound$ (1
value), and for each $x$, $S$ and $min$.
\begin{equation}
x_{norm} = lbound + S(x - min)
\end{equation}
where $x_{norm}$ is the normalized result, $lbound$ is the lower limit
of normalized values, $S$ is the scaling factor, and $min$ is the
minimum value of each $x$ in the dataset.

\item \beginred Matrix$<$ double $>$\& getTrainMatrix(): returns the
      $TrainSetData$ Matrix \endcolor
\item void setTrainMatrix( Matrix$<$ double $>$\& $inMatrix$ ): loads
      Matrix $inMatrix$ into the $TrainSetData$ Matrix
\end{itemize}

\item {\em TwoSet} accessors for the training set:
\begin{itemize}
\item \beginred bool setTrainTwoSet(): Constructs an internal {\em
      TwoSet} object from the training set, populating the `real'
      column of the {\em TwoSet} object with the output column from
      the training set {\em Matrix}, returns true if successful
      \endcolor
\item TwoSet\& getTrainTwoSet(): Returns the constructed {\em TwoSet}
      object from the training set
\item bool TrainTwoSetLoaded(): Returns true if the internal training
      {\em TwoSet} object was successfully constructed
\item bool saveTrainTwoSet( string\& $filename$ ): Takes $filename$ as
      argument (string), and saves the {\em TwoSet} object from the
      training set into a file, returns flag indicating if save
      operation was successful
\end{itemize}

\item Test set methods:
\begin{itemize}
\item void loadTest( string\& $filename$ ): Takes $filename$ as
      argument (string), and loads a file into the $TestSetData$
      Matrix, returns nothing
\item \beginred bool testLoaded(): Returns a flag to indicate if
      $TestSetData$ Matrix is loaded \endcolor
\item bool saveTest( string\& $filename$ ): Takes $filename$ as
      argument (string), and saves the $TestSetData$ Matrix into a
      file, returns flag indicating if save operation was successful
\item \beginred Matrix$<$ double $>$\& getTestMatrix(): returns the
      $TestSetData$ Matrix \endcolor
\item void setTestMatrix( Matrix$<$ double $>$\& $inMatrix$ ): loads
      Matrix $inMatrix$ into the $TestSetData$ Matrix
\end{itemize}

\item {\em TwoSet} accessors for the test set:
\begin{itemize}
\item \beginred bool setTestTwoSet(): Constructs an internal {\em
      TwoSet} object from the test set, populating the `real' column
      of the {\em TwoSet} object with the output column from the test
      set {\em Matrix}, returns true if successful \endcolor
\item TwoSet\& getTestTwoSet(): Returns the constructed {\em TwoSet}
      object from the test set
\item bool TestTwoSetLoaded(): Returns true if the internal test {\em
      TwoSet} object was successfully constructed
\item bool saveTestTwoSet( string\& $filename$ ): Takes $filename$ as
      argument (string), and saves the {\em TwoSet} object from the
      test set into a file, returns flag indicating if save operation
      was successful
\end{itemize}

\item General methods:
\begin{itemize}
\item void setInput( unsigned $I$ ): Set the number of input nodes to
      $I$
\item \beginred unsigned getInput(): Returns number of input nodes
      \endcolor
\item void setOutput( unsigned $O$ ): Set the number of output nodes
      to $O$
\item \beginred unsigned getOutput(): Returns number of output nodes
      \endcolor
\item \beginred unsigned getNumTrain(): Returns the number of
      examples in $TrainSetData$, the training set \endcolor
\item \beginred unsigned getNumTest(): Returns the number of examples
      in $TestSetData$, the test set \endcolor
\item \beginred void removeInputs( vector$<$ unsigned $>$\& $v$ ):
      Given a vector representing input positions within a dataset,
      removes the corresponding columns in $Raw$, $TrainSetData$, and
      $TestSetData$, thus effectively removing those inputs from a
      dataset \endcolor
\end{itemize}

{\bf Note:} In the following methods, a variate $x$ is normalized
according to:
\begin{equation}\label{normEqn}
x_{norm} = lbound + [({{x - min} \over {max - min}}) \times (hbound -
lbound)]
\end{equation}
where $x_{norm}$ is the normalized value, $lbound$ is the lower limit
of the normalized result, $hbound$ is the upper limit of the
normalized result, $min$ is the minimum value of $x$ in the dataset,
and $max$ is the maximum value of $x$ in the dataset.

\item Methods related to normalizing dataset variates
\begin{itemize}
\item void setInUpper( double $x$ ): Sets the upper limit of
      normalized input variates to $x$
\item double getInUpper(): Returns the upper limit of normalized input
      variates
\item void setInLower( double $x$ ): Sets the lower limit of
      normalized input variates to $x$
\item double getInLower(): Returns the lower limit of normalized input
      variates
\item void setOutUpper( double $x$ ): Sets the upper limit of
      normalized output variates to $x$
\item double getOutUpper(): Returns the upper limit of normalized
      output variates
\item void setOutLower( double $x$ ): Sets the lower limit of
      normalized output variates to $x$
\item double getOutLower(): Returns the lower limit of normalized
      output variates
\end{itemize}

\item Methods related to identifying output as discrete (held to 0 or
      1)
\begin{itemize}
\item void setDiscrete( bool $flag$ ): Sets output to be discrete if
      $flag$ is true, continuous if false
\item \beginred bool getDiscrete(): Returns the state of discrete
      output \endcolor
\item void setThreshold( double $x$ ): Sets the threshold value for
      discrete output to be $x$, any output $>= x$ will be 1, and any
      output $< x$ will be 0
\item double getThreshold(): Returns the threshold value
\end{itemize}

\item Methods related to transforming datasets:
\begin{itemize}
\item bool raw2train(): Converts the raw dataset to a training set by
      normalizing its input variates, and if nondiscrete, its output
      variates, returning true if the operation was successful
\item bool randomize( unsigned $n$ ): Randomizes the raw dataset into
      training and test sets which are then normalized, preserving the
      frequency of outcomes in the training and test sets, takes the
      number of exemplars $n$ to be placed in the test set as an
      argument {\em ( number of outputs must be only 1, and it must be
      discrete )}, returns true if successful
\item bool randomize( unsigned $n$, unsigned $d$ ): Randomizes the raw
      dataset into training and test sets which are then normalized,
      preserving the frequency of outcomes in the training and test
      sets, takes the numerator $n$ and denominator $d$ of a fraction
      representing the portion of examplars to be placed in the test
      set as arguments, {\em ( number of outputs must be only 1, and
      it must be discrete )}, returns true if successful
\item bool randomizeD( double $x$ ): Randomizes the raw dataset into
      training and test sets which are then normalized, preserving the
      frequency of outcomes in the training and test sets, takes the
      fraction $x$ as a double value from 0 to 1 representing the
      portion of examplars to be placed in the test set as argument
      {\em ( number of outputs must be only 1, and it must be discrete
      )}, returns true if successful
\end{itemize}

\item Methods related to calculating accuracy
\begin{itemize}
\item \beginred void metricsReport( ostream\& $streamObj$ ): Outputs
      all internal {\em TwoSet} metrics for a 1-output {\em Model},
      takes ostream object $streamObj$ as argument \endcolor
\end{itemize}

\item Methods related to logging to the history file:
\begin{itemize}
\item void setHistory( bool $flag$ ): Sets history logging, which {\em
      appends} all operations to the history file
\item bool getHistory(): Returns the current state of history logging
\item void setHistoryFilename( string\& $filename$ ): Sets the name of
      the history file to string $filename$
\item bool addHistory( ostringstream\& $outputStream$ ): appends the
      ostringstream object $outputStream$ to the history file.
\end{itemize}

\end{itemize}

\item {\em Private} members

\begin{itemize}
\item Matrix$<$ double $>$ $Raw$: the raw dataset
\item Matrix$<$ double $>$ $TrainSetData$: the training set
\item Matrix$<$ double $>$ $TestSetData$: the test set
\item TwoSet $TrainTwoSet$: the {\em TwoSet} object for the training
      set
\item TwoSet $TestTwoSet$: the {\em TwoSet} object for the test set
\item vector$<$ double $>$ $minima$: a vector containing the minimum
      values of the columns in $Raw$
\item vector$<$ double $>$ $maxima$: a vector containing the maximum
      values of the columns in $Raw$
\item unsigned $nInput$: the number of input nodes
\item unsigned $nOutput$: the number of output nodes
\item double $threshold$: the threshold for discrete output
\item double $inUpperLimit$: the upper limit for a normalized input
      variate
\item double $inLowerLimit$: the lower limit for a normalized input
      variate
\item double $outUpperLimit$: the upper limit for a normalized output
      variate
\item double $outLowerLimit$: the lower limit for a normalized output
      variate
\item bool $discreteFlag$: a flag to indicate if the output is to be
      considered discrete (held to 0,1)
\item bool $rawLoadedFlag$: a flag to indicate if $Raw$ loaded
\item bool $trainLoadedFlag$: a flag to indicate if $TrainSetData$
      loaded
\item bool $testLoadedFlag$: a flag to indicate if $TestSetData$
      loaded
\item bool $trainTwoSetFlag$: a flag to indicate if the training set
      {\em TwoSet} object has been successfully constructed
\item bool $testTwoSetFlag$: a flag to indicate if the test set {\em
      TwoSet} object has been successfully constructed
\item bool $historyFlag$: a flag to indicate if operations will be
      logged to a history file, this file is {\em appended}
\item bool $minimaxFlag$: a flag to indicate if minima and maxima
      vectors have been set
\item string $historyFilename$: the name of history file
\end{itemize}

\item {\em Private} Utility functions

\begin{itemize}
\item void minimax( Matrix$<$ double $>$\& $data$ ): computes the
      $minima$ and $maxima$ vectors from {\em Matrix} $data$
\item bool checkDiscrete( Matrix$<$ double $>$\& $data$ ): returns
      true if the output columns of {\em Matrix} $data$ are discrete,
      false if otherwise
\item void normalize( Matrix$<$ double $>$\& $data$ ): normalizes
      input variates of {\em Matrix} $data$, and if nondiscrete, its
      output variates {\em (note: $minima$ and $maxima$ vectors must
      have been previously set)}
\item void clear(): clears {\em DataSet Matrix} members and flags
\end{itemize}

\end{itemize}

\section{TwoSet}

{\em TwoSet} is a specialized standalone object that is designed to
calculate outcome metrics such as sensitivity, specificity, predictive
value positive, predictive value negative and ROC area from a {\em
one} output {\em Model}. The {\em TwoSet} data {\em Matrix} contains 2
columns, one for the known outcomes and one for the {\em Model's}
outputs.  The {\em Model} output {\em must be} discrete.  The format
of this data {\em Matrix} is:

\begin{center}
\begin{tabular}{cc}
$real_0$ & $test_0$ \\
$real_1$ & $test_1$ \\
$real_2$ & $test_2$ \\
$\ldots$ & $\ldots$ \\
$real_{n-1}$ & $real_{n-1}$ \\
\end{tabular}
\end{center}

where $real$ represents a known outcome, and $test$ is the {\em
Model's} output, for $n$ examples in the data {\em Matrix}, e.g.:

\begin{verbatim}
0.000000 0.002851
0.000000 0.002978
0.000000 0.002857
1.000000 0.571710
\end{verbatim}

For calculating ROC (receiver-operator characteristic) curve area, 2
methods are available.  One is the traditional method using trapezoids
and numerical calculation to calculate the area under the curve.  By
default, the number of trapezoids is set to the number of data points,
but that may be changed by the {\em setTrapThresholds()} method.
However, as should be evident, this method {\em underestimates} the
ROC area.  Wickens\footnote{Wickens, Thomas D., Elementary Signal
Detection Theory, Oxford University Press, New York, NY 2002,
pp. 60--74}\index{Wickens, Thomas D.} provides a statistical method to
calculate ROC area.  Here, \hyperref[invZarea]
{$z = Z(A) = \Phi^{-1}(A)$} is linearly modeled for sensitivity and
1 - specificity, and this linear model is used to calculate ROC area.
If enough data is available to bin and generate errors in the abscissa
and ordinate of this line, then a p-value will be generated to reflect
the probability of a least squares fit to that line, where a smaller
p-value is {\em worse} (closer to 1 is more linear).

The method $ROCarea( ostream\& )$ may be used to automatically choose
which method of calculating ROC curve are is used, depending on how
many data points are available.

Two errors may be thrown from this class.  The first is
$DivisionByZero$ where division by zero is encountered, generally in
calculating positive and negative predictive value.  The second is
$twoSetErr$, which will contain errors generated in statistical
calculation or linear modeling. $DivisionByZero$ always contains the
message ``division by zero'', whereas $twoSetErr$'s $what()$ member
will contain a message reflecting the type of error.  An example of
the
latter would be:

\begin{verbatim}
string reply;
cout << "What is the name of the file? ";
cin >> reply;
TwoSet data;
data.load( reply );
try {
   data.ROCarea( cout );
} catch ( TwoSet::twoSetErr& e ) { cout << e.what() << endl; }
\end{verbatim}

The {\em TwoSet} object consists of:

\begin{itemize}

\item {\em Public} methods

\begin{itemize}

The constructor TwoSet( Matrix$<$ double $>$\& $M$ ) constructs and
loads {\em Matrix} $M$ into a {\em TwoSet} object

\item Methods for loading Matrices and files into a {\em TwoSet}
      object:
\begin{itemize}
\item void setMatrix( Matrix$<$ double $>$\& $M$ ): Loads {\em Matrix}
      $M$ into a {\em TwoSet} object, returns nothing
\item void load( string\& $filename$ ): Takes $filename$ as argument
      (string), and loads a file into the {\em TwoSet} data {\em
      Matrix}, returns nothing
\item bool loaded(): Returns a flag to indicate if the {\em TwoSet}
      data {\em Matrix} is loaded
\end{itemize}

\item {\em TwoSet} accessors:
\begin{itemize}
\item Matrix$<$ double $>$\& getData(): Returns the {\em TwoSet} data
      {\em Matrix}
\item bool save( string\& $filename$ ): Takes $filename$ as argument
      (string), and saves the {\em TwoSet} object into a file, returns
      flag indicating if save operation was successful
\item unsigned getNumElements(): Returns the number of elements in the
      {\em TwoSet} data {\em Matrix}
\item TwoSet\& insertReal( vector$<$ double $>$\& $v$ ): Insert vector
      $v$ into the real column of the {\em TwoSet} data {\em Matrix},
      the number of elements of $v$ {\em must equal} the number of
      rows in the {\em TwoSet} data {\em Matrix}
\item vector$<$ double $>$\& real( vector$<$ double $>$\& $v$ ):
      Populate vector $v$ with the real column of the {\em TwoSet}
      data {\em Matrix}, the number of elements of $v$ {\em must
      equal} the number of rows in the {\em TwoSet} data {\em Matrix}
\item vector$<$ double $>$ real(): Get a new vector from the real
      column of the {\em TwoSet} data {\em Matrix}
\item TwoSet\& insertTest( vector$<$ double $>$\& $v$ ): Insert vector
      $v$ into the test column of the {\em TwoSet} data {\em Matrix},
      the number of elements of $v$ {\em must equal} the number of
      rows in the {\em TwoSet} data {\em Matrix}
\item vector$<$ double $>$\& test( vector$<$ double $>$\& $v$ ):
      Populate vector $v$ with the test column of the {\em TwoSet}
      data {\em Matrix}, the number of elements of $v$ {\em must
      equal} the number of rows in the {\em TwoSet} data {\em Matrix}
\item vector$<$ double $>$ test(): Get a new vector from the test
      column of the {\em TwoSet} data {\em Matrix}
\item double real( unsigned $r$ ): An accessor to set or retrieve the
      element in row $r$ in the real column of the {\em TwoSet} data
      {\em Matrix}
\item double test( unsigned $r$ ): An accessor to set or retrieve the
      element in row $r$ in the test column of the {\em TwoSet} data
      {\em Matrix}
\end{itemize}

\item Accessors and methods for {\em TwoSet} calculated metrics
\begin{itemize}
\item void setThreshold( double $x$ ): Sets threshold value to $x$,
      any test element $>=x$ will be counted as $1$, whereas any test
      element $<x$ will be counted as $0$
\item double getClassAcc(): Returns classification accuracy
\item double getSens(): Returns sensitivity, throws exception
      TwoSet::DivisionByZero\& $e$, where $e$.what() is ``division by
      zero''
\item double getSpec(): Returns specificity, throws exception
      TwoSet::DivisionByZero\& $e$, where $e$.what() is ``division by
      zero''
\item double getPVP(): Returns predictive value positive, throws
      exception TwoSet::DivisionByZero\& $e$, where $e$.what() is
      ``division by zero''
\item double getPVN(): Returns predictive value negative, throws
      exception TwoSet::DivisionByZero\& $e$, where $e$.what() is
      ``division by zero''
\item void ClassTable( ostream\& $outputStream$ ): Outputs to
ostream\&
      the classification table
\end{itemize}

\item ROC area calculation by trapezoidal method
\begin{itemize}
\item void setTrapThresholds( unsigned $n$ ): Sets number of
      trapezoids to calculate ROC area to $n$
\item unsigned getTrapThresholds(): Returns number of trapezoids to
      calculate ROC area
\item double getTrapROCarea(): Returns ROC area calculated by
      trapezoidal method
\end{itemize}

\item ROC area calculation by statistical method
\begin{itemize}
\item void setNbins( unsigned n ): If data can be binned, sets the
      number of bins to $n$
\item unsigned getNbins(): Returns number of bins specified by user
\item void setBinSize( unsigned $n$ ): If data can be binned, sets
      number of elements to be contained in a bin to $n$
\item unsigned getBinSize(): Return bin size specified by user
\item void NbinThreshold( unsigned $n$ ): Sets the number of data
      points below which data will not be binned
\item unsigned getNBinThreshold(): Returns the number of data points
      below which data will not be binned
\item bool getBinned(): Returns $true$ if data was binned
\item void setNbinsSetsSize( bool $flag$ ): Set $flag$ if number of
      bins sets bin size
\item bool getNbinsSetsSize(): Returns $true$ if number of bins sets
      bin size
\item double getStatROCarea(): Returns statistical ROC area, throws
      exception TwoSet::twoSetErr\& $e$, where $e$.what() contains a
      statistical or linear regression error message
\item double getStatP(): Returns p-value for fitted line (smaller is
      worse)
\item double getStatChi2(): Returns chi-squared value for fitted line
\end{itemize}

\item Methods and accessors which output to ostream a ROC report which
      uses either or both the trapezoidal and/or statistical method
\begin{itemize}
\item void ROCarea( ostream\& $outputStream$ ): Outputs to ostream\&
      ROC area, and a report detailing method and results of
      calculation
\item void setROCReportFlag( $flag$ ): Set $flag$ true if both
      statistical and trapezoidal methods are to be reported, false to
      use either the trapezoidal method if number of data points $<$
      threshold ($calcThresh$) or statistical method if $>=$ threshold
\item bool getROCReportFlag(): Returns flag indicating if either
      (false) or both (true) of the statistical and trapezoidal
      methods
      are used and reported, according to the above description of
      setROCReportFlag($\ldots$)
\item void setROCthresh( unsigned $n$ ): Sets threshold for number of
      non-zero, non-one sensitivity and 1 - specificity data points
      under which statistics will not be used
\item unsigned getROCthresh(): Returns number of non-zero, non-one
      sensitivity and 1 - specificity data points under which
      statistics will not be used
\item void setROCSearchFlag( bool $flag$ ): Set $flag$ true if best
      $p$ and ROC areas are searched and reported within a specified
      range of number of bins
\item bool getROCSearchFlag(): Returns flag indicating if best $p$ and
      ROC areas are searched and reported within a specified range of
      bins
\item void setMinROCSearchBins( unsigned $n$ ): Sets minimum number of
      bins to search for best $p$ and ROC areas (must be $>2$)
\item unsigned getMinROCSearchBins(): Returns minimum number of bins
      to search for best $p$ and ROC areas
\item void setMaxROCSearchBins( const unsigned $n$ ): Sets maximum
      number of bins to search for best $p$ and ROC areas
\item unsigned getMaxROCSearchBins(): Returns maximum number of bins
      to search for best $p$ and ROC areas
\end{itemize}

\item Accessors for the Kolmogorov-Smirnov test\footnote{Adapted from
      W. H. Press et. al, Numerical Recipes in C, 2nd edition,
      Cambridge
      University Press, Cambridge, 1992, pp. 623--626}
      \index{Numerical Recipes in C}
      \index{Kolmogorov-Smirnov!Test}
\begin{itemize}
\item double getKSD(): Returns Kolmogorov-Smirnov statistic D
\item double getKSP(): Returns \hyperref[ks_p]{Kolmogorov-Smirnov $p$-
value}
\end{itemize}

\end{itemize}

\item {\em Private} members

\begin{itemize}
\item Matrix$<$ double $>$ A: the {\em TwoSet} data {\em Matrix}
\item unsigned n: the number of elements in the {\em TwoSet} data {\em
      Matrix}
\item unsigned tp: the number of true positives
\item unsigned tn: the number of true negatives
\item unsigned fp: the number of false positives
\item unsigned fn: the number of false negatives
\item unsigned nThresholds: the number of thresholds for trapezoidal
      method
\item unsigned nBins: the number of bins for statistical ROC
\item unsigned binSize: the bin size for statistical ROC
\item unsigned binThresh: the number of data points under which data
      will not be binned
\item unsigned calcThresh: the number of data points under which
      statistics will not be used
\item double threshold: the threshold value
\item double statP: p-value for fitted line for statistical ROC
      calculation
\item double statChi2: chi-squared value for fitted line for
      statistical ROC calculation
\item double KSD: Kolmogorov-Smirnov test D value
\item double KSP: Kolmogorov-Smirnov test $p$-value
\item bool loadedFlag: a flag to indicate if the {\em TwoSet} data
      {\em Matrix} is loaded
\item bool thresholdFlag: a flag to indicate if the threshold is set
\item bool nBinFlag: flag indicates if number of bins determines bin
      size
\item bool binnedFlag: flag indicates if data was binned for
      statistical ROC calculation
\item bool statROCcalcFlag: flag indicates if statistical ROC area
      calculated
\item bool reportFlag: flag directs how to report statistical and/or
      trapezoidal ROC
\item bool searchFlag: flag indicates if largest p, largest ROC is
searched
\item bool KScalcFlag: flag indicates if Kolmogorov-Smirnov test
calculated
\end{itemize}

\item {\em Private} Utility functions
\begin{itemize}
\item void initialize(): Sets common initial values for constructors
\item void calculate(): Given a threshold, calculates tp, tn, fp, and
      fn
\item bool checkDiscrete( Matrix$<$ double $>$\& $M$ ): Returns true
      if the first column (representing real outcomes) of {\em Matrix}
      $M$ is discrete, false if otherwise
\item void statReport( ostream\& $outputStream$, unsigned $goodData$,
      unsigned $nBins$, double $AUC$, double $chi2$, double $p$ ):
      Utility method for ROCarea, outputs statistical report given
      ostream\& $outputStream$, number of non-zero, non-one data
      points $goodData$, number of bins $nBins$, area under the curve
      $AUC$, chi-square value $chi2$, and p-value $p$
\item void KScalc(): Utility method to calculate Kolmogorov-Smirnov
      test\footnote{Adapted from W. H. Press et. al, Numerical Recipes
      in C, 2nd edition, Cambridge University Press, Cambridge, 1992,
      pp.
      623--626} \index{Numerical Recipes in C} \index{Kolmogorov-
      Smirnov!Test}
\end{itemize}

\end{itemize}

\section{What you really need to know about programming a Model}

\subsection{Object Type}\label{objType}

All concrete derived objects of {\em Model} inherit its protected
member, (string) $objType$, which is the name of the object.
Typically, the class constructor defines this member, for example:

\begin{verbatim}
class BareProp : public Network {
public:
   BareProp() { objType = "BareProp"; } // default constructor
\end{verbatim}

Abstract classes {\em do not} define $objType$. This member is
accessed by the {\em Model} method getType().

\subsection{Copy constructor and overloaded = operator}\label{copy_fx}

All derived objects of {\em Model} {\bf must} have a copy constructor
and overloaded $=$ operator.  Typically, both call a protected utility
function so that both copy the same members.  The protected utility
function typically begins with a call to the copy method of its
immediate base class so that it inherits all of the members of its
base class(es).  For example, suppose a hypothetical neural network
object named "LittleNet" is derived from {\em Network}, and contains
the following members:

\begin{verbatim}
unsigned nNodes; // number of nodes

Matrix< double > W; // weight Matrix

vector< double > I, // input vector
                 y; // output vector
\end{verbatim}

The following would be defined in the header file:

\begin{verbatim}
class LittleNet : public Network {
public:
   // Copy constructor
   LittleNet( const LittleNet& rhs );

   // Overloaded = operator
   LittleNet& operator= ( const LittleNet& rhs );

private: // or protected if the object is a base class
   // Copy utility
   void copy( const LittleNet& rhs );
\end{verbatim}

Which would then be coded in the $C++$ file as:

\begin{verbatim}
// Copy constructor
LittleNet::LittleNet( const LittleNet& rhs )
{
   LittleNet::copy( rhs ); // use the copy utility
}

// Overloaded = operator
LittleNet& LittleNet::operator= ( const LittleNet& rhs )
{
   if ( &rhs != this ) // check for self-assignment
      LittleNet::copy( rhs ); // use the copy utility

   return *this; // enables A = B = C
}

// Copy utility
void LittleNet::copy( const LittleNet& rhs )
{
   Network::copy( rhs ); // call immediate base object copy
   nNodes = rhs.nNodes;
   W = rhs.W;
   I = rhs.I;
   y = rhs.y;
}
\end{verbatim}

{\large \bf Remember:} {\bf if you add any member to any existing
class, you} {\bf \em must} {\bf insert a corresponding line of code to
copy it in the copy utility function of that class.}

\subsection{Implementation of pure virtual methods in {\em Model}}

Ultimately, all models map the inputs of each exemplar to the outputs.
Thus, most of the {\em Model} base class methods pertain directly to
this process.  If you are writing a Model, you need to implement the
following pure virtual methods:

\subsubsection{setDataSet}
The pure virtual method

\begin{verbatim}
virtual void setDataSet( DataSet& $dataObj$ )
\end{verbatim}

loads a {\em DataSet} object into the derived {\em Model}.  The
protected member inherited from the base class {\em Model} is named
$theData$, so the first line of

\begin{verbatim}
void YourModel::setDataSet( DataSet& dataObj )
\end{verbatim}

where ``YourModel'' is the name of your derived {\em Model} object,
should be

\begin{verbatim}
theData = dataObj; // set "theData" to incoming DataSet
\end{verbatim}

The remainder of YourModel::setDataSet would then contain code that
tailors your model to the incoming {\em DataSet} object.  See
SimpleProp::setDataSet and BareProp::setDataSet for examples.

\subsubsection{train}
The pure virtual method

\begin{verbatim}
virtual double train()
\end{verbatim}

should train a {\em Model}, and return the final error (double) for
that model.  In the case of neural networks, training would involve
finding the weight vectors and Matrices that minimize final error. See
SimpleProp::train and BareProp::train for examples.

\subsubsection{reportAccuracy}
The pure virtual method

\begin{verbatim}
virtual void reportAccuracy( ostream& outputStream )
\end{verbatim}

should report the accuracy of the trained model to an ostream
outputStream.  {\em If you are programming an object that extends}
Network, {\em e.g. a neural network, you} {\bf do not} {\em need to
implement this virtual method, as} Network {\em implements
reportAccuracy for you.}

\subsubsection{outputHeader}
The pure virtual method

\begin{verbatim}
virtual void outputHeader( ostream& outputStream )
\end{verbatim}

should output a header to an ostream outputStream describing the
architecture of the derived {\em Model}.  See SimpleProp::outputHeader
and BareProp::outputHeader for examples.

\subsection{{\em Model} protected members}

The following are the protected members in {\em Model} that your
object needs to set or get:

\subsubsection{theData}

The protected member

\begin{verbatim}
DataSet theData
\end{verbatim}

represents the primary dataset, in the form of a DataSet object.  Your
setDataSet method needs to set this member.  See
SimpleProp::setDataSet and BareProp::setDataSet for examples.

\subsubsection{Train}

The protected member

\begin{verbatim}
Matrix< double > Train
\end{verbatim}

represents a Matrix that contains the training set input columns only,
typically constructed by a call within your setDataSet method to the
{\em Model} protected utility function extractInputMatrices().  See
SimpleProp::setDataSet and BareProp::setDataSet for examples.

\subsubsection{Test}

The protected member

\begin{verbatim}
Matrix< double > Test
\end{verbatim}

represents a Matrix that contains the test set input columns only,
typically constructed by a call within your setDataSet method to the
{\em Model} protected utility function extractInputMatrices().  See
SimpleProp::setDataSet and BareProp::setDataSet for examples.

\subsubsection{TrainOutput}

The protected member

\begin{verbatim}
Matrix< double > TrainOutput
\end{verbatim}

represents a Matrix that contains the training set output columns
only, typically constructed by a call within your setDataSet method to
the {\em Model} protected utility function extractOutputMatrices().
See BackProp::setDataSet for an example.\index{Backpropagation}

\subsubsection{TestOutput}

The protected member

\begin{verbatim}
Matrix< double > TestOutput
\end{verbatim}

represents a Matrix that contains the test set output columns only,
typically constructed by a call within your setDataSet method to the
{\em Model} protected utility function extractOutputMatrices().  See
BackProp::setDataSet for an example.\index{Backpropagation}

\subsubsection{builtFlag}

Set the protected member
\begin{verbatim}
bool builtFlag
\end{verbatim}

to {\em true} when the Model has been formally constructed, and its
architecture set.  See SimpleProp::setHidden and BareProp::setHidden
for examples.

\subsubsection{historyFlag}

Use the protected member

\begin{verbatim}
bool historyFlag
\end{verbatim}

to decide if operations will be logged to a history file using the
protected {\em Model} utility function addHistory.  See SimpleProp and
BareProp code for examples, as most methods use addHistory.

\subsubsection{errorType}

Use the protected member

\begin{verbatim}
bool errorType
\end{verbatim}

to determine if the type of error will be least mean squared
($errorType = 0$) or cross-entropy ($errorType = 1$).
\index{Error function!Cross-entropy}
\index{Error function!Least squares}
Typically, this member is used as an argument in {\em errorFunction}
constructors (see \hyperref[errFunc]{section \ref{errFunc}}, and
SimpleProp::trainSet() and BareProp::trainSet() as examples.)

\subsection{{\em Model} protected utility functions}

The following are the protected utility functions in {\em Model} that
your object will want to call:

\subsubsection{extractInputMatrices}

Call the utility function

\begin{verbatim}
void extractInputMatrices()
\end{verbatim}

in your setDataSet method to extract the input columns of the training
set into the protected {\em Model} Matrix $Train$, and the input
columns of the test set into the protected {\em Model} Matrix $Test$.

\subsubsection{extractOutputMatrices}

Call the utility function

\begin{verbatim}
void extractOutputMatrices()
\end{verbatim}

in your setDataSet method to extract the output columns of the
training set into the protected {\em Model} Matrix $TrainOutput$, and
the output columns of the test set into the protected {\em Model}
Matrix $TestOutput$.

\subsubsection{addHistory}

Pass an ostringstream object to the utility function

\begin{verbatim}
bool addHistory( ostringstream $outputStream$ )
\end{verbatim}

for that object to be added to the history file.  This method should
be called for every operation.  See SimpleProp and BareProp code for
examples, as most methods use addHistory.

\subsubsection{copy}

See section \ref{copy_fx} for details.

\section{Model}

In the following specification, methods and members that pertain
directly to coding a concrete derived {\em Model} implementation are
shown in \beginred red \endcolor.

{\bf Model} is the base object. {\em Model} provides access to the
dataset, and contains the virtual methods common to all modeling
techniques.  {\em Model} consists of:

\begin{itemize}

\item {\em Public} methods

\begin{itemize}

\item Copy constructor and overloaded $=$ operator, see section
\ref{copy_fx} for discussion:
\begin{itemize}
\item Model\& operator= ( const Model\& rhs )
\item Model( const Model\& rhs )
\end{itemize}

\item Methods related to accessing the dataset to model:
\begin{itemize}
\item \beginred virtual void setDataSet( DataSet\& $data$ ): A pure
      virtual method that loads the {\em DataSet} object $data$ into
      the {\em Model} \endcolor
\item DataSet\& getDataSet(): Returns a the {\em Model's} dataset in
      the form of a {\em DataSet} object
\end{itemize}

\item Methods related to training the model:
\begin{itemize}
\item \beginred virtual double train(): A pure virtual method that
      trains a {\em Model}, and returns the final error (double) for
      that model \endcolor
\item \beginred virtual void reportAccuracy( ostream\& $outputStream$
      ): A pure virtual method that reports the accuracy of the
      trained model to an ostream $outputStream$ \endcolor
\end{itemize}

\item Methods related to logging to files or reporting to user:
\begin{itemize}
\item \beginred virtual void outputHeader( ostream\& $outputStream$ ):
      A pure virtual method which outputs a header to an ostream
      $outputStream$ describing the architecture of the model, {\em
      Model} type dependant \endcolor
\item void setHistory( bool $flag$ ): Sets history logging, which {\em
      appends} all operations to the history file
\item bool getHistory(): Returns the current state of history logging
\item void setHistoryFilename( string\& $filename$ ): Sets the name of
      the history file to string $filename$
\item string getHistoryFilename(): Returns the name of the history
      file
\item void setLastop( bool $flag$ ): Sets last operation logging,
      which {\em overwrites} a file to report the building (training
      in the case of {\em Iterative}) of a specific model
\item bool getLastop(): Returns the current state of last operation
      logging
\item void setLastopFilename( string\& $filename$ ): Sets the name of
      the last operation logging file to $filename$
\item string getLastopFilename(): Returns the name of the last
      operation logging file
\end{itemize}

\item Methods related to setting the output error function
\begin{itemize}
\item void setLMSerror(): Sets the output error function to least mean
      squared, which is \index{Error function!Least squares}
\begin{equation}
E = \frac{1}{2} \sum_{i=1}^{m} (y_i - o_i)^2
\end{equation}
where $E$ is the error, $m$ the number of outputs, $y$ is the known
value and $o$ the Model's output.
\item void setXEerror(): Sets the output error function to be
      \hyperref[xentropyfunk]{cross-entropy},
      \index{Error function!Cross-entropy} which is
\begin{equation}
E = \sum_{i=1}^{m} -y\log(o)-(1-y)\log(1-o)
\end{equation}
{\em Note: for the cross-entropy error function, the output } {\bf
must} {\em be } {\bf discrete}.
\end{itemize}

\end{itemize}

\item {\em Protected} members

\begin{itemize}
\item \beginred DataSet $theData$: the primary dataset, in the form of
      a {\em DataSet} object \endcolor
\item \beginred Matrix$<$ double $>$ $Train$: a Matrix that contains
      the training set input columns {\em only}, constructed by a call
      to the protected utility function {\em extractInputMatrices()}
      \endcolor
\item \beginred Matrix$<$ double $>$ $Test$: a Matrix that contains
      the test set input columns {\em only}, constructed by a call to
      the protected utility function {\em extractInputMatrices()}
      \endcolor
\item \beginred Matrix$<$ double $>$ $TrainOutput$: a Matrix that
      contains the training set output columns {\em only}, constructed
      by a call to the protected utility function {\em
      extractOutputMatrices()} \endcolor
\item \beginred Matrix$<$ double $>$ $TestOutput$: a Matrix that
      contains the test set output columns {\em only}, constructed by
      a call to the protected utility function {\em
      extractOutputMatrices()} \endcolor
\item \beginred bool $builtFlag$: a flag to indicate if the {\em
      Model} has been formally constructed \endcolor
\item \beginred bool $historyFlag$: a flag to indicate if operations
      will be logged to a history file, this file is {\em appended}
      \endcolor
\item bool $lastopFlag$: a flag to indicate if model building
      (training in the case of {\em Iterative}) will be logged to a
      file describing the last operation, this file is {\em
      overwritten}
\item \beginred bool $errorType$: flag indicates type of error
function,
       $0$ if least mean squared, $1$ if cross-entropy \endcolor
       \index{Error function!Cross-entropy}
       \index{Error function!Least squares}
\item \beginred bool $builtFlag$: a flag to indicate the type of error
      function, 0 if least mean squared, 1 if cross-entropy \endcolor
      \index{Error function!Cross-entropy}
      \index{Error function!Least squares}
\item \beginred string $objType$: the name of the object, see
      \hyperref[objType]{section \ref{objType}} \endcolor
\item string $historyFilename$: the name of history file
\item string $lastopFilename$: the name of the file to log the last
      operation
\item string $errorLabel$: a string containing the name of the output
      error function
\end{itemize}

\item {\em Protected} Utility functions

\begin{itemize}
\item \beginred void extractInputMatrices(): Utility function which is
      intended to be called by the derived {\em Model} in its {\em
      setDataSet} method, and places the input columns of the training
      set in Matrix $Train$, and the input columns of the test set in
      Matrix $Test$ \endcolor
\item \beginred void extractOutputMatrices(): Utility function which
      is intended to be called by the derived {\em Model} in its {\em
      setDataSet} method, and places the output columns of the
      training set in Matrix $TrainOutput$, and the output columns of
      the test set in Matrix $TestOutput$ \endcolor
\item \beginred bool addHistory( ostringstream\& $outputStream$ ):
      Utility function to append the ostringstream object
      $outputStream$
      to the history file. \endcolor
\item \beginred void copy( const Model\& rhs ): Copy utility function,
      see section \ref{copy_fx} for discussion. \endcolor
\end{itemize}

\end{itemize}

\section{What you really need to know about programming an
Iterative Model}

{\em Iterative} inherits methods setDataSet, train, reportAccuracy and
outputHeader from {\em Model}.  If you are programming an {\em
Iterative Model}, there are 2 pure virtual methods in {\em Iterative}
that you must also implement:

\subsubsection{trainSet}

{\em Iterative} handles models that require multiple passes
(iterations) to train.  Objects derived from {\em Iterative} must thus
provide a method to adjust a model by a single pass through the
training set, and {\em Iterative} will call this method iteratively to
reach the final trained state.  Thus the method:

\begin{verbatim}
virtual double trainSet()
\end{verbatim}

changes the components of a {\em Model} (e.g. weight Matrices and
vectors in a neural network), and returns the set error.  See
SimpleProp::trainSet() and BareProp::trainSet() for examples.

\subsubsection{classAccuracy}

With each iteration, {\em Iterative} reports the set error and
classification accuracy for the training and/or test sets if the
outputs are discrete.  Thus, the pure virtual method:

\begin{verbatim}
virtual void classAccuracy( ostream& outputStream )
\end{verbatim}

must output to an ostream object the classification accuracy for the
training and/or test sets.  {\em If you are programming a} Network
{\em object, you} {\bf do not} {\em need to implement this virtual
method, as} Network {\em implements classAccuracy for you.}

\subsubsection{getGradMax}

Most iterative algorithms employ some form of gradient descent, and
the iterative class can report the maximum absolute gradient and use
it to determine cessation of training.  The pure virtual method:

\begin{verbatim}
virtual double getGradMax()
\end{verbatim}

must be coded to return the maximum absolute gradient of the derived
object. {\em If you are programming a} Network {\em object, you}
{\bf do not} {\em need to implement this virtual method, as} Network
{\em implements getGradMax for you.}  However, you must code
\hyperref[packutil]{Network::pack()} to create the $Network$ protected
member $stackG$ so that Network::getGradMax() returns the proper
value.

\subsubsection{runHeader}\label{runHeaderSection}

If you are adding parameters specific to an {\em Iterative} derived
object that should be output to screen or log files prior to an {\em
Iterative} run, you will need to implement a runHeader method in the
derived class:

\begin{verbatim}
virtual void runHeader( ostream& outputStream)
\end{verbatim}

which outputs to an ostream object the specific parameters to be
displayed.  {\em If you are programming a} Network {\em object, you
will need to add code to this method in} Network.

\hspace{10pt} \linebreak

As with any derivation of {\em Model}, a copy utility function, copy
constructor and overloaded $=$ operator must be implemented.  See
section \ref{copy_fx} for details.

\subsubsection{iteration}

You may access the iteration index in your derived object by the
protected member unsigned $iteration$.

\subsubsection{boundsErrorFlag}

If your object can generate a numerical bounds error (see
\hyperref[errFunc]{section \ref{errFunc}} for an example of a method
which generates a numerical bounds error), then use the protected
member bool $boundsErrorFlag$ to indicate if such an error has been
thrown.

\subsubsection{gradMaxFlag}

Use the protected member bool $gradMaxFlag$ to determine if a gradient
is to be constructed in your derived object.

\section{Iterative}

{\bf Iterative} extends {\em Model}.  Neural computational algorithms
are iterative, and thus use {\em Iterative} as a base.  {\em
Iterative} manages methods and parameters related to iterating through
a dataset.

In the following specification, methods and members that pertain
directly to coding a concrete derived {\em Iterative} implementation
are shown in \beginred red \endcolor.

\begin{itemize}

\item {\em Public} methods

\begin{itemize}

\item Copy constructor and overloaded $=$ operator, see section
\ref{copy_fx} for discussion:
\begin{itemize}
\item Iterative\& operator= ( const Iterative\& rhs )
\item Iterative( const Iterative\& rhs )
\end{itemize}

\item Methods pertaining to stopping training
\begin{itemize}
\item void setMaxIterations( unsigned $maxIterations$ ): Sets the
      maximum number of iterations through the training set to
      $maxIterations$
\item unsigned getMaxIterations(): Returns the maximum number of
      iterations through the training set
\item void setMinStop( bool $minStopFlag$ ): Flag sets if a minimum
      error threshold will be used to stop training
\item bool getMinStop(): Returns flag indicating if a minimum error
      threshold will be used to stop training
\item void setMinError( double $minError$ ): Sets minimum error
      threshold value to $minError$
\item double getMinError(): Returns minimum error threshold value
\item void setChangeStop( bool $changeStopFlag$ ): Flag sets if
      training stops after error increases over 1 iteration
\item bool getChangeStop(): Returns flag indicating if change over 1
      iteration will be used to stop training
\item void setChange( double $\Delta$ ): Sets limit of change value to
      $\Delta$
\item double getChange(): Returns limit of change value
\item void setWindowStop( bool $windowStopFlag$ ): Flag sets if
      training stops if set error increases over a specified window
      width
\item bool getWindowStop(): Returns flag indicating if set error
      increases over a window width will be used to stop training
\item void setWindow( unsigned $window$ ): Sets the value of the
      window width to $window$
\item unsigned getWindow(): Returns specified window width value

\item void setGradStop( bool $gradMaxFlag$ ): Flag sets if training
      will stop when the maximum absolute gradient decreases below the
      limit
\item bool getGradStop(): Returns flag indicating if training will
      stop when the maximum absolute gradient decreases below the
      limit
\item void setGradMaxLimit( double $limit$ ): Sets the value of the
      maximum absolute gradient limit to $limit$
\item double getGradMaxLimit(): Returns the value of the maximum
      absolute gradient limit

\end{itemize}

\item Methods pertaining to printing while training
\begin{itemize}
\item void setLogPrint( bool $logPrintFlag$ ): Sets if prints out by
      powers of 10 ( 1-10 by ones, 10-100 by 10s, 100-1000 by 100s,
      etc. )
\item bool getLogPrint(): Returns flag indicating if printing by
      powers of 10
\item void setPrintCount( unsigned $printCount$ ): Sets the linear
      print counter to $printCount$
\item unsigned getPrintCount(): Returns the linear print counter
\end{itemize}

\item General methods
\begin{itemize}
\item \beginred virtual void setDataSet( DataSet\& $data$ ): A pure
      virtual method inherited from {\em Model} that loads the {\em
      DataSet} object $data$ into the {\em Model} \endcolor
\item \beginred virtual double train(): A pure virtual method
      inherited from {\em Model} that trains a {\em Model}, and
      returns the final error (double) for that model \endcolor
\item \beginred virtual void reportAccuracy( ostream\& $outputStream$
      ): A pure virtual method inherited from {\em Model} that reports
      the accuracy of the trained model to an ostream $outputStream$
      \endcolor
\item \beginred virtual void outputHeader( ostream\& $outputStream$ ):
      A pure virtual method inherited from {\em Model} which outputs a
      header to an ostream $outputStream$ describing the architecture
      of the model, {\em Model} type dependant \endcolor
\item \beginred virtual double trainSet(): A pure virtual method which
      trains one iteration through the training set, {\em Iterative
      Model} type dependant, and returns the set error \endcolor
\item \beginred virtual void classAccuracy( ostream\& $outputStream$
      ): A pure virtual method which outputs classification accuracies
      to an ostream $outputStream$ for an {\em Iterative} output table
      \endcolor
\end{itemize}

\end{itemize}

\item {\em Protected} members

\begin{itemize}
\item \beginred unsigned $iteration$: the iteration index \endcolor
\item unsigned $maxIterations$: the maximum number of iterations
      through the training set
\item bool $logPrintFlag$: flag indicates if printing by powers of 10
\item unsigned $printCount$: the value for the linear print counter
\item bool $minStopFlag$: flag indicates if a minimum error threshold
      will be used to stop training
\item double $minError$: minimum error threshold value
\item bool $changeStopFlag$: flag indicates if change in error over 1
      iteration will be used to stop training
\item double $delta$: value of change over 1 iteration
\item bool $gradMaxFlag$: flag indicates if training will stop when
      maximum absolute gradient decreases below $gradMaxLimit$
\item double $gradMaxLimit$: value of the maximum absolute gradient
      limit
\item bool $windowStopFlag$: flag indicates if set error increases
      over a window width will be used to stop training
\item unsigned $window$: value of the window width
\item \beginred bool $boundsErrorFlag$: flag indicates if out of
      numerical bounds error encountered (e.g. $\log(0)$) \endcolor
\end{itemize}

\item {\em Protected} Utility functions

\begin{itemize}
\item \beginred void copy( const Iterative\& rhs ): Copy utility
      function, see section \ref{copy_fx} for discussion. \endcolor
\item virtual void runHeader( ostream\& $outputStream$ ): A pure
      virtual utility function that outputs parameters specific
      to derived objects preceeding an Iterative run
\item \beginred virtual double getGradMax(): A pure virtual utility
      function that returns the maximum absolute gradient value
      \endcolor
\end{itemize}

\end{itemize}

\section{What you really need to know about programming a Network
Model}

\subsection{Implementation of pure virtual methods in {\em Network}}
\label{NetworkMethodsSect}

{\em Network} inherits methods setDataSet, train, reportAccuracy and
outputHeader from {\em Model}, and trainSet from {\em Iterative}, so
you must code all of these methods.  (See the respective methods in
SimpleProp and BareProp as examples.)  If you are programming a {\em
Network Model}, there are 7 additional pure virtual methods in {\em
Network} that you must implement in the concrete derived object:

\subsubsection{df}

The method

\begin{verbatim}
virtual unsigned df()
\end{verbatim}

must be coded to return the number of degrees of freedom of your {\em
Network} object, which is equal to the number of nodal connections.
See SimpleProp::df and BareProp::df as examples.

\subsubsection{randomize}

The method

\begin{verbatim}
virtual void randomize()
\end{verbatim}

must be coded to randomize the weights in your {\em Network} object.
See SimpleProp::randomize and BareProp::randomize as examples.

\subsubsection{load and save}

The methods

\begin{verbatim}
virtual bool load( string& filename )
virtual bool save( string& filename )
\end{verbatim}

must be coded to load and save your {\em Network} derived object from
and to a file named ``filename''.  Typically the method save would use
your coded method outputHeader to write the {\em Network} object
architecture to a file, then the Matrices and vectors representing the
{\em Network} model.  As expected, load should do the opposite.  See
SimpleProp::load, SimpleProp::save, BareProp::load and BareProp::save
as examples.

\subsubsection{innerTrainSet}\index{Automatic stepsize
selection}\label{autoStep}

The method

\begin{verbatim}
virtual double innerTrainSet()
\end{verbatim}

must be coded in a {\em Network} object to calculate the gradient and
error.  The reason that another method in addition to the pure virtual
method trainSet inherited from {\em Iterative} must be coded is to
support automatic stepsize selection.  The strategy proceeds as
follows: trainSet is intended to call innerTrainSet to determine if
the set error would increase, and set the learning rate (``step
size'') accordingly lower if it would.  Thus, innerTrainSet calculates
the gradient and error, but does not necessarily update the weights.
If automatic stepsize {\em is not} selected, trainSet simply calls
innerTrainSet to calculate the gradient and error, and updates the
weights.  If automatic stepsize {\em is} selected, trainSet calls
innerTrainSet to calculate the gradient and error, and reduces the
stepsize before updating the weights if the error is found to
increase. See SimpleProp::innerTrainSet, BareProp::innerTrainSet and
BackProp::innerTrainSet as examples of the way trainSet and
innerTrainSet are coded to work together.

Richard Golden has suggested the following automatic stepsize
selection algorithm, which he terms the ``Crummy Back-tracking Line
Search Algorithm'':

\begin{itemize}
\item Let DELTAERROR be a small non-negative number (e.g., DELTAERROR
= 1e-10 or 0)
\item Let MAXLOOPS = a small positive integer (e.g., MAXLOOPS = 2 or 3
or 4)
\item Let GAMMA be some positive number less than 1 (e.g., GAMMA = 0.5
or 0.8)
\item Let LoopCounter = 0
\item Set ``candidate'' learning rate equal to 1.0
\item While Exit Loop is FALSE
\begin{itemize}
\item LoopCounter = LoopCounter + 1
\item Compute change to weights using current candidate learning rate
\item Evaluate the new ERROR at the weights you obtained to get a
``candidate error change''
\item If (``candidate error change'' $<$ DELTAERROR) OR (LoopCounter
$>$ MAXLOOPS), then EXITLOOP is set to TRUE
\item Else set current candidate learning rate = GAMMA [candidate
learning rate]
\item End While Loop
\end{itemize}
\end{itemize}

Don't be fooled by the name--this algorithm is far from ``crummy''.

{\bf Note:} This algorithm is only to be used with off-line or
batch/epoch learning.

\subsubsection{forward}\index{Forward propagation}

The method

\begin{verbatim}
virtual void forward( Matrix< double >& M, const unsigned r )
\end{verbatim}

must be coded to take a Matrix containing only the inputs of a
training or test set as one argument (typically the $Train$ or $Test$
inherited protected members from {\em Model}), the position of the
exemplar in the training or test set as a second argument, and derive
the {\em Network's} output.  {\em This is not returned by the method.
This is very important: in single output} Networks, {\em the
inherited} Network {\em protected member o would be set by this
method, and in multiple output} Networks, {\em the inherited} Network
{\em protected member vector$<$ double $>$ output would be set by this
method.  Also very important, if sigmoidal output nodes are supported,
the inherited} Network {\em protected member x would be set by this
method, and in multiple output} Networks, {\em the inherited} Network
{\em protected member vector$<$ double $>$ xs would be set by this
method.}\index{Transfer function!Sigmoidal}

\subsubsection{removeInputs}\label{removeInputs}

The method

\begin{verbatim}
virtual void removeInputs( vector< unsigned >& v )
\end{verbatim}

must be coded to take as its argument a vector of unsigned
representing positions of input nodes, and to remove those nodes from
the {\em Network} derived object.  Typically, this would be
accomplished first by a call to
{\em DataSet}::removeInputs( vector$<$ unsigned $>$\& $v$ ),
followed by a call to {\em Model}::extractInputMatrices(), followed by
code which removes input nodes in the structures representing weights
in the {\em Network} object, typically {\em Matrix} columns.  See
SimpleProp::removeInputs and BareProp::removeInputs as examples.

\subsection{Code that must be added to {\em RegressNet}}

If you are creating a regressable {\em Network} object (and
chances are that you are,) you {\em must} add code to the methods
RegressNet::copy\_network() and RegressNet::network\_name().  See
section \ref{regresscode} for details.

\subsection{{\em Network} protected members}

The following are the protected members in {\em Network} that your
object needs to set or get:

\subsubsection{trainingType}

The protected member

\begin{verbatim}
unsigned trainingType
\end{verbatim}

determines the type of gradient descent algorithm of your neural
network. $trainingType = 0$ is defined as canonical backpropagation.
If you add other types of training to the {\em Network} class, you
must define a unique $trainingType$ for them.  Currently coded
algorithms include $trainingType =$

\begin{enumerate}
\item Conjugate gradient descent\index{Conjugate gradient descent}
\item Shanno algorithm\index{Shanno algorithm}
\end{enumerate}

{\bf NOTE:} Look closely at the code for Network::getGradMax(). You
will find that this code assumes that if $trainingType$ is not 0, then
$stackG$ will be constructed.  If this is not true for your
$trainingType$, then you must modify the if block in
Network::getGradMax().

\subsubsection{o}

The protected member

\begin{verbatim}
double o
\end{verbatim}

is the output for a single output network.  {\em If your} Network {\em
object supports single outputs, it must set this protected member,}
typically in the method forward.\index{Forward propagation}
Since the method trainSet will typically call the method forward,
trainSet will generally access the member $o$.  See
SimpleProp::forward and BareProp::forward for examples in setting the
protected member $o$, and SimpleProp::trainSet and BareProp::trainset
for examples in accessing the protected member $o$.

\subsubsection{x}

The protected member

\begin{verbatim}
double x
\end{verbatim}

contains the $x$ term for the sigmoidal output $o$ in a single output
network:\index{Transfer function!Sigmoidal}

\begin{displaymath}
o = 1/(1+e^{-x})
\end{displaymath}

{\em If your} Network {\em object supports single sigmoidal outputs,
it must set this protected member,} typically in the method forward.
\index{Forward propagation} See SimpleProp::forward and
BareProp::forward for examples in setting the protected member $x$.

\subsubsection{output}

The protected member

\begin{verbatim}
vector< double > output
\end{verbatim}

is the output for a multiple output network.  {\em If your} Network
{\em object supports multiple outputs, it must set this protected
member,} typically in the method forward.\index{Forward propagation}
Since the method trainSet will typically call the method forward,
trainSet will generally access the member $output$.

\subsubsection{xs}

The protected member

\begin{verbatim}
vector< double > xs
\end{verbatim}

contains the $\vec{\bf x}$ terms for the sigmoidal outputs
$\vec{\bf o}$ in a multiple output network:
\index{Transfer function!Sigmoidal}

\begin{displaymath}
\vec{\bf o} = 1/(1+e^{-\vec{\bf x}})
\end{displaymath}

{\em If your} Network {\em object supports multiple sigmoidal outputs,
it must set this protected member,} typically in the method forward.
\index{Forward propagation} See BackProp::forward for an example in
setting the protected member $xs$.\index{Backpropagation}

\subsubsection{randomLimit}

The protected member

\begin{verbatim}
double randomLimit
\end{verbatim}

is used to specify the range of randomized weights, which will be
between $- randomLimit$ and $+ randomLimit$.  Access this protected
member in your randomize method.  See SimpleProp::randomize and
BareProp::randomize for examples.

\subsubsection{weightsSetFlag}

Your derived {\em Network} object must set the protected member

\begin{verbatim}
bool weightsSetFlag
\end{verbatim}

to $true$ when the weights are set, either by randomization or loading
weights from a file.  See SimpleProp::randomize, BareProp::randomize,
SimpleProp::load and BareProp::load for examples of setting
$weightsSetFlag$ to $true$, and SimpleProp::setHidden and
BareProp::setHidden for examples of setting $weightsSetFlag$ to
$false$.

\subsubsection{biasFlag}

Use the protected member

\begin{verbatim}
bool biasFlag
\end{verbatim}

to determine if your {\em Network} object has biases.

\subsubsection{eta}

Use the protected member

\begin{verbatim}
double eta
\end{verbatim}

for the learning rate $\eta$ of your {\em Network} object.

\subsubsection{batchEpochFlag}\index{Backpropagation!Off-line}

Use the protected member

\begin{verbatim}
bool batchEpochFlag
\end{verbatim}

to determine if your {\em Network} object uses batch/epoch learning or
not.  (See the description of classic off-line (``batch'') backprop
in \hyperref[learningAlgs]{section \ref{learningAlgs}}.)
\index{Backpropagation!Off-line}{\em Pay special attention to how the
learning rate is set.}  It is expected that the driver will guide the
user to the appropriate setting of the learning rate.  Note also that
``off- line'' learning is synonymous with `` batch'' or ``epoch''
learning, whereas ``on-line'' learning refers to updating the weights
with each exemplar.\index{Backpropagation!On-line}  Here, if
batchEpochFlag is {\em true}, ``offline'' learning is invoked, and if
batchEpochFlag is {\em false}, ``online'' learning is used.

\label{discussDecay}
\subsubsection{weightDecayFlag}\label{decayFlagSection}\index{Weight
decay}

Use the protected member

\begin{verbatim}
bool weightDecayFlag
\end{verbatim}

to determine if your {\em Network} object uses weight decay.  (See the
discussion of \hyperref[deferrorfunk]{equation \ref{deferrorfunk} in
section \ref{objFunc}} for an explanation of weight decay.)  The
learning rate equivalent of equation \ref{deferrorfunk} is

% Note the necessity of \left. even though it doesn't show
\begin{eqnarray*}\index{Weight decay}
\vec w_{t+1} & = & \vec w_t - \eta \left. \D{E'}w \right|_{w_t} \\
& = & \vec w_t - \eta \left. \D{E}w \right|_{w_t} - 2\eta\lambda\vec
w_t \\
& = & (1-2\eta\lambda)\vec w_t - \eta \left. \D{E}w \right|_{w_t}
\end{eqnarray*}

where $(1-2\eta\lambda)$ is ``weight decay''.\index{Weight decay}
$\lambda$ may be set by either of 2 methods ($N$ is the sample number
size, see \hyperref[objFunc]{equation \ref{deferrorfunk} in section
\ref{objFunc}} for a discussion of $\sigma_w$):

\begin{itemize}

\item The complexity of the population objective function $>>$ sample
objective function:

\begin{equation}
\lambda = \frac 1{N2\sigma_w^2}
\end{equation}

\item The complexity of the population objective function $\simeq$
sample objective function:

\begin{equation}
\lambda = \frac 1{2\sigma_w^2}
\end{equation}

\end{itemize}

It is anticipated that the driver will guide the user to set the
appropriate value of $\lambda$, and {\em insure that the user knows
what he or she's doing if they change it!}

\subsubsection{decay}\index{Weight decay}

You can use the protected member

\begin{verbatim}
double decay
\end{verbatim}

for the weight decay term $2\lambda$ of your {\em Network} object.
Note the multiple of 2 so that calculation of decay with the gradient
as in the right hand term of equation~\ref{gradcalc},
$(1/\sigma_w^2)[{\bf W} \;,\; {\bf b}]$, does not require the
additional step of multiplication by 2.

\subsubsection{regularizer}\index{Weight decay}

You can use the protected member

\begin{verbatim}
double regularizer
\end{verbatim}

if you need to access $\lambda$ directly, e.g. for of calculation of
your error function in equation \ref{deferrorfunk}, which like
$decayTerm$ only needs to be calculated once prior to iterative calls
to the TrainSet method, and is thus placed in the member function
runHeader, where the {\em Iterative} class will call it once prior to
iteratively calling the TrainSet method.

\subsubsection{decayTerm}\label{decayTermSect}\index{Weight decay}

In the protected member function runHeader in the {\em Network} class,

\begin{verbatim}
double decayTerm
\end{verbatim}

is defined to be $1-2\eta\lambda$ (or $1-\eta decay$), which only
needs to be calculated once prior to iterative calls to the TrainSet
method, and is thus placed in the member function runHeader, where the
{\em Iterative} class will call it once prior to iteratively calling
the TrainSet method.  You may therefore use $decayTerm$ as the weight
multiple for weight decay in your {\em Network} object if the decay is
not calculated with the gradient.  If the decay is calculated with the
gradient as in equation~\ref{gradcalc}, then $decay$ is used.  See
BareProp::TrainSet and SimpleProp::TrainSet for examples using
$decayTerm$.

\subsubsection{stackG}

The protected member

\begin{verbatim}
vector< double > stackG
\end{verbatim}

is a single vector which contains the entire weight gradient of your
{\em Network} derived object.  It is used by the protected utility
functions pack and unpack (\hyperref[packutil]{see below}).

\subsection{{\em Network} protected utility functions}

\subsubsection{Copy utility function, constructor and overloaded =
operator}

As with any derivation of {\em Model}, a copy utility function, copy
constructor and overloaded $=$ operator must be implemented.  See
section \ref{copy_fx} for details.

\subsubsection{runHeader utility function}

If you are adding parameters specific to a {\em Network} object that
should be output to screen or log files prior to an {\em Iterative}
run,
you will need to add code to the runHeader method in {\em Network}.
See \hyperref[runHeaderSection]{section \ref{runHeaderSection}}.

\subsubsection{pack and unpack}\label{packutil}

The methods

\begin{verbatim}
virtual void pack()
virtual void unpack()
\end{verbatim}

must be coded to convert the weight gradient structure of your {\em
Network} derived object to the member vector $stackG$ (pack), and vice
versa (unpack).  See SimpleProp::pack, SimpleProp::unpack,
BareProp::pack and BareProp::unpack as examples.

\subsection{Built-in gradient descent algorithms}

One of the really neat things about {\em Networks} that construct the
$stackG$ vector is that they may use the built-in advanced gradient
descent algorithms, Shanno's\index{Shanno algorithm} and conjugate
gradient descent\index{Conjugate gradient descent}.\footnote{Richard
M. Golden, Mathematical Methods for Neural Network Analysis and
Design, 1996, MIT Press, ISBN 0-262-07174-6, pp. 217--218, 221--222}
All the programmer needs to do is to construct the $stackG$ vector,
add the line of code
\begin{verbatim}
engine( trainingType, iteration );
\end{verbatim}
after the {\em Network} object gradient structure is calculated and
before the weights are updated, and these algorithms will
automagically be available for the user.  See SimpleProp::trainSet()
and BareProp::trainSet() for examples.

\subsubsection{Shanno algorithm}\index{Shanno algorithm}

Briefly, Shanno's algorithm proceeds as follows: let the computed
gradient be $\vec{g}$ ($stackG$), and the iteration be $t$ ({\bf
Note:} the iteration is {\em not} the exemplar index through the
training set. It is the actual iteration number for each calculated
gradient.  It should also be noted that this algorithm, as well as
conjugate gradient descent {\em is intended for off-line or
batch/epoch learning.}) Let the learning rate be $\gamma$.

\begin{itemize}
\item Step 0: Let $t=0$.
\item Step 1: Set $\vec{f}(t)=-\vec{g}_t$.  Go to step 3.
\item Step 2: If $t$ is a multiple of the degrees of freedom of the
      {\em Network}, then go to step 1.  Otherwise, compute
      $\vec{u}_t=\vec{g}_t-\vec{g}_{t-1}$ and the scalars $a_t$,
      $b_t$, $c_t$ such that
\begin{itemize}
\item $a_t=\frac{\vec{f}(t-1)^T\vec{g_t}}{\vec{f}(t-1)^T\vec{u}_t}$,
\item $b_t=\frac{\vec{u}^T\vec{g_t}}{\vec{f}(t-1)^T\vec{u}_t}$,
\item $c_t=\gamma+\frac{|\vec{u}_t|^2}{\vec{f}(t-1)^T\vec{u}_t}$,
\end{itemize}
and descent vector
\begin{displaymath}
\vec{f}(t)=-\vec{g}_t+a_t\vec{u}_t+(b_t-c_ta_t)\vec{f}(t-1)
\end{displaymath}
\item Step 3: Update the weights by
$\vec{w}(t+1)=\vec{w}(t)+\gamma\vec{f}(t)$.
Let $t=t+1$.  Go to step 2.
\end{itemize}

Interested readers are directed to Richard Golden's book, pp.
217--218. It should be noted that in the highly optimized code in
Network::shanno(unsigned $iteration$), $lastF$ refers to
$-\vec{f}(t-1)$.

\subsubsection{Conjugate gradient descent}\index{Conjugate gradient
descent}

Conjugate gradient descent is a special case of the Shanno algorithm,
and briefly proceeds as follows: let the computed gradient be
$\vec{g}$ ($stackG$), and the iteration be $t$ ({\bf Note:} again, the
iteration is {\em not} the exemplar index through the training set. It
is the actual iteration number for each calculated gradient.  It
should also be noted that this algorithm, as well as Shanno's, {\em is
intended for off-line or batch/epoch learning.}) Let the learning rate
be $\gamma$.

\begin{itemize}
\item Step 0: Let $t=0$.
\item Step 1: Set $\vec{f}(t)=-\vec{g}_t$.  Go to step 3.
\item Step 2: If $t$ is a multiple of the degrees of freedom of the
      {\em Network}, then go to step 1.  Otherwise, compute
      $\vec{u}_t=\vec{g}_t-\vec{g}_{t-1}$ and the scalar $b_t$
      such that
\begin{displaymath}
b_t=\frac{\vec{u}^T\vec{g_t}}{\vec{u}^T_t\vec{f}(t-1)}
\end{displaymath}
and descent vector
\begin{displaymath}
\vec{f}(t)=-\vec{g}_t+b_t\vec{f}(t-1)
\end{displaymath}
\item Step 3: Update the weights by
$\vec{w}(t+1)=\vec{w}(t)+\gamma\vec{f}(t)$.
Let $t=t+1$.  Go to step 2.
\end{itemize}

Interested readers are directed to Richard Golden's book, pp.
221--222. It should be noted that in the highly optimized code in
Network::CGD(unsigned $iteration$), $lastF$ refers to
$-\vec{f}(t-1)$.

\subsection{Calculation of the condition number}

The overloaded storeGrads($\ldots$) method, the reportCondNum(
ostream\&
$outputStream$) method and the $finalFlag$ flag may be used to compute
the condition number of a {\em Network} object.  The condition number
is the ratio of the largest to the smallest eigenvalue of the
symmetric matrix ${\bf B}$ computed by the outer product of the
average gradient vector and its transpose:

\begin{displaymath}
{\bf B} = {\bf g} {\cdot} {\bf g}^T
\end{displaymath}

The {\em Gnu Scientific Library (GSL)} eigenvalue routines were ported
for use in these methods.  Note that condition numbers of neural
networks are highly experimental, and in the current driver, the
condition number is calculated automatically only for logistic
regression (e.g. {\em Logistic}) type {\em Network} objects.  However,
implementation of condition number calculation is as general utility
methods in the {\em Network} class so that future programmers may
report
condition numbers of {\em Network} objects if desired.

Computation of the condition number for a {\em Network} derived object
proceeds as follows:

\begin{enumerate}
\item The overloaded storeGrads( unsigned $d$, unsigned $n$ ) method
is
      first used to initialize a gradient accumulator matrix ($grads$,
      which is transparent to the programmer and does not need to be
      included in added code) which will ultimately be used to compute
      the ${\bf B}$ matrix.  Arguments for this overloaded method are
      $d$, the dimension (size) of the gradient vector, and $n$, the
      number of exemplars in the dataset.
\item The gradient vector $stackG$ must be created using the $pack()$
      method.
\item Because methods involved in computation of the condition number
      are highly intensive, and because the condition number only
      requires
      the final average gradient vector, the $finalFlag$ flag is used
      to indicate a final iteration for the following step $\ldots$
\item The overloaded storeGrads( unsigned $n$ ) method is used to
      accumulate exemplar $n$'s gradient $stackG$ (which does {\em
      not} need to be passed as an argument, as it must be generated
      in code prior to use of this method.)  Note that the {\em
      Network} method reportCondNum(ostream\& $outputStream$) will
      automatically compute the average gradient, hence that
      particular exemplar's gradient is loaded (again, transparent to
      the programmer.)
\item reportCondNum(ostream\& $outputStream$) is used in methods which
      report results to the user, e.g. reportAccuracy( ostream\&
      $outputStream$ )
\end{enumerate}

See {\em Logistic} for example code, and see {\em SimpleProp} for
comments which indicate where code would be placed in a neural network
style implemented {\em Network} object.

\section{Network}

In the following specification, methods and members that pertain
directly to coding a concrete derived {\em Network} implementation are
shown in \beginred red \endcolor.

{\bf Network} extends {\em Iterative}. Algorithms with hidden layers,
weights, learning rates, transfer functions and other general
properties of neural networks are encapsulated in {\em Network}.

{\bf Note:} It is expected that in objects derived from {\em Network},
the method setDataSet inherited from {\em Model} be called {\em
before} the derived methods that specify object architecture, so that
the derived object may know the number of input nodes and size the
Matrices related to the dataset prior to completing its architecture.
This sequence of loading a dataset using the inherited {\em Model}
method setDataSet prior to establishing network architecture will be
invoked for all {\em Network} objects.

\begin{itemize}

\item {\em Public} methods

\begin{itemize}
\item Network\& operator= ( const Network\& rhs ): see section
      \ref{copy_fx} for discussion
\item Network( const Network\& rhs ): see section \ref{copy_fx} for
      discussion
\item \beginred virtual void setDataSet( DataSet\& $data$ ): A pure
      virtual method inherited from {\em Model} that loads the {\em
      DataSet} object $data$ into the {\em Model} \endcolor
\item \beginred virtual void outputHeader( ostream\& $outputStream$ ):
      A pure virtual method inherited from {\em Model} which outputs a
      header to an ostream $outputStream$ describing the architecture
      of the model, {\em Model} type dependant \endcolor
\item void setEta( double $\eta$ ): Sets the learning rate to $\eta$
\item double getEta(): Returns the learning rate $\eta$
\item void setBias( bool $flag$ ): Sets bias nodes in {\em Network}
\item bool getBias(): Gets boolean value indicating if bias nodes
      specified in {\em Network}
\item void setBatchEpoch( bool $flag$ ): Sets if batch/epoch learning
      is used in a {\em Network} object (See the description of
      classic off-line (``batch'') backprop in
      \hyperref[learningAlgs]{section \ref{learningAlgs}})
      \index{Backpropagation!Off-line}
\item bool getBatchEpoch(): returns if batch/epoch used
      \index{Backpropagation!Off-line}
\item void setWeightDecay( bool $flag$ ): Sets if weight decay is used
      in a {\em Network} object (See the discussion of
      \hyperref[discussDecay]{weight decay} in section
      \ref{decayFlagSection})\index{Weight decay}
\item bool getWeightDecay(): returns if weight decay used
      \index{Weight decay}
\item void setDecay( double $2\lambda$ ): Sets weight decay term to
      $2\lambda$ in a {\em Network} object\index{Weight decay}
\item double getDecay(): returns weight decay term $\lambda$
      \index{Weight decay}
\item void setTrainingType( unsigned $T$ ): Sets training type to $T$
\item unsigned getTrainingType(): returns training type
\item void setAutoStepSize( bool $flag$ ): Sets if automatic stepsize
      selection is used (See the discussion of
      \hyperref[autoStep]{automatic stepsize selection} in section
      \ref{autoStep})\index{Automatic stepsize selection}
\item bool getAutoStepSize(): returns if automatic stepsize selection
      is used
\item \beginred virtual unsigned df(): Returns the number of degrees
      of freedom, equal to the number of nodal connections \endcolor
\item void setRandomLimit( double $L$ ): Sets the limit for the random
      weights $-L$ to $+L$
\item \beginred virtual void randomize(): A pure virtual method which
      randomizes initial weights in a {\em Network} object \endcolor
\item getWeightsSet(): Returns a flag to indicate if weights set
\item \beginred virtual bool save( string\& $filename$ ): A pure
      virtual method which saves the {\em Network} object architecture
      and weights to a file named $filename$, returns true if
      successful \endcolor
\item \beginred virtual bool load( string\& $filename$ ): A pure
      virtual method which loads the {\em Network} object architecture
      and weights from a file named $filename$, returns true if
      successful \endcolor
\item \beginred virtual double trainSet(): A pure virtual method which
      trains one iteration through the training set, {\em $Network$}
      type dependant, and returns the set error \endcolor
\item \beginred virtual double innerTrainSet(): A pure virtual method
      which is called by trainSet() to calculate the gradient and
      return set error, used for automatic stepsize selection
      \endcolor
\item \beginred virtual void forward( Matrix$<$ double $>$\& $M$,
      unsigned $n$ ): A pure virtual method which forward propagates
      for one input row $n$ in a Matrix of inputs $M$ (e.g. either the
      inherited {\em Model} Matrix $Train$ or $Test$) \endcolor
      \index{Forward propagation}
\item virtual void reportAccuracy( ostream\& $outputStream$ ):
      Implementation of the pure virtual method inherited from {\em
      Model} that reports the accuracy of the trained model to an
      ostream $outputStream$
\item virtual void classAccuracy( ostream\& $outputStream$ ): A pure
      virtual method inherited from {\em Iterative} which outputs
      classification accuracies to an ostream $outputStream$ for an
      {\em Iterative} output table
\item \beginred virtual void removeInputs( vector$<$ unsigned $>$\&
      $v$ ): A pure virtual method which takes as its argument a
      vector of unsigned $v$ representing positions of input nodes,
      and removes those nodes from the {\em Network} object \endcolor
\end{itemize}

\item {\em Protected} members

\begin{itemize}
\item \beginred unsigned $trainingType$: the type of gradient descent
      algorithm, $trainingType = 0$ is reserved for canonical
      backpropagation ($trainingType = 1$ has been coded for conjugate
      gradient descent, and $trainingType = 2$ for Shanno algorithm)
      \endcolor
\item \beginred double $o$: the output for a single-output {\em
      Network} \endcolor
\item \beginred double $x$: the $x$ term for the sigmoidal output $o$
      in a single output network:\endcolor
      \index{Transfer function!Sigmoidal}
\begin{displaymath}
o = 1/(1+e^{-x})
\end{displaymath}
\item \beginred vector $<$ double $>$ $output$: the outputs vector for
      a multiple output {\em Network} \endcolor
\item \beginred vector $<$ double $>$ $xs$: the $\vec{\bf x}$ terms
      for the sigmoidal outputs $\vec{\bf o}$ in a multiple output
      network: \endcolor\index{Transfer function!Sigmoidal}
\begin{displaymath}
\vec{\bf o} = 1/(1+e^{-\vec{\bf x}})
\end{displaymath}
\item \beginred double $randomLimit$: the limit for the random weights
      \endcolor
\item \beginred bool $weightsSetFlag$: a flag to indicate if weights
      set \endcolor
\item \beginred bool $biasFlag$: a flag to indicate if bias nodes
      specified in {\em Network} \endcolor
\item \beginred bool $batchEpochFlag$: a flag to indicate if batch/
      epoch training is used in the {\em Network} object \endcolor (
      See the description of classic off-line (``batch'') backprop in
      \hyperref[learningAlgs]{section \ref{learningAlgs}})
      \index{Backpropagation!Off-line}
\item \beginred bool $weightDecayFlag$: a flag to indicate if weight
      decay is used in the {\em Network} object \endcolor (See the
      discussion of \hyperref[discussDecay]{weight decay} in section
      \ref{decayFlagSection})\index{Weight decay}
\item \beginred double $decay$: the weight decay term $2\lambda$
\endcolor
      \index{Weight decay}
\item \beginred double $regularizer$: the weight decay term $\lambda$
      \endcolor \index{Weight decay}
\item \beginred double $decayTerm$: the weight decay multiple
      $1-2\eta\lambda$ calculated in the protected member function
      runHeader (\hyperref[decayTermSect]{see section
      \ref{decayTermSect} for an explanation}) \endcolor
      \index{Weight decay}
\item \beginred double $eta$: the learning rate $\eta$ \endcolor
\item \beginred bool $automaticStepSizeFlag$: a flag to indicate if
      automatic step size algorithm is used \endcolor
\item \beginred double $o\_errAccumulate$: stores average output error
      accumulation over the examplars used in innerTrainSet() method
      \endcolor
\item \beginred double $oldErrorAccumulate$: stores average output
      error accumulation for previous run over the examplars, used in
      innerTrainSet() method \endcolor
\item \beginred double $deltaError$: deltaError constant for automatic
      stepsize selection \endcolor
\item double $currGradMax$: current absolute maximum gradient for
      shanno/CGD
\item \beginred double $gamma$: gamma constant for automatic stepsize
      selection \endcolor
\item \beginred unsigned $maxLoops$: maxLoops constant for automatic
      stepsize selection \endcolor
\item \beginred vector $<$ double $>$ $stackG$: a single vector which
      contains the entire weight gradient of the {\em Network} derived
      object \endcolor
\item vector $<$ double $>$ $lastG$: $\vec{g}(t-1)$ for shanno/CGD
\item vector $<$ double $>$ $lastF$: $\vec{f}(t-1)$ for shanno/CGD
\item \beginred bool $finalFlag$: a flag to indicate that the train
      methods (innerTrainSet() or trainSet()) are at the final
      iteration,
      for code to be run only during the final iteration \endcolor
\item Matrix $<$ double $>$ $grads$: a matrix to hold exemplar
gradients
      to be used to compute the B matrix to calculate the condition
      number

\end{itemize}

\item {\em Protected} Utility functions

\begin{itemize}
\item \beginred void copy( const Network\& rhs ): Copy utility
      function, see section \ref{copy_fx} for discussion. \endcolor
\item virtual void runHeader( ostream\& $outputStream$ ): A pure
      virtual utility function inherited from the Iterative class
      that outputs Network specific parameters preceeding an
      Iterative run
\item \beginred virtual void pack(): A pure virtual utility function
      that converts the entire weight gradient structure of the {\em
      Network} derived object to the protected member $stackG$
      \endcolor
\item \beginred virtual void unpack(): A pure virtual utility function
      that converts the protected member $stackG$ back to the weight
      gradient structure of the {\em Network} derived object\endcolor
\item virtual double getGradMax(): A pure virtual utility function
      inhertied from {\em Iterative} that returns the maximum absolute
      gradient
\item void engine( unsigned $type$, unsigned $iteration$ ): an utility
      function which calls the appropriate training algorithm
      depending on training type $type$, and passes it the current
      iteration
\item void engine( unsigned $type$, unsigned $iteration$ ): chooses
which
      kind of advanced gradient descent algorithm (CGD or shanno)
      based
      on $type$, and passes $iteration$ to it
\item void CGD( unsigned $iteration$ ): implements conjugate gradient
      descent (see Golden pp. 221--222)\footnote{Richard M. Golden,
      Mathematical Methods for Neural Network Analysis and Design,
      1996, MIT Press, ISBN 0-262-07174-6}
\item void shanno( unsigned $iteration$ ): implements Shanno algorithm
      (see Golden pp. 217--218)
\item void storeGrads( unsigned $d$, unsigned $n$ ): (used in
      computing condition number) overloaded utility method (see
      below), with 2 arguments $d$ and $n$, initializes $grads$ matrix
      with dimension $d$ of packed gradient (size), and number of
      exemplars $n$
\item void storeGrads( unsigned $n$ ): (used in computing condition
      number) overloaded utility method (see above), with 1 argument
      $n$ accumulate grads matrix with that exemplar $n$'s gradient
      $stackG$ (note that $stackG$ must have been previously packed
      using pack()
\item void reportCondNum( ostream\& $outputStream$ ): utility method
      to report out the condition number to $outputStream$

\end{itemize}

\end{itemize}

\section{BareProp}\index{Backpropagation|(}

The {\bf BareProp} class is a {\em Network} object that serves to
illustrate neural network design in neUROn2++.  {\em BareProp} is a
one output node, one hidden layer neural network without biases.

{\bf Note:} It is expected that the method setDataSet inherited from
{\em Model} be called {\em before} the {\em BareProp} method
setHidden, so that the {\em BareProp} object may know the number of
input nodes and size the Matrices related to the dataset prior to
completing its architecture.  This sequence of loading a dataset using
the {\em Model} method setDataSet prior to establishing network
architecture will be invoked for all {\em Network} objects.

\begin{itemize}

\item {\em Public} methods

\begin{itemize}
\item BareProp\& operator= ( const BareProp\& rhs ): see section
      \ref{copy_fx} for discussion
\item BareProp( const BareProp\& rhs ): see section \ref{copy_fx} for
      discussion
\item virtual void setDataSet( DataSet\& $data$ ): Implementation of
      the pure virtual method inherited from {\em Model} that loads
      the {\em DataSet} object $data$ into the {\em BareProp} {\em
      Model}
\item virtual void outputHeader( ostream\& $outputStream$ ):
      Implementation of the pure virtual method inherited from {\em
      Model} which outputs a header to an ostream $outputStream$
      describing the architecture of the {\em Bareprop} {\em Model}
\item void setHidden( unsigned $H$ ): Sets number of hidden nodes to
      $H$, which is all that is required to specify the architecture
      for a {\em BareProp} object
\item virtual unsigned df(): Implementation of the pure virtual method
      inherited from {\em Network} which returns the number of degrees
      of freedom, equal to the number of nodal connections
\item virtual void randomize(): Implementation of the pure virtual
      method inherited from {\em Network} which randomizes initial
      weights
\item virtual bool save( string\& $filename$ ): Implementation of the
      pure virtual method inherited from {\em Network} which saves the
      {\em BareProp} object architecture and weights to a file named
      $filename$, returns true if successful
\item virtual bool load( string\& $filename$ ): Implementation of the
      pure virtual method inherited from {\em Network} which loads the
      {\em BareProp} object architecture and weights from a file named
      $filename$, returns true if successful
\item virtual double trainSet(): Implementation of the pure virtual
      method inherited from {\em Iterative} which trains one iteration
      through the training set, and returns set error
\item virtual double innerTrainSet(): Implementation of the pure
      virtual method inherited from {\em Network} which is called by
      trainSet() to calculate the gradient and return set error, used
      for automatic stepsize selection
\item virtual void forward( Matrix$<$ double $>$\& $M$, unsigned $n$
      ): Implementation of the pure virtual method inherited from {\em
      Network} which forward propagates for one input row $n$ in a
      Matrix of inputs $M$ (e.g. either the inherited {\em Model}
      Matrix $Train$ or $Test$)\index{Forward propagation}
\item virtual void removeInputs( vector$<$ unsigned $>$\&
      $v$ ): Implementation of the pure virtual method which takes as
      its argument a vector of unsigned $v$ representing positions of
      input nodes, and removes those nodes from the {\em BareProp}
      object
\end{itemize}

\item {\em Private} members

\begin{itemize}
\item unsigned $nHidden$: number of hidden nodes
\item Matrix$<$ double $>$ $hW$: hidden weight Matrix
\item Matrix$<$ double $>$ $hWup$: hidden weight update Matrix
\item Matrix$<$ double $>$ $hG$: hidden weight gradient Matrix
\item vector$<$ double $>$ $y$: vector containing dataset outputs
\item vector$<$ double $>$ $I$: vector for single exemplar inputs
\item vector$<$ double $>$ $hO$: hidden output vector
\item vector$<$ double $>$ $oW$: output weight vector
\item vector$<$ double $>$ $oG$: output weight gradient vector
\item vector$<$ double $>$ $h\_err$: hidden error vector
\item double $o\_err$: output error term
\end{itemize}

\item {\em Private} Utility functions

\begin{itemize}
\item void copy( const BareProp\& rhs ): Copy utility function, see
      section \ref{copy_fx} for discussion.
\item virtual void pack(): Implementation of the pure virtual utility
      function inherited from {\em Network} that converts the entire
      weight gradient structure to the {\em Network} inherited
      protected member $stackG$
\item virtual void unpack(): Implementation of the pure virtual
      utility function inherited from {\em Network} that converts the
      {\em Network} inherited protected member $stackG$ back to the
      weight gradient structure
\end{itemize}

\end{itemize}

\section{SimpleProp}

The {\bf SimpleProp} class is a {\em Network} object that serves to
illustrate neural network design in neUROn2++.  {\em SimpleProp} is a
one output node, one hidden layer neural network with biases.

{\bf Note:} It is expected that the method setDataSet inherited from
{\em Model} be called {\em before} the {\em SimpleProp} method
setHidden, so that the {\em SimpleProp} object may know the number of
input nodes and size the Matrices related to the dataset prior to
completing its architecture.  This sequence of loading a dataset using
the {\em Model} method setDataSet prior to establishing network
architecture will be invoked for all {\em Network} objects.

\begin{itemize}

\item {\em Public} methods

\begin{itemize}
\item SimpleProp\& operator= ( const SimpleProp\& rhs ): see section
      \ref{copy_fx} for discussion
\item SimpleProp( const SimpleProp\& rhs ): see section \ref{copy_fx}
      for discussion
\item virtual void setDataSet( DataSet\& $data$ ): Implementation of
      the pure virtual method inherited from {\em Model} that loads
      the {\em DataSet} object $data$ into the {\em SimpleProp} {\em
      Model}
\item virtual void outputHeader( ostream\& $outputStream$ ):
      Implementation of the pure virtual method inherited from {\em
      Model} which outputs a header to an ostream $outputStream$
      describing the architecture of the {\em Bareprop} {\em Model}
\item void setHidden( unsigned $H$ ): Sets number of hidden nodes to
      $H$, which is all that is required to specify the architecture
      for a {\em SimpleProp} object
\item virtual unsigned df(): Implementation of the pure virtual method
      inherited from {\em Network} which returns the number of degrees
      of freedom, equal to the number of nodal connections
\item virtual void randomize(): Implementation of the pure virtual
      method inherited from {\em Network} which randomizes initial
      weights
\item virtual bool save( string\& $filename$ ): Implementation of the
      pure virtual method inherited from {\em Network} which saves the
      {\em SimpleProp} object architecture and weights to a file named
      $filename$, returns true if successful
\item virtual bool load( string\& $filename$ ): Implementation of the
      pure virtual method inherited from {\em Network} which loads the
      {\em SimpleProp} object architecture and weights from a file
      named
      $filename$, returns true if successful
\item virtual double trainSet(): Implementation of the pure virtual
      method inherited from {\em Iterative} which trains one iteration
      through the training set, and returns set error
\item virtual double innerTrainSet(): Implementation of the pure
      virtual method inherited from {\em Network} which is called by
      trainSet() to calculate the gradient and return set error, used
      for automatic stepsize selection
\item virtual void forward( Matrix$<$ double $>$\& $M$, unsigned $n$
      ): Implementation of the pure virtual method inherited from {\em
      Network} which forward propagates for one input row $n$ in a
      Matrix of inputs $M$ (e.g. either the inherited {\em Model}
      Matrix $Train$ or $Test$)\index{Forward propagation}
\item virtual void removeInputs( vector$<$ unsigned $>$\&
      $v$ ): Implementation of the pure virtual method which takes as
      its argument a vector of unsigned $v$ representing positions of
      input nodes, and removes those nodes from the {\em SimpleProp}
      object
\end{itemize}

\item {\em Private} members

\begin{itemize}
\item unsigned $nHidden$: number of hidden nodes
\item unsigned $nH$: $nHidden - 1$, used for indexing with biases
\item Matrix$<$ double $>$ $hW$: hidden weight Matrix
\item Matrix$<$ double $>$ $hWup$: hidden weight update Matrix
\item Matrix$<$ double $>$ $hG$: hidden weight gradient Matrix
\item vector$<$ double $>$ $in\_bias$: input bias vector to add to
      input Matrix
\item vector$<$ double $>$ $y$: vector containing dataset outputs
\item vector$<$ double $>$ $I$: vector for single exemplar inputs
\item vector$<$ double $>$ $hO$: hidden output vector
\item vector$<$ double $>$ $oW$: output weight vector
\item vector$<$ double $>$ $oG$: output weight gradient vector
\item vector$<$ double $>$ $h\_err$: hidden error vector
\item double $o\_err$: output error term
\end{itemize}

\item {\em Private} Utility functions

\begin{itemize}
\item void copy( const SimpleProp\& rhs ): Copy utility function, see
      section \ref{copy_fx} for discussion.
\item virtual void pack(): Implementation of the pure virtual utility
      function inherited from {\em Network} that converts the entire
      weight gradient structure to the {\em Network} inherited
      protected member $stackG$
\item virtual void unpack(): Implementation of the pure virtual
      utility function inherited from {\em Network} that converts the
      {\em Network} inherited protected member $stackG$ back to the
      weight gradient structure
\end{itemize}

\end{itemize}

\section{BackProp}

The {\bf BackProp} class is a {\em Network} object that allows
multiple output nodes and multiple hidden layers in a neural network,
and biases may or may not be present.

{\bf Note:} It is expected that the {\em Model} method setDataSet be
called {\em before} the {\em BackProp} method setHidden, so that the
{\em BackProp} object may know the number of input nodes and size the
Matrices related to the dataset prior to completing its architecture.
This sequence of loading a dataset using the {\em Model} method
setDataSet prior to establishing network architecture will be invoked
for all {\em Network} objects.

\begin{itemize}

\item {\em Public} methods

\begin{itemize}
\item BackProp\& operator= ( const BackProp\& rhs ): see section
      \ref{copy_fx} for discussion
\item BackProp( const BackProp\& rhs ): see section \ref{copy_fx}
      for discussion
\item virtual void setDataSet( DataSet\& $data$ ): Implementation of
      the pure virtual method inherited from {\em Model} that loads
      the {\em DataSet} object $data$ into the {\em BackProp} {\em
      Model}
\item virtual void outputHeader( ostream\& $outputStream$ ):
      Implementation of the pure virtual method inherited from {\em
      Model} which outputs a header to an ostream $outputStream$
      describing the architecture of the {\em Backprop} {\em Model}
\item void setHidden( const vector$<$ unsigned $>$\& $v$ ): Sets the
      hidden layers, with the number of hidden nodes on each layer
      indicated by the corresponding vector $v$ element
\item virtual unsigned df(): Implementation of the pure virtual method
      inherited from {\em Network} which returns the number of degrees
      of freedom, equal to the number of nodal connections
\item virtual void randomize(): Implementation of the pure virtual
      method inherited from {\em Network} which randomizes initial
      weights
\item virtual bool save( string\& $filename$ ): Implementation of the
      pure virtual method inherited from {\em Network} which saves the
      {\em BackProp} object architecture and weights to a file named
      $filename$, returns true if successful
\item virtual bool load( string\& $filename$ ): Implementation of the
      pure virtual method inherited from {\em Network} which loads the
      {\em BackProp} object architecture and weights from a file named
      $filename$, returns true if successful
\item virtual double trainSet(): Implementation of the pure virtual
      method inherited from {\em Iterative} which trains one iteration
      through the training set, and returns set error
\item virtual double innerTrainSet(): Implementation of the pure
      virtual method inherited from {\em Network} which is called by
      trainSet() to calculate the gradient and return set error, used
      for automatic stepsize selection
\item virtual void forward( Matrix$<$ double $>$\& $M$, unsigned $n$
      ): Implementation of the pure virtual method inherited from {\em
      Network} which forward propagates for one input row $n$ in a
      Matrix of inputs $M$ (e.g. either the inherited {\em Model}
      Matrix $Train$ or $Test$)\index{Forward propagation}
\item virtual void removeInputs( vector$<$ unsigned $>$\&
      $v$ ): Implementation of the pure virtual method which takes as
      its argument a vector of unsigned $v$ representing positions of
      input nodes, and removes those nodes from the {\em BackProp}
      object
\end{itemize}

\item {\em Private} members

\begin{itemize}
\item vector$<$ unsigned $>$ $nLayer$: describes the number of hidden
      nodes on each layer indicated by the corresponding vector
      element in $nLayer$
\item unsigned $nLayers$: number of hidden layers
\item vector$<$ double $>$ $in\_bias$: input bias vector to add to
      input Matrix if bias nodes present
\item vector$<$ double $>$ $I$: vector for single exemplar inputs
\item vector$<$ double $>$ $y$: vector containing dataset outputs
\item vector$<$ double $>$ $o\_err$: vector containing output errors
\item vector$<$ Matrix$<$ double $>$ $>$ $Weights$: vector of weight
      Matrix for all hidden layers and output layer
\item vector$<$ Matrix$<$ double $>$ $>$ $WeightsUp$: vector of
      Matrices to hold updated weights one for each weights Matrix
\item vector$<$ Matrix$<$ double $>$ $>$ $WeightsAccumulate$: vector
      of Matrices to hold accumulated weight updates for batch/epoch
      learning\index{Backpropagation!Off-line}
\item vector$<$ vector$<$ double $>$ $>$ $HOutputs$: vector containing
      a vector of outputs for each hidden layer
\item vector$<$ vector$<$ double $>$ $>$ $HErrors$: vector containing
      a vector of errors for each hidden layer
\end{itemize}

\item {\em Private} Utility functions

\begin{itemize}
\item void copy( const BackProp\& rhs ): Copy utility function, see
      section \ref{copy_fx} for discussion.
\item virtual void pack(): Implementation of the pure virtual utility
      function inherited from {\em Network} that converts the entire
      weight gradient structure to the {\em Network} inherited
      protected member $stackG$
\item virtual void unpack(): Implementation of the pure virtual
      utility function inherited from {\em Network} that converts the
      {\em Network} inherited protected member $stackG$ back to the
      weight gradient structure
\end{itemize}

\end{itemize}
\index{Backpropagation|)}

\section{Logistic}\index{Logistic regression}

The {\bf Logistic} class is a {\em Network} object that implements
logistic regression. This form of data modeling is mathematically
equivalent to a single output node neural network with a bias, the
sigmoidal transfer function, no hidden nodes, and the cross-entropy
error function.  As a result, logistic regression is implemented as a
derived Network class.  This strategy also allows the user to perform
stepwise forward and reverse regression on the built logistic
regression model using the \hyperref[regressnet] {RegressNet class}.

Note that the error is \hyperref[xentropyfunk]{cross-entropy} by
definition, that the gradient descent method is \hyperref[
batchepochdef]
{batch-epoch} by definition, and a bias node (y-intercept) is present
by definition.  At the end of a run, the $\beta$ weights and standard
error are reported and $p$-values for the Wald test that the $\beta$
weights are nonzero.\footnote{see Hosmer and Lemeshow, Applied
Logistic Regression, 2nd edition, John Wiley \& Sons, Inc., New York,
2000, pp 36--42}\index{Hosmer \& Lemeshow}

{\bf Note:} It is expected that the method setDataSet inherited from
{\em Model} be called {\em before} the {\em Logistic} method randomize
() (which sets initial weights to zero or random values, depending on
the value of the flag {\em randInitCondFlag} which is set by the
method setInitCond()), so that the {\em Logistic} object may know the
number of input nodes and size the vectors and Matrices related to the
dataset prior to completing its architecture.  This sequence of
loading a dataset using the {\em Model} method setDataSet prior to
establishing network architecture will be invoked for all {\em
Network} objects.

\begin{itemize}

\item {\em Public} methods

\begin{itemize}
\item Logistic\& operator= ( const Logistic\& rhs ): see section
      \ref{copy_fx} for discussion
\item Logistic( const Logistic\& rhs ): see section \ref{copy_fx} for
      discussion
\item virtual void setDataSet( DataSet\& $data$ ): Implementation of
      the pure virtual method inherited from {\em Model} that loads
      the {\em DataSet} object $data$ into the {\em Logistic} {\em
      Model}
\item virtual void outputHeader( ostream\& $outputStream$ ):
      Implementation of the pure virtual method inherited from {\em
      Model} which outputs a header to an ostream $outputStream$
      describing the architecture of the {\em Logistic} {\em Model}
\item virtual unsigned df(): Implementation of the pure virtual method
      inherited from {\em Network} which returns the number of degrees
      of freedom, equal to the number of inputs $+$ 1 for the y-
      intercept
\item virtual void randomize(): Implementation of the pure virtual
      method inherited from {\em Network} which randomizes initial
      weights
\item virtual void reportAccuracy( ostream\& $outputStream$ ):
      overrides (calls base {\em Network} reportAccuracy() method
      first) {\em Network} method inherited from {\em Iterative},
      adding binary logistic regression specific footer information,
      including reporting $\beta$ weights, standard error, and Wald
      test ($p$ for $\beta$ weights nonzero)
\item virtual bool save( string\& $filename$ ): Implementation of the
      pure virtual method inherited from {\em Network} which saves the
      {\em Logistic} object $\beta$ weights and y-intercept to a file
      named $filename$, returns true if successful
\item virtual bool load( string\& $filename$ ): Implementation of the
      pure virtual method inherited from {\em Network} which loads the
      {\em Logistic} object $\beta$ weights and y-intercept from a
      file named $filename$, returns true if successful
\item virtual double trainSet(): Implementation of the pure virtual
      method inherited from {\em Iterative} which trains one iteration
      through the training set, and returns set error
\item virtual double innerTrainSet(): Implementation of the pure
      virtual method inherited from {\em Network} which is called by
      trainSet() to calculate the gradient and return set error, used
      for automatic stepsize selection
\item virtual void forward( Matrix$<$ double $>$\& $M$, unsigned $n$):
      Implementation of the pure virtual method inherited from {\em
      Network} which applies the sigmoidal function for one input row
      $n$ in a Matrix of inputs $M$ (e.g. either the inherited {\em
      Model} Matrix $Train$ or $Test$)\index{Forward propagation}
\item virtual void removeInputs( vector$<$ unsigned $>$\&
      $v$ ): Implementation of the pure virtual method which takes as
      its argument a vector of unsigned $v$ representing positions of
      input nodes, and removes those nodes from the {\em Logistic}
      object
\end{itemize}

\item {\em Private} members

\begin{itemize}
\item vector$<$ double $>$ $y$: vector containing dataset outputs
\item vector$<$ double $>$ $I$: vector for single exemplar inputs
\item vector$<$ double $>$ $W$: vector for $\beta$ weights (last is
      y-intercept)
\item vector$<$ double $>$ $G$: vector for $\beta$ weight gradients
\item double $o\_err$: output error term
\end{itemize}

\item {\em Private} Utility functions

\begin{itemize}
\item void copy( const Logistic\& rhs ): Copy utility function, see
      section \ref{copy_fx} for discussion.
\item virtual void pack(): Implementation of the pure virtual utility
      function inherited from {\em Network} that converts the $\beta$
      weight gradient structure to the {\em Network} inherited
      protected member $stackG$, as the $\beta$ weight structure is
      already
      a single vector, this simply copies it
\item virtual void unpack(): Implementation of the pure virtual
      utility function inherited from {\em Network} that converts the
      {\em Network} inherited protected member $stackG$ back to the
      weight gradient structure, which, since it is a vector, just
      copies
      it
\end{itemize}

\end{itemize}

\section{DFA}

The {\bf DFA} class is a {\em Model} object that encodes discriminant
function analysis.  It is an abstract class, whose derived concrete
classes are {\em LDFA} for linear discriminant function analysis, and
{\em QDFA} for quadratic discriminant function analysis.\footnote{ An
excellent treatment of discriminant function analysis can be found in
Richard O. Duda and Peter E. Hart, Pattern Classification and Scene
Analysis, 1973, John Wiley \& Sons, New York, ISBN 0-471-22361-1,
pp. 24-32}

\begin{itemize}

\item {\em Public} methods

\begin{itemize}
\item DFA\& operator= ( const DFA\& rhs ): see section
      \ref{copy_fx} for discussion
\item DFA( const DFA\& rhs ): see section \ref{copy_fx} for
      discussion
\item virtual void setDataSet( DataSet\& $data$ ): Implementation of
      the pure virtual method inherited from {\em Model} that loads
      the {\em DataSet} object $data$ into the {\em DFA} {\em Model}
\item virtual void outputHeader( ostream\& $outputStream$ ):
      Implementation of the pure virtual method inherited from {\em
      Model} which outputs a header to an ostream $outputStream$
      describing the architecture of the {\em DFA} {\em Model}
\item \beginred virtual double train(): A pure virtual method
      inherited from {\em Model} that trains a {\em Model}, and
      returns the final error (double) for that model \endcolor
\item \beginred virtual void reportAccuracy( ostream\& $outputStream$
      ): A pure virtual method inherited from {\em Model} that reports
      the accuracy of the trained model to an ostream $outputStream$
      \endcolor
\end{itemize}

\item {\em Protected} members

\begin{itemize}
\item Matrix$<$ double $>$ $D0$, $D1$: Matrices for inputs of 1-output
      datasets
\item vector$<$ Matrix$<$ double $>$ $>$ $D$: a vector of Matrices for
      inputs of multiple output datasets
\item vector$<$ double $>$ $X$: holds inputs for 1 exemplar
\item vector$<$ double $>$ $U0$, $U1$: mean vectors for 1-output
      datasets
\item vector$<$ vector$<$ double $>$ $>$ $U$: a vector of mean vectors
      for multiple output datasets
\item double $P0$, $P1$: a priori probabilities for 1-output datasets
\item vector$<$ double $>$ $P$: a priori probabilities for multiple
      output datasets
\item double $K0$, $K1$: constants for 1-output datasets
\item vector$<$ double $>$ $K$: constants for multiple output datasets
\item double $d0$, $d1$: discriminant function results for 1-output
      datasets
\item vector$<$ double $>$ $d$: discriminant function results for
      multiple output datasets
\item vector$<$ unsigned $>$ $trainClasses$: outputs for multiple
      output training sets
\item vector$<$ unsigned $>$ $testClasses$: outputs for multiple
      output test sets
\end{itemize}

\item {\em Protected} Utility functions

\begin{itemize}
\item \beginred void copy( const DFA\& rhs ): Copy utility
      function, see section \ref{copy_fx} for discussion. \endcolor
\end{itemize}

\end{itemize}

\section{LDFA}

The {\bf LDFA} class is a {\em DFA} object that encodes linear
discriminant function analysis.

\begin{itemize}

\item {\em Public} methods

\begin{itemize}
\item LDFA\& operator= ( const LDFA\& rhs ): see section
      \ref{copy_fx} for discussion
\item LDFA( const LDFA\& rhs ): see section \ref{copy_fx} for
      discussion
\item virtual double train(): Implementation of the pure virtual
      method inherited from {\em Model} that trains a {\em LDFA} {\em
      Model}
\item virtual void reportAccuracy( ostream\& $outputStream$ ):
      Implementation of the pure virtual method inherited from {\em
      Model} that reports the accuracy of the trained {\em LDFA} model
      to an ostream $outputStream$
\end{itemize}

\item {\em Private} members

\begin{itemize}
\item	Matrix$<$ double $>$ C: the common covariance Matrix
\item Matrix$<$ double $>$ S: inverse of the common covariance Matrix
\end{itemize}

\item {\em Private} Utility functions

\begin{itemize}
\item void copy( const LDFA\& rhs ): Copy utility function, see
      section \ref{copy_fx} for discussion.
\end{itemize}

\end{itemize}

\section{QDFA}

The {\bf QDFA} class is a {\em DFA} object that encodes quadratic
discriminant function analysis.

\begin{itemize}

\item {\em Public} methods

\begin{itemize}
\item QDFA\& operator= ( const QDFA\& rhs ): see section
      \ref{copy_fx} for discussion
\item QDFA( const QDFA\& rhs ): see section \ref{copy_fx} for
      discussion
\item virtual double train(): Implementation of the pure virtual
      method inherited from {\em Model} that trains a {\em QDFA} {\em
      Model}
\item virtual void reportAccuracy( ostream\& $outputStream$ ):
      Implementation of the pure virtual method inherited from {\em
      Model} that reports the accuracy of the trained {\em QDFA} model
      to an ostream $outputStream$
\end{itemize}

\item {\em Private} members

\begin{itemize}
\item Matrix$<$ double $>$ $C0$, $C1$: covariance Matrices for
      1-output datasets
\item vector$<$ Matrix$<$ double $>$ $>$ $C$: a vector of covariance
      Matrices for multiple output datasets
\item Matrix$<$ double $>$ $S0$, $S1$: covariance Matrix inverses for
      1-output datasets
\item vector$<$ Matrix$<$ double $>$ $>$ $S$: a vector of covariance
      Matrices for multiple output datasets
\item double $Det0$, $Det1$: covariance Matrix determinants for
      1-output datasets
\item vector$<$ double $>$ $Det$: covariance Matrix determinants for
      multiple output datasets
\end{itemize}

\item {\em Private} Utility functions

\begin{itemize}
\item void copy( const QDFA\& rhs ): Copy utility function, see
      section \ref{copy_fx} for discussion.
\end{itemize}

\end{itemize}

\chapter{Stepwise regression}\label{regresschap}
\setcounter{figcount}{1}

Statistical analysis of neural computational models is based on the
observation that given 2 fully trained neural networks with final
errors $E({\bf y}_{reduced})$ and $E({\bf y}_{full})$ that are
sufficiently close may be discriminated by a chi-square probability
distribution, where

\begin{displaymath}
G^2 = -2N [E({\bf y}_{full}) - E({\bf y}_{reduced})].
\end{displaymath}

with $G^2$ the chi-square probability $p$ with degrees of freedom the
difference in free parameters (connections between nodes) between the
2 neural models, and $N$ the number of exemplars in the dataset (see
section \ref{Wilks}.)  $E$ is expressed as the cross-entropy error
(see equation \ref{xentropyfunk}):
\index{Error function!Cross-entropy}

\begin{equation}
e({\bf t}^k, {\bf a}^{(m)}) =
-{\bf t}^k .* {\bf log} [{\bf a}^{(m)}_k] - ({\bf 1} - {\bf t}^k) .*
{\bf log} [1 - {\bf a}^{(m)}_k].
\end{equation}

In stepwise regression, the 2 networks ${\bf y}_{reduced}$ and
${\bf y}_{full}$ are chosen to differ by 1 input variable (which may
be represented by 1 or more input nodes.)

\section{Stepwise reverse regression}\label{revreg}

Stepwise reverse regression proceeds by training the full network, and
recording that network's final error.  Input variables are removed 1
at a time, the subnetworks so generated are trained to completion and
their errors recorded, and the variable with the largest $p$-value
(least significant) is then recorded with its $p$-value.  Using the
final error of the network which is deficient the one node with the
largest $p$-value as a starting point, input variables are then
removed 2 at a time, with each pair containing the 1 variable with the
largest $p$-value and 1 of the remaining input variables.  The pair of
variables with the largest $p$-value (least significant) is then
recorded with its $p$-value, and the process repeats until a chosen
threshold $p$-value is reached.

\hyperref[revfig]{Figure~\ref{regresschap}.\arabic{figcount}}
demonstrates stepwise reverse regression.

\begin{figure}[htb] \addtocounter{figcount}{1}
\begin{center}
\begin{picture}(4.2,2.2)
\put(0,0){\framebox(4.2,2.2)}
% for dvi
\put(0.1,0.1){\includegraphics[width=4in,height=2in]{reverse.bmp}}
% for dvipdfm
\put(0.1,0.1){\special{pdf: epdf width 4in height 2in (reverse.pdf)}}
\end{picture}
\end{center}
\caption{Stepwise reverse regression}\label{revfig}
\end{figure}

\section{Stepwise forward regression}\label{forreg}

Stepwise forward regression proceeds by training a baseline network
with no input nodes, and recording that network's final error.  Input
variables are added 1 at a time, the subnetworks so generated are
trained to completion and their errors recorded, and the variable with
the smallest $p$-value (most significant) is then recorded with its
$p$-value.  Using the final error of the network which is the one node
with the smallest $p$-value as a starting point, input variables are
then added 2 at a time, with each pair containing the 1 variable with
the smallest $p$-value and 1 of the remaining input variables.  The
pair of variables with the smallest $p$-value (most significant) is
then recorded with its $p$-value, and the process repeats until a
chosen threshold $p$-value is reached.

\hyperref[forfig]{Figure~\ref{regresschap}.\arabic{figcount}}
demonstrates stepwise forward regression.

\begin{figure}[htb] \addtocounter{figcount}{1}
\begin{center}
\begin{picture}(4.2,2.2)
\put(0,0){\framebox(4.2,2.2)}
% for dvi
\put(0.1,0.1){\includegraphics[width=4in,height=2.1in]{forward.bmp}}
% for dvipdfm
\put(0.1,0.1){\special{pdf: epdf width 4in height 2.1in(forward.pdf)}}
\end{picture}
\end{center}
\caption{Stepwise forward regression}\label{forfig}
\end{figure}

\section{The RegressNet class}\label{regressnet}

{\bf RegressNet} is a compound class based on the {\em Network} object
that implements stepwise regression analysis on neural computational
models.

\begin{itemize}

\item {\em Public} methods

\begin{itemize}
\item RegressNet\& operator= ( const RegressNet\& rhs ): see section
      \ref{copy_fx} for discussion
\item RegressNet( const RegressNet\& rhs ): see section \ref{copy_fx}
      for discussion
\item void setInputStructure( const vector$<$ vector$<$ unsigned $>$
      $>$\& $v$): Sets a vector of vector of unsigned $v$ as the
      structure of input nodes representing the variables for the
      stepwise regression. For example, in a 3 variable dataset, if
      variable 0 was represented by input nodes 0 and 1, variable 1 by
      input node 2, and variable 2 by input nodes 3, 4 and 5,
      vector[0] in $v$ would contain (0,1), vector[1] in $v$ would
      contain (2), and vector[3] in $v$ would contain (3,4,5).  All
      input nodes {\em must} be represented in $v$.  See section
      \ref{parsefx} for an utility method that parses a string
      description
      of input variables, and returns a vector of vector of unsigned.
\item void setNetwork( Network* $networkObj$, const double $E$ ): sets
      the {\em Network} object $networkObj$ in the {\em RegressNet}
      object, and the final error $E$ of the trained $networkObj$
\item void reverse\_regress(): performs stepwise reverse regression
\item void forward\_regress(): performs stepwise forward regression
\end{itemize}

\item {\em Private} members

\begin{itemize}
\item Network $*netPtr$: pointer to the incoming Network
\item Network $*netCopyPtr$: pointer to a copy of the incoming Network
      to generate subnetworks
\item bool $historyFlag$: indicates logging to history file
\item double {\em e\_in}: incoming Network's error
\item vector$<$ vector$<$ unsigned $>$ $>$ {\em variable\_defs}: the
      data structure representing input variables
\item typedef multimap$<$ double, unsigned, less$<$ double $>$ $>$
      $ptable$: type definition of a table used to hold stepwise
      regression
      results, p-value first, followed by variable number
\item ostringstream $out$: to simplify writing to screen and history
file
\item strstream $errorOut$: to build an exception message
\item string $errorString$: the exception message
\end{itemize}

\item {\em Private} Utility functions

\begin{itemize}
\item void report( ostringstream\& $o$): a utility function which
outputs
      an ostringstream to screen and history file, takes ostringstream
      $o$
      as its argument
\item void print\_input\_structure( const vector$<$ vector$<$ unsigned
      $>$ $>$ \& $v$): a utility function which prints the input
      variable structure to screen and history file, takes data
      structure representing input variables $v$ as its argument
\item void print\_regression\_table( const ptable\& $table$ ): a
      utility function which prints out a stepwise regression table,
      takes regression table multimap type $ptable table$ as argument
\item bool copy\_network(): a utility function which copies the
      incoming Network object $*netPtr$ into $*netCopyPtr$ so that the
      latter can be made feature deficient for subnetworks. \beginred
      If you are creating a regressable {\em Network} object (and
      chances are that you are,) you {\em must} add code to this
      method.\endcolor
\item string network\_name() const: a utility function which returns
      the name of the type of {\em Network} object.  \beginred
      If you are creating a regressable {\em Network} object (and
      chances are that you are,) you {\em must} add code to this
      method.\endcolor
\item inline double Wilks( const double $N$, const double $E_{full}$,
      const double $E_{sub}$ ) const: a utility function which
      calculates the chi-square parameter for 2 networks with $N$
      exemplars and which differ by their number of input nodes, the
      cross-entropy error of the fully trained network with all input
      nodes being $E_{full}$, and the fully trained feature deficient
      subnetwork's error being $E_{sub}$
      \index{Error function!Cross-entropy}
\item void copy( const RegressNet\& rhs ): copy utility function, see
      section \ref{copy_fx} for discussion.
\end{itemize}

\end{itemize}

\section{Exceptions}

The {\em RegressNet} class throws exceptions of type $RegressNetErr$\&
$e$, where $e$.what() may contain:

\begin{itemize}
\item RegressNet error: maximum node value not $correct$, where
      $correct$ is 1 less than the number of input nodes.  This
      indicates that the vector of vector of unsigned representing the
      variable structure is incorrectly specified.
\item RegressNet error: couldn't copy Network, unknown type.  This
      indicates that the {\em Network} object has not been properly
      coded in RegressNet::copy\_network().
\item RegressNet error: stepwise regression not fully specified.  This
      indicates that {\em Network} object and its final error have not
      been set in the {\em RegressNet} object.
\end{itemize}

\section{Example}

An example driver may thus appear as:

\begin{verbatim}
RegressNet regressionObj; // object to regress a Network

// Say you've already trained a Model pointed to by modelPtr, and
//    it was a Network Model, and the last error was lastError
Network* netPtr = dynamic_cast< Network* >( modelPtr );
regressionObj.setNetwork( netPtr, lastError );

vector< vector< unsigned > > variable_defs; // variable definitions
string variable_string; // for entry of variable definitions

cout << "Please enter a string representing the variables: ";
cin >> variable_string;
try {
   variable_defs = util::variable_parse( variable_string );
}
catch ( util::utilErr& e ) {
   cout << e.what() << endl;
}

if ( !variable_defs.empty() ) // parser returned with a good result
{
   try { // try to set the input structure in the stepwise regression
   object
      regressionObj.setInputStructure( variable_defs );
   }
   catch ( RegressNet::RegressNetErr& e ) {
      cout << e.what() << endl;
   }
}

try {
   regressionObj.reverse_regress();
}
catch ( RegressNet::RegressNetErr& e ) {
   cout << e.what() << endl;
}
\end{verbatim}

\section{Adding code}\label{regresscode}

\beginred If you are creating a regressable {\em Network} object (and
chances are that you are,) you {\em must} add code to the methods
RegressNet::copy\_network() and RegressNet::network\_name().\endcolor
Code added to RegressNet::copy\_network() will take the form:

\begin{verbatim}
else if ( typeid( *netPtr ) == typeid( YourNet ) )
   // Create new object with copy ctor & set pointer to the copy
   netCopyPtr = new YourNet( *dynamic_cast< YourNet* >( netPtr ) );
\end{verbatim}

and code added to RegressNet::network\_name() will take the form:

\begin{verbatim}
else if ( typeid( *netPtr ) == typeid( YourNet ) )
   network_type = "YourNet";
\end{verbatim}

where ``YourNet'' is to be substituted with the name of your {\em
Network} derived object.

Your own {\em Network} derived object must contain code in its copy()
constructor for {\em \bf all} of your derived object member variables
and structures, and in its removeInputs($\ldots$) method for all
member structures that need to be altered to remove input nodes.  See
\hyperref[removeInputs]{section \ref{NetworkMethodsSect}} for a
discussion of removeInputs($\ldots$).

\section{Compiling RegressNet with MSVC++}

Because RegressNet uses Runtime Type Information (RTTI) to determine
the type of {\em Network} object, if you are compiling with MSVC++,
you must set the compiler options accordingly (g++ does this
automatically for you.)  Under the Project menu, choose the Settings
option.  Choose the C/C++ tab, and choose C++ Language in the Category
Box.  Then, click the box for ``Enable Run-Time Type Information
(RTTI)'', as shown in
\hyperref[rttifig]{Figure~\ref{regresschap}.\arabic{figcount}}.

\begin{figure}[htb] \addtocounter{figcount}{1}
\begin{center}
\begin{picture}(4.2,2.6)
\put(0,0){\framebox(4.2,2.6)}
% for dvi
\put(0.1,0.1){\includegraphics[width=4in,height=2.431in]{snap.bmp}}
% for dvipdfm
\put(0.1,0.1){\special{pdf: epdf width 4in height 2.431 in
(snap.pdf)}}
\end{picture}
\end{center}
\caption{Setting RTTI in MSVC++}\label{rttifig}
\end{figure}

\chapter{The Driver}

You of course may design your own driver for the neUROn2++
environment.
Naturally, a command line interface driver is under development.  The
following is the user's guide to that driver.

\setcounter{figcount}{1}

\section{Specify dataset}

This menu choice includes all of the procedures related to specifying
the dataset for the model.  A dataset must be specified {\em before}
any other step.

\subsection{Number of input nodes}

This menu choice allows you to set the number of input nodes.  Its
default is 1 input node, although that should {\em never} be used.

\subsection{Number of output nodes}

This menu choice allows you to set the number of output nodes.  Its
default is 1 output node.

\subsection{Load raw dataset}

This menu choice allows you to load a {\em raw} dataset from a file.
The file must be in the same directory as the program.  A raw dataset
is one whose inputs and outputs have not yet been normalized.
Generally, neural computational procedures require a dataset to be
normalized, whereas statistical procedures do not.  All datasets are
of the format:

\begin{center}
\begin{tabular}{cccccccc}
$i_{11}$ & $i_{12}$ & $\ldots$ & $i_{1I}$ & $o_{11}$ & $o_{12}$ &
$\ldots$ & $o_{1O}$ \\
$i_{21}$ & $i_{22}$ & $\ldots$ & $i_{2I}$ & $o_{21}$ & $o_{22}$ &
$\ldots$ & $o_{2O}$ \\
$i_{31}$ & $i_{32}$ & $\ldots$ & $i_{3I}$ & $o_{31}$ & $o_{32}$ &
$\ldots$ & $o_{3O}$ \\
$\ldots$ & $\ldots$ & $\ldots$ & $\ldots$ & $\ldots$ & $\ldots$ &
$\ldots$ & $\ldots$ \\
$i_{N1}$ & $i_{N2}$ & $\ldots$ & $i_{NI}$ & $o_{N1}$ & $o_{N2}$ &
$\ldots$ & $o_{NO}$ \\
\end{tabular}
\end{center}

Where $i$ is an input to the model, and $o$ is an output, and for $N$
examples (``exemplars'') in the set, with $I$ total inputs and $O$
total outputs.

These columns may be separated in the file by spaces, commas, or any
character which is non-numeric (1 through 9, `.', e, E, - or + are not
allowed as column separators.).

\subsection{Convert a raw dataset into a training set}

This menu choice will convert your raw dataset into a training set, by
normalizing its inputs and outputs to those specified in the menu
choice ``specify variate bounds and types.''  You will be asked if you
want to change the number of thresholds used to calculate ROC area.

You will also be asked if you want to save your training set into
a file.  {\bf This is generally recommended,} as you may want to load
the training set at some point in the future to resume training.
If you do save your training set, you will be asked for the filename
to save its scaling factors.  Scaling factors are used in the equation:
\begin{equation}\label{driverScale}
x_{norm} = lbound + S(x - min)
\end{equation}
where $x_{norm}$ is the normalized result, $lbound$ is the lower limit
of the normalized $x$ value (1 value for all), $S$ is the scaling
factor, and $min$ is the minimum value of each $x$ in the dataset. In
this way, should the model be used in the future, you will be able to
compute the normalized $x_{norm}$ value from the $x$ value that the
user supplies.

\subsection{Randomize a raw dataset into training and test sets}

This menu choice will randomly divide your raw dataset into training
and test sets.  As it will attempt to keep the fraction of outcomes
equal in both sets the same, the number of output nodes {\em must be}
1, and the output {\em must be} discrete (e.g. held to 0 or 1).  You
specify if an output is discrete in the menu choice ``specify variate
bounds and types.''

After the dataset has been randomized, you will be asked if you want
to change the number of thresholds used to calculate ROC area, and if
you want to save the newly created training and test sets to files.

As with converting a raw dataset into a training set, you will be
asked if you want to save your training set and scaling factors into a
file (see \hyperref[driverScale]{equation \ref{driverScale}}).  You
will also be asked if you want to save your test set into a file.
Again, {\bf this is generally recommended.}

You will be given the following choices to specify the number of
exemplars to be placed in the test set, with the remainder to be
placed in the training set (``$N1/N2$'' cross-validation):

\subsubsection{As a whole number}

This choice allows you to specify the number of exemplars in the test
set as a whole number.

\subsubsection{As a fraction with a numerator and denominator}

This choice allows you to specify the number of exemplars in the test
set as a fraction of the number of exemplars in the raw dataset.  You
will be asked for the numerator and the denominator of that fraction.

\subsubsection{As a decimal fraction}

This choice allows you to specify the number of exemplars in the test
set as a fraction of the number of exemplars in the raw dataset.  You
will be asked for the decimal notation of that fraction (e.g. $0.25$
would place $\frac{1}{4}$ of the exemplars from the raw dataset into
the test set, and the remainder into the training set.)

\subsection{Load training set}

This menu choice allows you to load a training set from a file.  The
file must be in the same directory as the program.  After the training
set has been loaded, you will be asked if you want to change the
number of thresholds used to calculate ROC area.

\subsection{Save training set}

This menu choice allows you to save your training set to a file, which
will be saved in the same directory as the program.  You will also be
asked for the filename to save its scaling factors.
See \hyperref[driverScale]{equation \ref{driverScale}} for an
explanation of scaling factors.

\subsection{Save scaling factors}

This menu choice allows you to save the scaling factors of the
training set, so that should the model be used in the future, you will
be able to compute the normalized $x_{norm}$ values from the $x$
values that the user supplies.  See \hyperref[driverScale]{equation
\ref{driverScale}} for an explanation of scaling factors.

\subsection{Load test set}

This menu choice allows you to load a test set from a file.  The file
must be in the same directory as the program.  After the test set has
been loaded, you will be asked if you want to change the number of
thresholds used to calculate ROC area.

\subsection{Save test set}

This menu choice allows you to save your test set to a file, which
will be saved in the same directory as the program.

\subsection{Log dataset operations to history file}

Set this option to ``yes'' if you want the operations you perform in
the ``Specify dataset'' menu to be logged to the history file {\bf
neuron.log}.  This file will appear in the same directory as the
program.

\subsection{Specify variate bounds and types}

In the following menu choices, variates are normalized according to
the equation:

\begin{equation}\label{USERinputScale}
x_{norm} = lbound + [({{x - min} \over {max - min}}) \times (hbound -
lbound)]
\end{equation}

Where $lbound$ is the lower boundary of a normalized value, $hbound$
is the upper boundary, $min$ is the minimum value of a variable {\em
in the training set}, $max$ is the maximum value of a variable {\em in
the training set}, $x$ is the variable's value and the normalized
value of $x$ is $x_{norm}$.

\subsubsection{Output is discrete}

Set this to ``yes'' if your outputs will always be either 0 or 1.

\subsubsection{Threshold for discrete output}

Set this to the number over which model outputs will be defined as 1,
and under which model outputs will be defined as 0 for discrete
output.  For example, if the threshold is $0.5$, a numeric model
output of $0.6$ will be transformed to 1, and a numeric model ouput of
$0.3$ will be transformed to 0.  The default is $0.5$.

\subsubsection{Lower limit of input variates}

Choose this menu item to set $lbound$ for inputs in equation
\ref{USERinputScale}.  The default is $-0.9$.

\subsubsection{Upper limit of input variates}

Choose this menu item to set $hbound$ for inputs in equation
\ref{USERinputScale}.  The default is $+0.9$.

\subsubsection{Lower limit of output variates}

Choose this menu item to set $lbound$ for outputs in equation
\ref{USERinputScale}.  The default is $+0.1$ if the output is not discrete
(and obviously 0 if it is).

\subsubsection{Upper limit of output variates}

Choose this menu item to set $hbound$ for outputs in equation
\ref{USERinputScale}.  The default is $+0.9$ if the output is not discrete
(and obviously 1 if it is).

\subsection{Number of thresholds to calculate ROC area}

Choose this menu item to set the number of thresholds over which ROC
area is calculated.  The output {\em must be} discrete, and the number
of outputs {\em must be} 1 for ROC area to be calculated.

\section{Specify model}

Choose this menu item to specify a {\em neural computational} model.
A dataset must be specified {\em before} this choice, and this menu
choice must obviously be specified before ``Use model''.  Discriminant
function analysis {\em does not} require this menu choice, and in
fact, it is useless for discriminant function analysis.

\subsection{Bias nodes}

Choose this menu item to specify whether or not bias nodes will be
included in your neural computational model.  In general, for a
real-world model, you should {\em always} include bias nodes.  The
default is ``yes''.

\subsection{Number of hidden layers / nodes on each layer}

Choose this option to specify the number of hidden layers and number
of nodes on each layer.  The default is 1 layer with 1 node, although
neural computational models will generally have $>$ 1 hidden node.
Most models will, however, only have 1 hidden layer.

\subsection{Output error function}

Choose this option to specify the output error function.  ``LMS'' is
least-mean-squared:

\begin{equation}
E = \frac{1}{2} \sum_{i=1}^{m} (y_i - o_i)^2
\end{equation}
where $E$ is the error, $m$ the number of outputs, $y$ is the known
value and $o$ the model's output.

You may also choose the cross-entropy error function ``X-entropy'':

\begin{equation}
E = \sum_{i=1}^{m} -y\log(o)-(1-y)\log(1-o)
\end{equation}

For the cross-entropy error function, the output {\em must be}
discrete.  Note that for sigmoidal output nodes, where for
output $\vec{\bf{o}}$,

\begin{displaymath}
\vec{\bf{o}} = \frac{1}{1+e^{-\vec{\bf{x}}}}
\end{displaymath}

if any element of $\vec{\bf{x}}$ is very large or very small,
the program will round to $0$ or $1$, and the $\log$ terms in the
cross-entropy error function will cause an error.  In such cases,
an approximation is made (see the document {\bf manifest.pdf}
for a detailed explanation of that approximation), and the following
error will be generated:

\begin{verbatim}
WARNING: Numerical out of bounds encountered when calculating error
\end{verbatim}

\subsection{Load network from a file}

This menu choice allows you to load a neural computational model from
a file.  The file must be in the same directory as the program.

\subsection{Binary logistic regression}

This menu choice specifies binary logistic regression, which is
mathematically equivalent to a single output node neural network with
a bias, the sigmoidal transfer function, no hidden nodes, and the
cross-entropy error function. {\em \beginred Always run a binary
logistic regression model} {\bf before} {\em starting to build a
neural network. \endcolor}  If your binary logistic model has
acceptable goodness-of-fit, you don't need to build a neural network.

Note that at the end of a binary logistic run, $\beta$ weights will be
reported, along with their standard errors and $p$-values for the Wald
test that the $\beta$ weights are nonzero.\footnote{see Hosmer and
Lemeshow, Applied Logistic Regression, 2nd edition, John Wiley \&
Sons, Inc., New York, 2000, pp 36--42}  Note also that the indexing
begins with input number 0, which is often referred to as the $y$-
intercept.  {\em \beginred In this case, $\beta_0$ is} {\bf not} {\em
the $y$-intercept. \endcolor}  It is the $\beta$ weight of the first
variate, indexed to 0.  The $y$-intercept is clearly labeled as
``$y$-intercept'' in the report table.

\subsection{Log last operation to file}

Set this option to ``yes'' (the default) if you want the last neural
computational operation performed to be logged to the file {\bf
model.txt}.  This file will be {\em overwritten} with each run, and
will appear in the same directory as the program.

\subsection{Log to history file}

Set this option to ``yes'' (the default) if you want the neural
computational operations that are performed logged to the file {\bf
neuron.log}.  This file will be {\em appended} with each operation,
and will appear in the same directory as the program.

\section{Use model}

Choose this menu item to use a specified {\em neural computational}
model.  A dataset must be specified {\em before} this choice, and the
``Specify Model'' menu must also have been entered before ``Use
model''.

\subsection{Randomize weights}

Choose this option to randomize the weights of the neural
computational model.  Prior to training the model, weights {\em must
be} set, either by randomization or by loading the model from a file
in the ``Specify model'' menu.  By default, weights are randomized
between $-0.5$ and $+0.5$.

\subsection{Set learning rate}

Choose this option to set the learning rate $\eta$ of the neural
computational model.  The default is $\eta = 1/N$ for on-line
backpropagation (where $N$ is the number of exemplars), and $\eta = 1$
for off-line (``batch'' or ``epoch'') training.

You may choose the learning rate to be ``automatic'' (see the Manifest
documentation for a fuller explanation of automatic stepsize
selection.) Briefly, automatic stepsize selection is designed to
prevent oscillation around a local error minimum.  Due to the increase
in time required for this feature, it may be most optimal to train to
near a local error minimum before setting automatic stepsize selection
{\em on}.

\subsection{Set stopping conditions}

Choose this option to set conditions by which the neural computational
model will stop training.  These include:

\subsubsection{Set maximum number of iterations}

The model will stop training after the specified number of iterations.
The default is 1000000, {\em which is arbitrary.  Do not consider any
neural computational model to be trained by stopping at any arbitrary
threshold.  You} must {\em inspect the model yourself to see if it has
stopped training.}

\subsubsection{Lower error limit}

Use this choice to set the lower error limit that, if the error of the
model decreases below, will result in cessation of training.  For
example, you could set this to the limit of your compiler.  By
default, no lower error limit is set.

\subsubsection{Limit of change in error over 1 iteration}

Use this choice to set the lower error limit that, if the change in
error over 1 iteration is below, will result in cessation of training.
Thus, in

\begin{equation}
\Delta E = E_{t+1} - E_t
\end{equation}

Where $E_t$ is the error after a training iteration, and $E_{t+1}$ is
the error after the next training iteration, if $\Delta E$ is less
than the specified value, training stops.  By default, this condition
is not specified.

\subsubsection{Stop if increase over error window width}

Allows you to set a number of iterations that after which the error
increases, training will stop.  Thus, in

\begin{equation}
E_{n+\Delta} > E_{n}
\end{equation}

Where $n$ is the current iteration, $E$ the error at an iteration, and
$\Delta$ the error window width, training will stop if the inequality
is satisfied.  By default, this condition is not specified.

\subsubsection{Maximum absolute gradient limit}

All neural networks update their weights in the following form:

\begin{equation}
{\bf w}_{t+1} = {\bf w}_t - \eta {\bf g}_t
\end{equation}

where $\bf{w}_t$ is the weight structure at time $t$, $\eta$ is the
learning rate, and $\bf{g}_t$ is the gradient.  This menu choice
allows you to stop training when the maximum absolute value of the
elements in the gradient structure $\bf{g}_t$ decreases below the
chosen limit.  The default is $10^{-6}$.

\vspace{9 pt}{\bf \large \beginred READ THIS! \endcolor}

\beginred The maximum absolute gradient limit should be the primary
means by which you judge that training has ended, {\em not} by
the error minimum. \endcolor

For canonical backpropagation, using the maximum absolute gradient
limit will significantly slow training (not by the number of
iterations it takes to train, but by the amount of computational time
consumed per iteration.)  Thus, one strategy is to turn the maximum
gradient choice off until you are near a local error minimum as judged
loosely by the decrease in the error, then to turn on this menu
choice.

Don't be surprised if the maximum absolute gradient fluctuates
(increases and decreases.)  Remember, there is more than one node and
one layer in even the simplest neural network.  Keep an eye on the
overall trend of the maximum absolute gradient to insure it is
decreasing.

\subsection{Batch Epoch}

Choose this option to set batch/epoch learning.

The {\em on-line} backpropagation learning algorithm is given by the
formula:
\begin{equation}
{\bf w}_{t+1} = {\bf w}_t - \eta {\bf g}_k({\bf w}_t)
\end{equation}
where $k$ is the exemplar index, ${\bf w}$ is the weight, $t$ the
iteration, ${\bf g}$ is the error gradient with respect to the weight,
and $\eta$ is the learning rate.  By default, if on-line learning is
chosen, $\eta = 1/N$, where $N$ is the number of exemplars.

The {\em off-line}, also known as ``batch'' or ``epoch''
backpropagation learning algorithm is given by the formula:
\begin{equation}
{\bf w}_{t+1} = {\bf w}_t - \eta {\bf g}({\bf w}_t)
\end{equation}
where the {\em average gradient} ${\bf g}$ is defined by the formula:
\begin{equation}
{\bf g} = (1/N) \sum_{k=1}^N {\bf g}_k.
\end{equation}
By default, if off-line learning is chosen, $\eta = 1$.

See the neUROn2++ Design Manifest for a more detailed description of
on-line and off- line learning.

If batch/epoch is set {\em on}, off-line, ``batch'' or ``epoch''
learning is used. If batch/epoch is set {\em off}, on-line learning is
used.

\subsection{Weight Decay}

The {\em error function}, $E_k$, for the $k$th exemplar is
defined by the formula:
\begin{equation}\label{USERdeferrorfunk}
E_k({\bf w}) =
e(y^k, {\bf o}^{(m)}_k) +
\frac{1}{2\sigma_w^2} \sum_{i=1}^m |{\bf w}|^2.
\end{equation}

where $y$ is the known or target outcome, and ${\bf o}$ is the model's
output.  The second term on the right hand side of
(\ref{USERdeferrorfunk}) is required to guarantee reliable statistical
inference. It basically `` stabilizes'' the learning process by
preventing the weights in the network from becoming too large or too
small.  $\sigma_w$ is the standard deviation of the weights {\em \bf
chosen by the user.}  Thus, a connection weight is expected to lie
between $-1.96 \sigma_w$ and $+1.96 \sigma_w$ with probability equal
to about 95\%. Adjustment of this parameter can ``change results'' of
the statistical analysis so $\sigma_w$ is constrained to be exponents
of 10 (e.g. 10, 100, 1000, $\ldots$ $10^{12}$.  Once chosen,
$\sigma_w$ should not be changed.

A reasonable choice for $\sigma_w$ might be $\sigma_w = 100$. This
basically constrains the ``learning process'' to find connection
strengths for the ANN such that the connection strength weight ranges
between $-200$ and $+200$. The choice of $\sigma_w$ should be reported
in conjunction with all statistical analysis results. Note that as
$\sigma_w \rightarrow \infty$, the effect of the summation term
becomes negligible.

The learning rate equivalent of equation \ref{USERdeferrorfunk} is

% Note the necessity of \left. even though it doesn't show
\begin{eqnarray*}
\vec w_{t+1} & = & \vec w_t - \eta \left. \D{E'}w \right|_{w_t} \\
& = & \vec w_t - \eta \left. \D{E}w \right|_{w_t} - 2\eta\lambda\vec w_t \\
& = & (1-2\eta\lambda)\vec w_t - \eta \left. \D{E}w \right|_{w_t}
\end{eqnarray*}

where $(1-2\eta\lambda)$ is ``weight decay''. $\lambda$ may
be set by either of 2 methods:

\begin{itemize}

\item The complexity of the population objective function $>>$ sample
objective function:

\begin{equation}
\lambda = \frac 1{N2\sigma_w^2}
\end{equation}

\item The complexity of the population objective function $\simeq$
sample objective function:

\begin{equation}\label{USERsameComplexEqn}
\lambda = \frac 1{2\sigma_w^2}
\end{equation}

\end{itemize}

In general, as we don't know the complexity of the population
objective function a priori, we choose equation \ref{USERsameComplexEqn}
to calculate $\lambda$.

The following choices allow you to specify weight decay:

\subsubsection{Don't change anything}

Allows the uninformed user (which, assuming you've read this section,
you are {\em not},) to exit without changing the way in which weight
decay is calculated.

\subsubsection{lambda = 1 / ( 2 * sigma$^\wedge$2 )}

Specifies that the weight decay term be calculated by

\begin{displaymath}
\lambda = \frac 1{2\sigma_w^2}
\end{displaymath}

and sets $\sigma_w$ to be an exponent of 10 from $10^1$ to $10^{12}$.

\subsubsection{lambda = 1 / ( 2 * N * sigma$^\wedge$2 )}

Specifies that the weight decay term be calculated by

\begin{displaymath}
\lambda = \frac 1{N2\sigma_w^2}
\end{displaymath}

and sets $\sigma_w$ to be an exponent of 10 from $10^1$ to $10^{12}$.

\subsubsection{No weight decay}

Turns off weight decay.

\vspace{9 pt}{\bf \large \beginred READ THIS! \endcolor}

Using too large of a weight decay term will degrade network modeling
performance, but it is critical to use weight decay in order for
statistical inference, for example, stepwise regression.  Thus,
setting a proper weight decay value is of utmost importance.  Here is
a good strategy to determine the proper weight decay value:

{\em (Note that the default value for the weight decay term is
$\frac 1{2\sigma_w^2}$ with $\sigma_w = 100$.  This is to prevent the
uninformed user from making incorrect statistical inferences.
However, you are no longer the uninformed user. In the first step, you
will turn off weight decay.)}

\begin{enumerate}

\item Turn weight decay off.

\item Experiment with different neural network architectures until you
find one that is as near as possible a strict local minimum, and with
as high ROC area as possible.

\item Save the neural network from Step 2.

\item Load the neural network from Step 2, and turn on weight decay.
Use the choice for $\lambda = \frac 1{2\sigma_w^2}$, and choose a low
initial value for $\sigma$, for example $10^2$.

\item Train the neural network to as near a possible a strict local
minimum (you'll be starting with the final weights from Step 2.)

\item Compare the ROC area to that from Step 2.  If it is similar, go
to step 8.

\item If the ROC area has substantially degraded, increase $\sigma$,
load the neural network from Step 2, and go to Step 5.

\item Save the neural network, and record $\sigma_w$, as you will cite
it when you publish your results. Perform stepwise regression with
this value in the weight decay term, beginning with the final weights
from the network you saved in this Step.

\end{enumerate}

\subsection{Print counter}

This menu choice allows you to indicate whether printing will be
linear (prints after a set number of iterations) or logarithmic.
Logarithmic is recommended.

\subsection{Train model}

This menu choice initiates neural computational model training.  You
will choose the type of gradient descent.  For an explanation of
conjugate gradient descent and the Shanno algorithm, see Golden, pp.
217--218 and 221--222.\footnote{Richard M. Golden, Mathematical
Methods for Neural Network Analysis and Design, 1996, MIT Press, ISBN
0-262-07174-6, pp. 217--218, 221--222}  A brief explanation of these
methods may also be found in the design manifest.

The menu choices are ordered by robustness of the algorithm, from less
to more.  Thus, conjugate gradient descent is less robust than the
Shanno algorithm.

\vspace{9 pt}{\bf \large \beginred READ THIS! \endcolor}

Conjugate gradient descent and the Shanno algorithm generally are much
faster at finding a local error minimum than canonical backpropagation
{\em once one is near the local error minimum}.  Thus, your training
strategy should be to use canonical backpropagation until you are near
a local error minimum (as judged by the maximum absolute gradient),
and then apply either conjugated gradient descent or the Shanno
algorithm. Empirically try each advanced algorithm, and choose the one
to use based on your trials.

\subsubsection{Canonical backpropagation}

This menu choice implements standard, plain vanilla backpropagation.

\subsubsection{Conjugate gradient descent}

This menu choice implements conjugate gradient descent.  See the
design manifest and Golden pp. 221--222.

\subsubsection{Shanno algorithm}

This menu choice implements the Shanno algorithm.  See the design
manifest and Golden pp. 217--218.

\subsection{Save the network to a file}

This menu choice allows you to save your trained neural computational
model with its weights to a file.  You may train the model repeatedly
prior to saving it, if you so choose.  If you have one output, and it
is discrete, you will also be asked if you want to save the network's
guesses to a file. These guesses will be saved as a 2 column file,
with the first column being the known value, and the second column
being the network's guess. You will be asked if you want to save
guesses for the training set, and if it is loaded, the test set.

\subsection{Stepwise regression}\label{stepregress}

Stepwise regression allows you to determine which input variables are
statistically significant to your trained neural computational model.
For a more complete discussion, see the neUROn2++ Design Manifest.
If you have not just finished training a neural network, you will be
asked for the final error of the network model to regress.

\subsubsection{Specify variables}

A variable may be represented by more than 1 input node.  For example,
ethnicity may be represented by 4 nodes: African-American, Asian,
Hispanic and Caucasian.  Or a variable may be represented by 2 nodes:
1 in which the value is present or absent, and the other by the value
of the variable itself.  Variable structure may be input from the
command line or loaded from file.

Input characters include:
\begin{itemize}
\item {\bf ;} delimits variables
\item {\bf ,} delimits nodes
\item {\bf -} specifies a range of nodes
\item spaces are ignored
\end{itemize}

For example, ``0-3,5; 4,8; 6,7'' indicates that variable 0 is composed
of input nodes 0,1,2,3 and 5, variable 1 is composed of input nodes 4
and 8, and variable 2 is composed of input nodes 6 and 7.

\subsubsection{Stepwise reverse regression}

Stepwise reverse regression proceeds by training the full network, and
recording that network's final error.  Input variables are removed 1
at a time, the subnetworks so generated are trained to completion and
their errors recorded, and the variable with the largest $p$-value
(least significant) is then recorded with its $p$-value.  Using the
final error of the network which is deficient the one node with the
largest $p$-value as a starting point, input variables are then
removed 2 at a time, with each pair containing the 1 variable with the
largest $p$-value and 1 of the remaining input variables.  The pair of
variables with the largest $p$-value (least significant) is then
recorded with its $p$-value, and the process repeats until a chosen
threshold $p$-value is reached.

Figure~\arabic{figcount} demonstrates stepwise reverse regression.

\begin{figure}[htb]\addtocounter{figcount}{1}
\begin{picture}(4.2,2.2)
\put(0,0){\framebox(4.2,2.2)}
% for dvi
\put(0.1,0.1){\includegraphics[width=4in,height=2in]{reverse.bmp}}
% for dvipdfm
\put(0.1,0.1){\special{pdf: epdf width 4in height 2in (reverse.pdf)}}
\end{picture}
\makebox[4in]{\rule{0in}{2in}}
\caption{Stepwise reverse regression}
\end{figure}

\subsubsection{Stepwise forward regression}

Stepwise forward regression proceeds by training a baseline network
with no input nodes, and recording that network's final error.  Input
variables are added 1 at a time, the subnetworks so generated are
trained to completion and their errors recorded, and the variable with
the smallest $p$-value (most significant) is then recorded with its
$p$-value.  Using the final error of the network which is the one node
with the smallest $p$-value as a starting point, input variables are
then added 2 at a time, with each pair containing the 1 variable with
the smallest $p$-value and 1 of the remaining input variables.  The
pair of variables with the smallest $p$-value (most significant) is
then recorded with its $p$-value, and the process repeats until a
chosen threshold $p$-value is reached.

Figure~\arabic{figcount} demonstrates stepwise forward regression.

\begin{figure}[htb]\addtocounter{figcount}{1}
\begin{picture}(4.2,2.2)
\put(0,0){\framebox(4.2,2.2)}
% for dvi
\put(0.1,0.1){\includegraphics[width=4in,height=2.1in]{forward.bmp}}
% for dvipdfm
\put(0.1,0.1){\special{pdf: epdf width 4in height 2.1in(forward.pdf)}}
\end{picture}
\makebox[4in]{\rule{0in}{2.1in}}
\caption{Stepwise forward regression}
\end{figure}

\section{Discriminant function analysis}

Choose this menu item to use discriminant function analysis to model
your dataset(s).

\subsection{Linear discriminant function analysis }

Choose this menu item to model your dataset(s) with {\em LDFA}, linear
discriminant function analysis.

\subsection{Quadratic discriminant function analysis }

Choose this menu item to model your dataset(s) with {\em QDFA},
quadratic discriminant function analysis.

\subsection{Log last operation to file}

Set this option to ``yes'' (the default) if you want the last
statistical operation performed to be logged to the file {\bf
model.txt}.  This file will be {\em overwritten} with each modeling
procedure, and will appear in the same directory as the program.

\subsection{Log to history file}

Set this option to ``yes'' (the default) if you want the statistical
operations that are performed logged to the file {\bf neuron.log}.
This file will be {\em appended} with each operation, and will appear
in the same directory as the program.


\chapter{Utility Scripts}\label{utility_chap}
\setcounter{figcount}{1}

\section{commenttex.pl}\index{Utility scripts!Convert LaTeX in a C++ file}

It is very useful to include \LaTeX mathematical equations in the
documentation of code, e.g.:

\begin{verbatim}
// Calculate the error term for the output unit
// (Freeman & Skapura p. 102, #6)
// $\delta^o_{pk} = (y_{pk} - o_{pk}){f^o_k}'(net^o_{pk})$ for mean squared error
// Note that the sign is changed ( o - y ) instead of ( y - o ) to conform to the
//    Methodology section
// Methodology equation 2.8 (t=y, a=o)
// $\delta^{(m)}_k = -({\bf t}_k - {\bf a}^{(m)})$ for mean squared error,
// or equation 2.9
//
$\delta^{(m)}_k = -({\bf t}_k - {\bf a}^{(m)}) ./ [{\bf a}^{(m)} .* ({\bf 1} - {\bf a}^{(m)}]$
//    for x-entropy error,
//    multiplied by the ${\bf Df}^{(j)}_k$ term in equation 2.7
// Note that if the error is x-entropy, then dE/do = (y-o)/(o(1-o))
//    and the o(1-o) in the denominator cancels d_sigmoidal
//   ($={f^o_k}'$, $={\bf Df}^{(j)}_k$) = o(1-o)
//    hence splitting this code into 2 lines with an if:
o_err = o - y[ example ];
if ( errorType == 0 ) // error is LMS
   o_err *= d_sigmoidal()( o );
\end{verbatim}

The Perl script {\em commenttex.pl} will create a \LaTeX document with
the \LaTeX equations identified by \$$\ldots$\$ output in normal
\LaTeX code, and everything else ``verbatim''.  Once that \LaTeX
document is then processed, you can easily `see' the mathematical
formulae.

As an example, suppose you wanted to `see' the above code from
bareprop.cpp. You would enter the following lines (assuming you have a
working \LaTeX installation on your computer):

\begin{verbatim}
commenttex.pl bareprop.cpp > bareprop.tex
latex bareprop
dvipdfm bareprop
\end{verbatim}

(You would of course not enter the `\$s', as they are the command
prompts.) The first line creates the \LaTeX file, the second line runs
the \LaTeX file, and the final line creates an Adobe pdf file (and is
not strictly necessary, but certainly convenient.)  The code would
then
appear as:

\vspace{3 pt}
$\ldots$
\begin{verbatim}
      // Calculate the error term for the output unit
      // (Freeman & Skapura p. 102, #6)
\end{verbatim}
\hspace{30 pt}//\hspace{4 pt}
$\delta^o_{pk} = (y_{pk} - o_{pk}){f^o_k}'(net^o_{pk})$ for mean squared error
\begin{verbatim}
      // Note that the sign is changed ( o - y ) instead of ( y - o )
      to conform to the
      //    Methodology section
      // Methodology equation 2.8 (t=y, a=o)
\end{verbatim}
\hspace{30 pt}//\hspace{4 pt}
$\delta^{(m)}_k = -({\bf t}_k - {\bf a}^{(m)})$ for mean squared error,
\begin{verbatim}
      // or equation 2.9
\end{verbatim}
\hspace{30 pt}//\hspace{4 pt}
$\delta^{(m)}_k = -({\bf t}_k - {\bf a}^{(m)}) ./ [{\bf a}^{(m)} .* ({\bf 1} - {\bf a}^{(m)}]$
\begin{verbatim}
      //    for x-entropy error,
\end{verbatim}
\hspace{30 pt}//\hspace{20 pt}multiplied by the ${\bf Df}^{(j)}_k$
term in equation 2.7
\begin{verbatim}
      // Note that if the error is x-entropy, then dE/do = (y-o)/(o(1- o))
      //    and the o(1-o) in the denominator cancels d_sigmoidal
\end{verbatim}
\hspace{30 pt}//\hspace{15 pt}($={f^o_k}'$, $={\bf Df}^{(j)}_k$) = o(1-o)
\begin{verbatim}
      //    hence splitting this code into 2 lines with an if:
      o_err = o - y[ example ];
      if ( errorType == 0 ) // error is LMS
         o_err *= d_sigmoidal()( o );
\end{verbatim}

\section{mkdataset.pl}\index{Utility scripts!Convert Excel file to dataset}

Most often, datasets begin with a Microsoft Excel file.  Often, these
files contain blank spaces in the columns, and it is your job to
recode the file so that no blank spots remain.  The Perl script
{\em mkdataset.pl} was written to help you do just that.  Columns with
blank entries will be replaced by 0, 0 where the entry is blank,
and 1, {\em data value} where the entry is populated.  Fully populated
columns are simply printed.  For example, the following file:

\begin{verbatim}
Age, Gender, Voting preference
32, 1, 1
58, , 0
63, 0, 0
37, 0,
60, 0, 1
\end{verbatim}

Would be converted to:

\begin{verbatim}
32, 1, 1, 1
58, 0, 0, 0
63, 1, 0, 0
60, 1, 0, 1
\end{verbatim}

\hyperref[noelim]{See below} for an explanation of why the line

\begin{verbatim}
37, 0,
\end{verbatim}

was eliminated.

Obviously, you'll still need to first transform the columns from an
Excel file into their numerical counterparts (e.g. if M and F are in
a column, you would need to change them to 0 and 1.)  But this Perl
script will take you much of the way towards a dataset that you can
then load into neUROn2++.

This Perl script expects a comma delimited file, so first save your
Excel spreadsheet in the .csv format (choose File, Save As, and then
in the ``Save as type'' box, choose ``CSV (Comma delimited)''.)

Then run the Perl script as follows:
\begin{verbatim}
mkdataset.pl [options] [FILE]
\end{verbatim}

Note that the filename is passed as the last argument.

Options include:

\begin{center}
\begin{tabular}{ll}
-h, -help or ?: & Prints this list \\
-quiet: & Suppress messages \\
-noheader: & No first line of labels will be processed \\
-key [FILE]: & Create a key file from the label line \\
-noelim: & Do not eliminate rows with the last column blank \\
-inputs [FILE]: & Create a file with the input node data structure \\
-o [FILE]: & Output file \\
\end{tabular}
\end{center}

\subsection{Notes:}

\subsubsection{-noheader}

Typically, an Excel file will have a first row that consists of the
labels for the columns.  If this option is specified, no such first
row will be processed.

\subsubsection{-key [FILE]}

If a first row is present that consists of the labels for the columns,
you may specify a file that will record those labels, prettily printed
with column numbers.

\subsubsection{-noelim}\label{noelim}

Typically the last column of an Excel spreadsheet represents the
outcome variable (note this is not the case with multiple output
datasets.)  Also typically you'll want to eliminate those rows that
have this last column blank. If this option is specified, rows with
blank last columns will {\em not} be eliminated.

\subsubsection{-inputs [FILE]}

Writes the input node data structure to file suitable for input
to neUROn2++ stepwise regression.

\subsection{Typical usage:}

A typical incarnation of {\em mkdataset.pl} for an Excel comma
delimited file {\em myfile.csv} with the first row containing the
labels of the columns would be

\begin{verbatim}
mkdataset.pl -key key.txt -o output.csv myfile.csv
\end{verbatim}

which would write the key file to {\em key.txt} and the output file to
{\em output.csv}.

\section{neuron2html.pl}\index{Utility scripts!Generate html/Javascript file}\label{neuron2html.pl}

Created by Ashay Kparker for SimpleProp networks, and modified by
Gaurav Bansal for BackProp and Binary Logistic Regression models,
{\em neuron2html.pl} is one of the most powerful Perl scripts for the
neUROn2++ user.

Once you've fully trained your model, you'll want to make it usable by
others.  The World Wide Web offers a way to widely distribute your
trained neural network in a user-friendly manner.  Running
{\em neuron2html.pl} (assuming, of course, you have Perl installed,) you'll be asked to
enter:
\begin{itemize}
\item The name of the file with the scaling factors that you saved
      from running part 1 of the driver, ``Specify dataset'' (now you
      know why it's so important to save those scaling factors!)
\item The name of the file that you saved with the weights of your
      fully trained model from running part 3 of the driver, ``Use
      model'' (now you know why you should save every trained model!)
\item A specification file that you create that will tell the Perl
      program what to use as the names of your variables
\item The name of the output html file (be sure to name it with the
      extension .htm or .html!)
\end{itemize}

\subsection{The specification file}

You'll need to make a file that will tell the Perl program what type
of variables you have, and what to name them.  This also includes your
outcome variable.  Each line in the file refers to a variable.  In
general, the first letter (or 2) of each line refers to the kind of
variable you have, and the remaining entries on that line tell the
program how to name that variable.  Entries on a line are separated by
the delimeter {\bf \%\%}.

\subsubsection{Numerical variables}

Numerical variables are specified in the manner:

N \%\% Numerical Field Name \%\% Numerical Field Units

$ $ % make a blank line here

Example:
\begin{verbatim}
N %% Age %% years
\end{verbatim}

\subsubsection{Categorical variables}

Categorical variables are specified in the manner:

C \%\% Categorical Field Name \%\% Choice 1 \%\% $\ldots$ \%\% Choice
n

$ $ % make a blank line here

Example:
\begin{verbatim}
C %% Ethnicity %% African American %% Hispanic %% Caucasian %% Asian
\end{verbatim}

\subsubsection{Binary variables}

Binary variables are specified in the manner:

B \%\% Boolean Field Name \%\% Unit True \%\% Unit False

$ $ % make a blank line here

Example:
\begin{verbatim}
B %% Digital rectal exam %% Abnormal %% Normal
\end{verbatim}

\subsubsection{Variables whose values may be unknown}\label{k_tag}

Any of the tags above may be preceded by a K tag to specify if the
variable following it is known or unknown in a not fully populated
dataset.  This addresses the coding of unknown variables in a not
fully populated dataset by adding an additional binary variable to
indicate if the succeeding variable is present or absent, e.g. if the
first variable is present or absent, and the second variable is length
in inches: \\
1 3 indicates length known, and length = 3 inches \\
1 0 indicates length known, and length = 0 inches \\
0 0 indicates length unknown \\

For known/unknown numerical variables:

K \%\% N \%\% Numerical Field Name \%\% Numerical Field Units

$ $ % make a blank line here

Example:
\begin{verbatim}
K %% N %% Age %% years
\end{verbatim}

$ $ % make a blank line here

For known/unknown categorical variables:

K \%\% C \%\% Categorical Field Name \%\% Choice 1 \%\% $\ldots$ \%\%
Choice n

$ $ % make a blank line here

Example:
\begin{verbatim}
K %% C %% Ethnicity %% African American %% Hispanic %% Caucasian %%
Asian
\end{verbatim}

$ $ % make a blank line here

For known/unknown binary variables:

K \%\% B \%\% Boolean Field Name \%\% Unit True \%\% Unit False

$ $ % make a blank line here

Example:
\begin{verbatim}
K %% B %% Digital rectal exam %% Abnormal %% Normal
\end{verbatim}

\subsubsection{Reporting the outcome}

After the user presses the ``predict'' button, a window will pop up
with the model's predicted outcome.  The `O' and `R' tags specify what
that window will say.

$ $ % make a blank line here

To specify what the discrete output means, use the `O' tag:

O \%\% 0 Response \%\% 1 response

$ $ % make a blank line here

Example:
\begin{verbatim}
O %% Cancer absent %% Cancer present
\end{verbatim}

$ $ % make a blank line here

To specify how to label the odds for the prediction, use the `R' tag:

R \%\% Odds ratio label

$ $ % make a blank line here

Example:
\begin{verbatim}
R %% Odds of cancer on prostate biopsy
\end{verbatim}

\subsection{Example specification file}

Suppose that you had trained a model with a dataset that had age,
gender and ethnicity as input variables, and voting preference
(Democrat or Republican) as an outcome variate.  For simplicity, we'll
say that the dataset was fully populated, so you didn't need to encode
for missing variates.  Your $1^{st}$ column in the dataset was age,
encoded in years old, and your $2^{nd}$ column was gender, with female
encoded as `0', and male encoded as `1'. Your $3^{rd}$, $4^{th}$,
$5^{th}$ and $6^{th}$ columns were encoded with ethnicity, with a `1'
in the $3^{rd}$ column representing African-American, a `1' in the
$4^{th}$ column representing Hispanic, a `1' in the $5^{th}$ column
representing Caucasian, and a `1' in the $6^{th}$ column representing
Asian.  (All of the columns beside the column marked `1' in the
ethnicity columns would be marked `0'.)  In the $7^{th}$ column of
your dataset was voting preference, with `0' Democrat, and `1'
Republican.

Your specification file might then look like:

\begin{verbatim}
N %% Age %% years
B %% Gender %% Male %% Female
C %% Ethnicity %% African American %% Hispanic %% Caucasian %% Asian
O %% Democrat %% Republican
R %% Odds of voting Republican
\end{verbatim}

The generated html file and its pop-up window when the user pressed
the ``predict'' button would then appear as in
\hyperref[html_out]{Figure~\ref{utility_chap}.\arabic{figcount}}.

\begin{figure}[htb] \addtocounter{figcount}{1}
\begin{center}
\begin{picture}(4.2,4.4)
\put(0,0){\framebox(4.2,4.4)}
% for dvi
\put(0.1,1.3){\includegraphics[width=4in,height=2.971in]{html.bmp}}
\put(0.1,0.1){\includegraphics[width=4in,height=1.067in]{htmlout.bmp}}
% for dvipdfm
\put(0.1,1.3){\special{pdf: epdf width 4in height 2.971in(html.pdf)}}
\put(0.1,0.1){\special{pdf: epdf width 4in height
1.067in(htmlout.pdf)}}
\end{picture}
\end{center}
\caption{Example generated html file and its pop-up window}\label{
html_out}
\end{figure}

\subsection{Things to tweak}

Aside from adding some explanatory text to the created html file,
there are really only 2 things you'll likely need to tweak in the file:

\subsubsection{The title}

You can change the title of the main window by altering
\begin{verbatim}
<title>Neural Net</title>
\end{verbatim}
which is very close to the top of the file.  Just change `Neural Net'
to anything you like.

You should only need to change 1 thing in the javascript portion of
the file--the part starting with
\begin{verbatim}
<script language = "JavaScript"><!--
\end{verbatim}
and ending with
\begin{verbatim}
//--></script>
\end{verbatim}
which, you should notice, is a very large part of your generated html
file.

You can change the title of the pop-up output window in the line:
\begin{verbatim}
OutputWindow.document.write("<html><head><title>Neural Net</title></
head>");
\end{verbatim}
Just change `Neural Net' in that line to anything you like.

\subsubsection{Default values}

The Perl program will take a stab at default values by using the
minimum values in your original dataset reconstructed by using the
scaling factors file. For populated variates, it will give sensible
defaults.  For unpopulated variates, it will let everything be
unknown, which you will likely not want.  In either case, you'll want
to be able to change the default values.  To do that, in the form
section of the generated html file (the part which starts with
$<$form$>$,) for numerical variables just change the number in
``value=$''\ldots''$'', and for binary and categorical variables, just
move the word ``CHECKED'' to your default variable.

\section{neuron2palm.pl}\index{Utility scripts!Generate Palm code}

This Perl program is a wizard which creates code for a Palm project
that runs your trained neural network.  It takes the same files as
{\em neuron2html.pl}, so before you use this program, read the
\hyperref[neuron2html.pl]{section on converting your trained neural
network to html/Javascript}.

This wizard assumes that you have the Palm OS Developer Suite installed
on your computer.  If you don't, detailed instructions on how to install
it may be found in the {\bf readme.html} file on your neUROn2++ CD.

3 files will be created: a C++ Palm program (AppMain.cpp), its header
file (AppResources.h) and its resource file (AppResources.xrd).  These
files are then copied into a template Palm OS Developer Suite C++ project
that you have created, and then you modify the resource file (the Palm
GUI) with the Palm OS Resource Editor and make your project with the Palm
OS Developer suite.

\subsection{Making your Palm OS project}

Follow these steps to make your Palm OS project:

\subsubsection{Make a blank Palm OS Developer Suite project}\label{
blankproj}

First, go to \href{https://www.developerpavilion.com/palmos/}
{https://www.developerpavilion.com/palmos/}, and select the link
``Manage My Applications'', then ``Get a Creator ID'', then register
a unique 4 character ID for your application.

\begin{enumerate}
\item Open the Palm OS Developer Suite.  Choose menu {\bf Window}
      $\rightarrow$ {\bf Open Perspective} $\rightarrow$
      {\bf Palm OS C/C++ Development}
\item Choose menu {\bf File} $\rightarrow$ {\bf New} $\rightarrow$
      {\bf Project}
\item Choose ``Managed Make 68K C/C++ Project'' under Wizard:
      Palm OS Development
\item Select ``Next'', and type in the name of your project in the box
\item Select ``Next'', and in the next window, enter your unique creator
      ID in the ``Creator ID'' box.  {\bf \em Do not use STRT!}
\item Accept the other defaults, and press ``Next''.  Choose ``Simple
      Application'', and then press ``Next''.  Press ``Finish''.
\end{enumerate}

\subsection{Run neuron2palm.pl}

Enter the same files, with the same descriptions, as
\hyperref[neuron2html.pl]{neuron2html.pl.  Read that section if you
haven't already.}

\subsection{Copy created files to the project directories}

Open the ``src'' tree under your project in the left hand pane of the
Palm OS Developer Suite window.  Right click on ``AppMain.c'' and
choose ``Delete''.

Copy the {\em neuron2palm.pl} created files {\bf AppMain.cpp} and
{\bf AppResources.h} to your project's ``src'' directory, and
{\bf AppResources.xrd} to your project's ``rsc'' directory, overwriting
the files there.  Copy the files in {\em utilities/perl/Craig's Palm icons}
on the neUROn2++ CD into your project's ``rsc/bitmaps'' directory.

\subsection{Tweak your GUI}

Open the Palm OS Resource Editor program.  This program is separate
from the Palm OS Developer Suite. Select menu choice {\bf File}
$\rightarrow$ {\bf Open$\ldots$}, and find your resource file
{\bf AppResources.xrd} in your project's ``rsc'' directory.  Move around
the elements on {\bf Form 1001 (Main)} to your tastes, and update the
reference on {\bf Form 1007 (About)}. Select menu choice {\bf File}
$\rightarrow$ {\bf Save}.

\subsection{Make your project}

Return to the Palm OS Developer Suite.  Right click on your project name
in the left hand pane of the window, and choose ``Refresh''.  Right
click on your project name again, and choose ``Rebuild Project''.

\subsection{Test your project}

Open the Palm OS Emulator program.  This program is separate
from the Palm OS Developer Suite.  Right click on the emulator, choose
``New'', and enter the name and location of your ROM file.  Remember
to soft reset the emulator (also by right clicking on it) between each
run!  Go back to the Palm OS Developer Suite, and select menu choice
{\bf Run} $\rightarrow$ {\bf Run}.  Click on the ``New'' button, and
the name of your Palm program should automatically appear under
{\em Palm OS Application} in the left hand window, and a checked box
next to the name of your application in the right hand pane under the
{\em Project} tab.  Select the {\em Target} tag, and choose {\bf Palm
OS Emulator} for {\em Device}.  Click on the ``Apply'' button, then
the ``Run'' button.

\subsubsection{Note}

Note that if you change your numeric defaults (and you probably
should, as they are encoded as the minima for those variates,) you
need to change the C++ code {\bf in 2 places}.  The first place is
near the top of the C++ program, where the {\em \_input} array is
initialized--just change the values of that array.  It's immediately
below the lines:

\begin{verbatim}
// Initialize INPUT vector !!! Change these values !!!
// Remember to make sure that for binary variables, the initial
//    value is set in the GUI, and for categorical variables,
//    make sure the initial value is the first category.
\end{verbatim}

The second place where you need to change the value of your numeric
variates is in the {\em MainFormHandleEvent} function, and these lines
are preceded by the comment:

\begin{verbatim}
// !!! Change to match your ... field name and default !!!
\end{verbatim}

where $\ldots$ is replaced by your numeric field name.  Immmediately
below these lines, you'll find a {\em WriteToField} function call,
just change the last value inside of the quotation marks to match what
you entered in the {\em \_input} array.

\section{neuron2iphone.rb}\index{Utility scripts!Generate iPhone code}

This Ruby/Tk program is a wizard which creates code for a iPhone/iPod Touch
project that runs your trained neural network.  You need to have Ruby installed
with the Tk bindings on your computer.  Fortunately, modern MacOS installations
such as Leopard come with Ruby and the Tk bindings, and since you {\em have to}
use Xcode to make iPhone/iPod Touch applications, you'll be able to run {\em neuron2iphone.rb}
without difficulty on the same platform that you'll be developing your iPhone/iPod Touch
application.  However, if you are determined to run {\em neuron2iphone.rb} on a different
operating system than the one on which you'll be developing iPhone/iPod Touch apps,
just be sure your Ruby installation includes the Tk bindings.

{\em neuron2iphone.rb} takes the same files as {\em neuron2html.pl}, so before you
use the program, read the \hyperref[neuron2html.pl]{section on converting your
trained neural network to html/Javascript}. {\em Note: the wizard is currently only
implemented for SimpleProp neural networks and Logistic Regression models.}

Three files will be created by {\em neuron2iphone.rb}, the ViewController .h file, the
.m file, and a detailed instruction file explaining step-by-step how to create an
iPhone/iPod Touch application with Xcode.  As the only way to deploy your application
is using the iTunes App Store and getting your app into the store is not trivial,
the iPhone Developer instructions for submitting an application are included 
in the {\em scripts} directory of the neUROn distro.  Running {\em neuron2iphone.rb}
is fairly easy and self explanatory as it has a graphical user interface.

\chapter{Old File Structure}\index{History!neUROn++|(}

What follows is a description of the structure of neUROn++ files prior
to the design of neUROn2++.

\section{Files specifying input data}

The following files specify data to be modeled:

\subsection{The input file ({\em filename}) }

In neUROn++, the input file is specified by each exemplar as a row
delimited by spaces and truncated by a carriage return, containing a
list of the input variables followed by the output variable, e.g.:

\begin{center}
\begin{tabular}{cccc}
$i_{11}$ & $i_{12}$ & $\ldots$ & $o_1$ \\
$i_{21}$ & $i_{22}$ & $\ldots$ & $o_2$ \\
$i_{31}$ & $i_{32}$ & $\ldots$ & $o_3$ \\
$\ldots$ & $\ldots$ & $\ldots$ & $\ldots$ \\
$i_{n1}$ & $i_{n2}$ & $\ldots$ & $o_n$ \\
\end{tabular}
\end{center}

For example, the exclusive-or problem is encoded:

\begin{center}
\begin{tabular}{ccc}
0 & 0 & 0 \\
0 & 1 & 1 \\
1 & 0 & 1 \\
1 & 1 & 0 \\
\end{tabular}
\end{center}

It is the job of the program to count the number of rows as the number
of exemplars.  The program is thus required to identify blank lines
and discard them (as some word processing programs, for example,
append blank lines to the end of a file without notifying the user.)

\subsection{The raw training file ({\em filename}\_train.txt)
}\label{rawFile}

The raw training file is generated by the ``big lotto'' module, which
randomizes the raw data set into training and test sets, with a
specified number of exemplars in each, and insures that, while random,
the frequency of outcomes in the test and training sets are similar.

Format of the raw training file is (space delimited) number of
exemplars, number of input nodes, number of output nodes, threshold
value\footnote{threshold $T$ is for discrete output values: if network
guess $>= T$, output is 1, if $< T$, output is 0} followed by lines
containing each exemplar as specified in the preceeding section, e.g.:

\begin{center}
\begin{tabular}{cccc}
{\em $N_x$ exemplars} & {\em $N_o$ outputs} & {\em $N_i
$ inputs} & {\em threshold T} \\
\end{tabular}

\begin{tabular}{cccc}
$i_{11}$ & $i_{12}$ & $\ldots$ & $o_1$ \\
$i_{21}$ & $i_{22}$ & $\ldots$ & $o_2$ \\
$i_{31}$ & $i_{32}$ & $\ldots$ & $o_3$ \\
$\ldots$ & $\ldots$ & $\ldots$ & $\ldots$ \\
$i_{n1}$ & $i_{n2}$ & $\ldots$ & $o_n$ \\
\end{tabular}
\end{center}

An example of an exclusive-or raw training file is:

\begin{verbatim}
75  1  2  0.5
1  0  1
0  1  1
0  1  1
1  0  1
0  1  1
1  0  1
0  1  1
1  0  1
0  1  1
1  0  1
\end{verbatim}
$\ldots$

\subsection{The raw test file ({\em filename}\_test.txt) }

The raw test file is also generated by the ``big lotto'' module, and
conforms to the specification of the raw training file, while
containing the test data rather than the training data.

\subsection{training\_set}

This file contains the number of examples in the training set,
followed by the normalized training set.  If $lbound$ is the lower
boundary of a normalized value, $hbound$ is the upper boundary, $min$
is the minimum value of a variable, $max$ is the maximum value of a
variable, and $x$ is the variable's value, then the normalized value
of $x$, $x_{norm}$ is:

\begin{equation}\label{inputScale}
x_{norm} = lbound + [({{x - min} \over {max - min}}) \times (hbound -
lbound)]
\end{equation}

Note that $x_{norm} = lbound$ when $x = min$ and $x_{norm} = hbound$
when $x = max$.

Note also that outputs are normalized (to the output bounds specified
in the header file) if ``discrete output'' is specified {\em off} in
the header file, but set to 0 or 1 if ``discrete output'' is specified
{\em on} in the header file.

An example of an exclusive-or raw training\_set file is:

\begin{verbatim}
75
 9.000000e-01 -9.000000e-01 1.000000e+00
-9.000000e-01  9.000000e-01 1.000000e+00
-9.000000e-01  9.000000e-01 1.000000e+00
 9.000000e-01 -9.000000e-01 1.000000e+00
-9.000000e-01  9.000000e-01 1.000000e+00
 9.000000e-01 -9.000000e-01 1.000000e+00
-9.000000e-01  9.000000e-01 1.000000e+00
 9.000000e-01 -9.000000e-01 1.000000e+00
-9.000000e-01  9.000000e-01 1.000000e+00
 9.000000e-01 -9.000000e-01 1.000000e+00
-9.000000e-01  9.000000e-01 1.000000e+00
\end{verbatim}
$\ldots$

\subsection{test\_set}

Like training\_set, test\_set is derived from {\em
filename}\_test.txt, and conforms to the specification of
training\_set, while containing the test data rather than the training
data.

\section{Files specifying network architecture}

The following files pertain to network architecture:

\subsection{minimax}

The file minimax contains the minimum and maximum values for each
variable, and is generated by the prepare module.  Its format is 2
lines, with the first line the mininum values, and the second the
maximum values.  The final value on each line is the output variable.
Each value is separated by blank space, with the lines truncated by
carriage returns.  File format is thus:

\begin{center}
\begin{tabular}{cccccc}
$i1_{minimum}$ & $i2_{minimum}$ & $i3_{minimum}$ & $\ldots$ &
$in_{minimum}$ & $o_{minimum}$
\\ $i1_{maximum}$ & $i2_{maximum}$ & $i3_{maximum}$ & $\ldots$ &
$in_{maximum}$ & $o_{maximum}$ \\
\end{tabular}
\end{center}

For example, the minimax file for the exclusive-or problem is:

\begin{verbatim}
0.000000e+00 0.000000e+00 0.000000e+00
1.000000e+00 1.000000e+00 1.000000e+00
\end{verbatim}

\begin{center}
\begin{tabular}{cccccccc}
$i1_{minimum}$ & $i2_{minimum}$ & $\ldots$ & $in_{minimum}$ &
$o1_{minimum}$ & $o2_{minimum}$ & $\ldots$ & $om_{minimum}$ \\
$i1_{maximum}$ & $i2_{maximum}$ & $\ldots$ & $in_{maximum}$ &
$o1_{maximum}$ & $o2_{maximum}$ & $\ldots$ & $om_{maximum}$ \\
\end{tabular}
\end{center}

Where the number of input nodes and the number of output nodes is
specified in the header file.

\subsection{hidden\_weights}

The file hidden\_weights contains the hidden weight matrix, in the
form of multiple lines separated by carriage returns, with each line
representing one hidden node, and each line containing the input
weights for that node.  All weights are bounded by bounds specified in
the header file.  File format is thus:

\begin{center}
\begin{tabular}{ccccc}
$w_{h_1i_1}$ & $w_{h_1i_2}$ & $w_{h_1i_3}$ & $\ldots$ & $w_{h_1i_n}$
\\
$w_{h_2i_1}$ & $w_{h_2i_2}$ & $w_{h_2i_3}$ & $\ldots$ & $w_{h_2i_n}$
\\
$w_{h_3i_1}$ & $w_{h_3i_2}$ & $w_{h_3i_3}$ & $\ldots$ & $w_{h_3i_n}$
\\
$\ldots$ & $\ldots$ & $\ldots$ & $\ldots$ & $\ldots$ \\
$w_{h_mi_1}$ & $w_{h_mi_2}$ & $w_{h_mi_3}$ & $\ldots$ & $w_{h_mi_n}$
\\
\end{tabular}
\end{center}

This is a file of critical importance, as after the program is run, it
represents a part of the modeling solution.

An example is from a 2 input network with 3 hidden nodes:

\begin{verbatim}
1.022590e-01 -4.725589e-01
4.030996e-01 -4.004863e-01
-1.634707e-01 2.717012e-01
\end{verbatim}

\subsection{hidden\_weights.full}

This file is generated in neUROn++ for the ``black input'' feature,
where input nodes are held to 0, and the network is retrained.  In
this implementation, the hidden\_weights file is copied to
hidden\_weights.full, and then subsequent hidden\_weights files
represent full training of feature deficient (e.g. specific input
variables held to 0) subnetworks.  It has the same format as the
hidden\_weights file.

This type of file structure appeared early on in neUROn++, when single
input variables were commonly held to 0, and the network retrained.
Subsequently, neUROn++ was reprogrammed so that input variables were
held to 0 in a combinatorial fashion or in a stepwise regression
algorithm.

\subsection{hidden\_bias}

This file contains one line, with the weight values of the bias nodes
on the hidden layer separated by spaces.  It thus takes the format:

\begin{center}
\begin{tabular}{ccccc}
$w_{h_1b}$ & $w_{h_2b}$ & $w_{h_3b}$ & $\ldots$ & $w_{h_mb}$ \\
\end{tabular}
\end{center}

\subsection{hidden\_bias.full}

Like hidden\_weights.full, hidden\_bias.full was designed for the
``black input feature'' in the initial stages of neUROn++.  It has the
same utility as the hidden\_bias.full file.

\subsection{output\_weights}

Like hidden\_weights, output\_weights contains the output weight
matrix.  It is a single line (for a single output node) containing the
output weights separated by spaces, and takes the form:

\begin{center}
\begin{tabular}{ccccc}
$w_{o_1}$ & $w_{o_2}$ & $w_{o_3}$ & $\ldots$ & $w_{o_n}$ \\
\end{tabular}
\end{center}

\subsection{output\_weights.full}

Like hidden\_weights.full, output\_weights.full was designed for the
``black input feature'' in the initial stages of neUROn++.  It has the
same utility for the output layer as the hidden\_weights.full file
does for the hidden layer.

\subsection{output\_bias}

The file output\_bias contains one value (for a single output node) as
the weight of the bias on the output node.

\subsection{output\_bias.full}

Like hidden\_bias.full, output\_bias.full was designed for the ``black
input feature'' in the initial stages of neUROn++.  It has the same
utility for the output layer as the hidden\_bias.full file does for
the hidden layer.

\section{Training multiple subnetworks}

Using the ``black input'' feature, subnetworks are generated and
trained with individual or sets of input nodes held to 0.  These
trained subnetworks are then statistically compared to the full
network to determine the significance of the input node(s) that
is(are) ``blacked out''.  The method to specify generation of such
subnetworks is either automatic, as in stepwise regression analysis,
or
specified by the user, as in combinatorial analysis.  In both, it is
necessary to define a specific feature (or variable) to be ``blacked''
out as a set of the input nodes which characterize it.  For example,
input node 17 may represent patient age, and input node 16 represents
if patient age is available in the database (not available is 0,
available is 1).  In this example, both nodes 16 {\em and} 17 must be
simultaneously held to 0 to ``black out'' the feature.  The file
class\_define specifies how input nodes are grouped into features to
be ``blacked out''.

The file format of the class\_define file is as follows:

\begin{quote}
\begin{em}
feature$_1$, feature$_2$,..., feature$_N$
\end{em}
\end{quote}

Where ``feature'' is a number, multiple numbers with a plus (+) in
between indicating ``and'', or two numbers with a dash (-) in between,
indicating ``through''.

For example:
\begin{verbatim}
1-4,5,6-10,11,12,13+15,14
\end{verbatim}

indicates that there are 15 input nodes divided into 7 features.  The
first feature is input nodes 1 through 4, the second feature is node
5, the 3rd feature is input nodes 6 through 10, the 4th is node 11,
the 5th is node 12, the 6th is nodes 13 and 15, and the 7th is node
14.

\section{Reports}

The following files contain reports generated during training and
testing:

\subsection{error\_file}\label{errorFile}

The file error\_file is the master record file.  It should contain
{\em all} of the information present in the header file organized in a
readable fashion, and all calculated errors at specified intervals.
It is used as documentation for an experiment that consists of fully
training a neural network.

Example:

\begin{verbatim}
Mon Jul 10 14:12:36  2000
=============================
Neural network error file:

  Target error not available. AUTO_ERROR ON. ERROR_WINDOW_WIDTH 1000
  STOP_AT_MIN ON
  STOP_INCREASE ON
  MOMENTUM OFF
  WEIGHT DECAY OFF
  2 INPUTS; 1 OUTPUT
Number of hidden nodes on the 1 hidden layers are:
	3
  BIAS ON
  Transfer function output low bound 0.100000, high bound 0.900000
    Input low bound -0.900000, high bound 0.900000
  DISCRETE_OUTPUT ON
  XENTROPY ON
  BATCH/EPOCH LEARNING DISABLED
  BLACK_INPUT OFF

	  (degree of freedom is 13)

75 examples in training set

N         Set Error     XE Error      CErr Train  CErr Test

       1  1.282783e-01  7.064047e-01  50.666667   48.000000
       2  1.274456e-01  7.030068e-01  50.666667   48.000000
       3  1.269191e-01  7.008692e-01  50.666667   48.000000
       4  1.265779e-01  6.994885e-01  50.666667   48.000000
       5  1.263515e-01  6.985742e-01  50.666667   48.000000
       6  1.261979e-01  6.979550e-01  50.666667   48.000000
       7  1.260917e-01  6.975270e-01  50.666667   48.000000
       8  1.260168e-01  6.972256e-01  50.666667   48.000000
       9  1.259631e-01  6.970095e-01  50.666667   48.000000
      10  1.259239e-01  6.968520e-01  50.666667   48.000000
      20  1.257945e-01  6.963322e-01  50.666667   48.000000
      30  1.257459e-01  6.961384e-01  50.666667   48.000000
      40  1.257018e-01  6.959635e-01  50.666667   48.000000
      50  1.256583e-01  6.957919e-01  50.666667   48.000000
      60  1.256143e-01  6.956188e-01  50.666667   48.000000
      70  1.255687e-01  6.954400e-01  50.666667   48.000000
      80  1.255205e-01  6.952518e-01  50.666667   48.000000
      90  1.254686e-01  6.950500e-01  50.666667   48.000000
     100  1.254120e-01  6.948302e-01  50.666667   48.000000
     200  1.240699e-01  6.896549e-01  24.000000   28.000000
     300  1.130212e-01  6.450682e-01  24.000000   28.000000
     400  9.205705e-02  5.416243e-01  24.000000   28.000000
     500  4.299956e-02  3.227091e-01  0.000000   0.000000
     600  4.685616e-03  9.736120e-02  0.000000   0.000000
     700  1.376876e-03  5.197709e-02  0.000000   0.000000
     800  6.277901e-04  3.487233e-02  0.000000   0.000000
     900  3.535863e-04  2.607810e-02  0.000000   0.000000
    1000  2.251775e-04  2.076367e-02  0.000000   0.000000
    2000  2.373450e-05  6.694852e-03  0.000000   0.000000
    3000  8.330674e-06  3.959491e-03  0.000000   0.000000
    4000  4.188223e-06  2.805047e-03  0.000000   0.000000

!! Error changing decreased under MIN_DELTA_ERR=1.000000e-03


To reach an overall Xentropy error of 2.638119e-03 required 4218
iterations.

----------------------------------------------------------------------
-------
For the TEST SET after iteration 4217:
--------------------------------------
true positive rate = 12    true negative rate = 13
false positive rate = 0    false negative rate = 0

sensitivity = 1.000000                    specificity = 1.000000
positive predictive value = 1.000000      negative predictive value =
1.000000
----------------------------------------------------------------------
---------

Training time (HHHH:MM:SS) 0000:00:04=============================
Mon Jul 10 14:12:42  2000

\end{verbatim}

\subsection{black\_report}

This file has an identical format to error\_file, and is generated
when an input node or set of input nodes is held to 0, as in the
``black input'' feature.

\subsection{xentropy.full}\index{Error function!Cross-entropy}

This file contains the cross-entropy error calculated for a fully
trained network.  Like many features, it was designed early in
neUROn++'s life cycle, prior to combinatorial or stepwise regression
analysis.

\subsection{out-file}\label{outFile}

The file outfile contains one line with the number of examples in the
test set, followed by the number of output nodes, followed by the
threshold value.  Subsequent lines are pairs of the known value
followed by the network's guess.  For a one-output network, this takes
the form:

\begin{center}
\begin{tabular}{ccc}
{\em $N_x$ exemplars} & {\em $N_o$ outputs} & {\em threshold T} \\
\end{tabular}

\begin{tabular}{cc}
$known_1$ & $guess_1$ \\
$known_2$ & $guess_2$ \\
$known_3$ & $guess_3$ \\
$\ldots$ & $\ldots$ \\
$known_n$ & $guess_n$ \\
\end{tabular}
\end{center}

For example,

\begin{verbatim}
25 1 0.500000
1.000000 0.996711
1.000000 0.998407
1.000000 0.998407
1.000000 0.996711
1.000000 0.996711
1.000000 0.998407
1.000000 0.996711
\end{verbatim}
$\ldots$

\chapter{The Old Header File}

Prior to full C++ implementation in neUROn2++, as much of the network
description as possible was encoded in pre-processor definitions in a
header file, so that the compiled executable is as tailored as
possible to the data model.  The following section details these
pre-processor definitions.

\section{Global defines}

\subsection{Limits of computability}

The TOLERANCE preprocessor definition specifies the precision of the
compiler, e.g. 1e-8, 1e-12, etc.

\section{Neurocomputer architecture}

The following pre-processor definitions encode neurocomputer
architecture.

\subsection{BIAS}

Bias nodes are nodes held to 1 on a layer.  In neUROn++, BIAS ON
indicates bias nodes on all layers, whereas BIAS OFF indicates bias
nodes on no layers.

\begin{verbatim}
#define BIAS {1, 3, 4}
\end{verbatim}

would specify bias nodes on layers 1, 3 and 4.

\subsection{Layers}

In neUROn++, up to 3 hidden layers could be specified.  The following
pre-processor definitions describe layer architecture:

\begin{itemize}
\item NUMBER\_NODES\_MAX: specifies the maximum number of nodes on
      each layer, and is used for initializing arrays
\item NUMBER\_HDLAYER: specifies the number of hidden layers
\item NUMBER\_HIDDEN: specifies the number of hidden nodes in the 1st
      layer
\item NUMBER\_HIDDEN1: specifies the number of hidden nodes in the 2nd
      layer
\item NUMBER\_HIDDEN2: specifies the number of hidden nodes in the 3rd
      layer
\item NUMBER\_OUTPUT: specifies the number of output nodes
\item NUMBER\_INPUT: specifies the number of input nodes
\end{itemize}

\subsection{Activation function}

In neUROn++, the activation function is defined in the header file, as
is its derivative for the backpropagation algorithm.  The activation
function and its derivative are defined in macros:
\index{Backpropagation}

\begin{itemize}
\item TRANSFER\_FUNCTION(x)
\item DIFF\_TRANSFER(o)
\end{itemize}

Because backpropagation uses the sigmoidal activation function:
\index{Transfer function!Sigmoidal}

\begin{equation}
t_k = {1 \over {1 + e^{-net_k}}}
\end{equation}

and its derivative:

\begin{equation}
f\prime(net_k) = o_k(1-o_k)
\end{equation}

The following macro definitions are:

\begin{verbatim}
#define TRANSFER_FUNCTION(x) 1.0 / (1.0 + exp( -(x) ))
#define DIFF_TRANSFER(o) (o) * ( 1 - (o) )
\end{verbatim}

\section{Preparation}

The following pre-processor definitions describe data and model
processing prior to modeling.

\subsection{Data set}

The pre-processor definition TRAINING\_MAX specifies the maximum
number of exemplars in the training data set.

\subsection{Scaling}

\subsubsection{Input values}

As described in equation \ref{inputScale}, input values are scaled
according to the following formula:

\begin{displaymath}
x_{norm} = lbound + [({{x - min} \over {max - min}}) \times (hbound -
lbound)]
\end{displaymath}

Where $lbound$ is the lower boundary of the input value to be scaled,
$hbound$ is the upper boundary, $min$ is the minimum value of the
variable, $max$ is the maximum value of the variable, and $x$ is the
variable's value.  Thus, $lbound$ and $hbound$ for input variables are
specified in the header file as:

\begin{itemize}
\item INPUT\_LBOUND = $lbound$
\item INPUT\_HBOUND = $hbound$
\end{itemize}

\subsubsection{Output values}

In neUROn++, {\em nondiscrete} (e.g. the pre-processor definition
DISCRETE\_OUTPUT is set OFF) output values were scaled using:

\begin{itemize}
\item OUTPUT\_LBOUND
\item OUTPUT\_HBOUND
\end{itemize}

where OUTPUT\_LBOUND and OUTPUT\_HBOUND represented the low and high
bounds respectively of the output value.

\subsection{Randomizing initial weights}

Initial weights are randomized to a value bounded by $-bound$ to
$+bound$.  This $bound$ is specified by the pre-processor definition
RANDOM\_THRESHOLD, which is generally set to 0.5.  Thus, weights are
generally initialized to values between $-0.5$ and $+0.5$.

\section{Training control}

The following pre-processor definitions describe global control of
training, including report generation and how training ends.

\subsection{Reporting}

The following pre-processor definitions describe report generation:

\begin{itemize}
\item PRINT\_COUNT
\item WITH\_CLASS
\item OUTCOMES\_FILE
\item TEST\_SET\_STATS
\end{itemize}

PRINT\_COUNT specifies at which iteration error outcomes are reported.
It may take the form of a number, e.g. a PRINT\_COUNT of 1000 would
specify error reportage every 1000 iterations, or an equation, e.g.:

\begin{verbatim}
#define PRINT_COUNT (int) ( pow(10, (int) log10(set_count) ) )
\end{verbatim}

Specifies a highly useful exponential count.

WITH\_CLASS specifies if classification error in the training and test
sets are reported at PRINT\_COUNT.  These are useful to generate
graphs of neurocomputer training and identify such behavior as
overlearning.  Classification error is described as:

\begin{equation}
CE = {N_{correct} \over N_{total}} * 100\%
\end{equation}

Where $CE$ is the classification error (in either training or test
sets), $N_{correct}$ is the number of correct guesses, and $N_{total}$
is the total number of guesses.

If OUTCOMES\_FILE is specified ON, ``out-file'' is generated as
specified in section \ref{outFile}.

If TEST\_SET\_STATS is specified ON, specificity, positive predictive
value, and negative predictive value are reported in ``error\_file''
as specified in section \ref{errorFile}.

\subsection{Ending training}

The following pre-processor definitions control how training ends.

\begin{itemize}
\item ITERATIONS\_MAX specifies the maximum number of iterations
      at which training ends.  Even if the error is not at a minimum,
      training will end at the number of iterations specified by this
      pre-processor definition.

\item AUTO\_ERROR controls the next 3 definitions.  If AUTO\_ERROR is
      OFF, {\em all} of these definitions will be ignored.  If
      AUTO\_ERROR is ON, {\em any} any of these definitions is
      sufficient to end training.

\begin{itemize}
\item ERROR\_WINDOW\_WIDTH defines the number of iterations at which
      training will end if either the STOP\_INCREASE or the
      STOP\_AT\_MIN conditions are satisfied.

\item STOP\_INCREASE specifies that training will end if the error at
      any iteration $+$ ERROR\_WINDOW\_WIDTH is greater than the error
      at that iteration, e.g.:

	\begin{displaymath}
	E_{n+\Delta} > E_{n}
	\end{displaymath}
	where $n$ is the iteration, $E$ is the error, and $\Delta$ the
	ERROR\_WINDOW\_WIDTH.

\item If STOP\_AT\_MIN is ON, training will end if the change in error
      over the number of iterations specified by ERROR\_WINDOW\_WIDTH
      is less than MIN\_DELTA\_ERR, and

\item MIN\_DELTA\_ERR specifies the minimum change in error at which
      training will end over any ERROR\_WINDOW\_WIDTH.  Thus, training
      will end if

	\begin{displaymath}
	 E_{\Delta min} > E_{n+\Delta} - E_{n}
	\end{displaymath}

	where $n$ is the iteration, $E$ is the error, $\Delta$ is
	ERROR\_WINDOW\_WIDTH, and $E_{\Delta min}$ the
	MIN\_DELTA\_ERR.

\end{itemize}
\end{itemize}

\section{Training}

The following pre-processor definitions specify how weights are
updated within a neural network, generally referred to as
``training''.

\subsection{Global definitions}

The following pre-processor definitions apply to any neural network
training methodology.

\subsubsection{Batch/Epoch}\index{Backpropagation!Off-line}

``Online'' learning refers to the scheme whereby weights are updated
with each training exemplar.\index{Backpropagation!On-line} Thus, if
$i$ refers to an individual training exemplar, then network weights
$w_i$ are updated by:

\begin{displaymath}
w_i \leftarrow w_i + \Delta w_i
\end{displaymath}

for each training exemplar, where $\Delta w_i
$ is the weight update value.

``Offline'' learning refers to the scheme whereby the weight updates
are calculated for the entire training set, summed, and applied to the
network's weights {\em after} the entire training set is processed:
\index{Backpropagation!Off-line}

\begin{displaymath}
w_i \leftarrow w_i + \sum \Delta w_i
\end{displaymath}

``Offline'' learning is often referred to as either ``Batch'' or ``
Epoch'' learning.\index{Backpropagation!Off-line}  If the
pre-processor variable BATCH\_EPOCH is ON, offline learning is
specified, if BATCH\_EPOCH is OFF, online learning is specified.
\index{Backpropagation!On-line}

\subsubsection{Discrete output}

If the pre-processor definition DISCRETE\_OUTPUT is ON, then output(s)
are defined to be 0 if the output value is less than a threshold
value, and 1 if the output value is greater than or equal to the
threshold value:

\begin{displaymath}
O \leftarrow O \geq T |^{1}_{0}
\end{displaymath}

where $O$ is the output value and $T$ is the threshold.  In neUROn++,
the threshold $T$ is specified in the raw training file as described
in section \ref{rawFile}.

\subsection{Backpropagation}\index{Backpropagation}

Backpropagation is specified by the pre-processor definition
CLASSIC\_BP.  Treatment of the backpropagation algorithm may be found
in a variety of texts such as Timothy Masters ``Practical Neural
Network Recipes in C++'' (Academic Press 525 B St, Suite 1900, San
Diego CA 92101-4495 www.apnet.com, www.academicpress.com ISBN
0-12-479040-2) 1993.

Briefly, canonical backpropagation proceeds as follows.  The input
vector is applied to the network, and outputs of the first hidden
layer are calculated by the activation function $t_k$ which is set to
be the the sigmoidal:\index{Transfer function!Sigmoidal}

\begin{equation}
t_k = {1 \over {1 + e^{-net_k}}}
\end{equation}

The outputs node values of the first hidden layer are applied to the
inputs nodes of the second hidden layer, etc.\ until the output
layer's
output node values are calculated.

The output node's weights are then updated by:

\begin{equation}
w_{jk} \leftarrow w_{jk} + \eta\delta_ko_j
\end{equation}

for learning rate $\eta$, $j$ output units and $k$ hidden units, where

\begin{equation}
\delta_k = (t_k - o_k)f\prime(net_k)
\end{equation}

$o_k$ is the known outcome value in the training set, and the the
derivative $f\prime(net_k)$ is:

\begin{equation}
f\prime(net_k) = o_k(1-o_k)
\end{equation}

Weights are then sequentially updated in the hidden layers in the
direction of {\em from} the ouput layer {\em to} the input layer,
hence the term ``backpropagation''.

\subsubsection{Learning rate}

The learning rate $\eta$ is specified by the preprocessor definition
ETA.

\subsubsection{Momentum}

Momentum adds a fraction of the previous iteration's weight update to
the current weight update:

\begin{equation}
\Delta w_t \leftarrow \Delta w_t + \alpha \Delta w_{t-1}
\end{equation}

where $\Delta w_t$ refers to the current weight update, and $\Delta
w_{t-1}$ to the previous weight update.  The ``momentum'' term is
$\alpha$.  The weight update rule in backpropagation then becomes:

\begin{equation}
\Delta w_t \leftarrow \eta\delta x + \alpha \Delta w_{t-1}
\end{equation}

The following pre-processor definitions specify momentum:
\begin{itemize}
\item MOMENTUM, if ON specifies that a momentum term will be applied.
\item ALPHA is the value of the momentum term.
\end{itemize}

\subsubsection{Weight decay}\index{Weight decay}

Weight decay subtracts a fraction of the previous iteration's weight
(not weight update as in momentum) from the current weight update:

\begin{equation}
\Delta w_t \leftarrow \Delta w_t - d w_{t-1}
\end{equation}

where $\Delta w_t$ refers to the current weight update, and $w_{t-1}$
to the previous weight.  The ``decay'' term is $d$.  The weight update
rule in backpropagation then becomes:

\begin{equation}
\Delta w_t \leftarrow \eta\delta x - d w_{t-1}
\end{equation}

The following pre-processor definitions specify weight decay:
\begin{itemize}
\item WEIGHT\_DECAY, if ON specifies that a weight decay term will be
      applied.
\item DECAY is the value of the weight decay term.
\end{itemize}

\subsubsection{Cross-entropy}\index{Error function!Cross-entropy}

Cross-entropy is an error function applied {\em only at the output
layer.}  It is defined by:

\begin{equation}
E = - t\log(o)+(1-t)\log(1-o)
\end{equation}

where $t$ is an output node's value derived by the network,
and $o$ is the known value of that output in the data set.

If the pre-processor definition XENTROPY is ON, then the cross-entropy
error function is calculated at the output node.  {\em It should be
clear that it is a requirement for XENTROPY that DISCRETE\_OUTPUT is
ON.}

\subsection{Shanno's algorithm}

This serves as an excellent example of the importance of
documentation.  A previous graduate student coded Shanno's memoryless
algorithm from Richard Golden's ``Mathematical Methods for Neural
Network Analysis and Design'' (1996) in neUROn++.  Experiments
demonstrated that training using this method resulted in speed
improvements up to 2 orders of magnitude.  Not only was the algorithm
itself coded, however, a very elegant method of automatically
switching between backpropagation and Shanno's algorithm was also
encoded.\index{Backpropagation}

Unfortunately, because the code was so poorly documented, subsequent
code modifications failed to include this module.  And also, as a
result, I am unable to reconstruct the meaning of the pre-processor
definitions here, only to list them as they appear in neUROn++ header
file.

\begin{verbatim}
#define SHANNO     OFF        // If SHANNO is ON, XENTROPY must be ON!

#define FAIL_LIMIT  2         // ensure shanno's algo working under
control
#define SWITCH_IT  60         // number of iterations for BP training
#define SHOW_SWITCH OFF       // ON to see the engine switch--hurt
performance!
#define SHANNO_MIN_DERROR 1e-4 // control of sensitivity of Shanno
switch

#define ALPHA_SHANNO 0.0001
#define BETA_SHANNO  0.999    // They are used for decide proper step
\end{verbatim}

\section{Feature extraction}

In neUROn++, feature extraction using Wilk's GLRT was originally
encoded in combinatorial fashion, and specified in the header file.
Later, with virtually no additional overhead, stepwise regression was
added as a set of internal modules.  Because these modules call the
training routines, the efficiency of the training routines do not
suffer by their presence in the compiled executable, only that the
executable is somewhat larger as a result.  The combinatorial feature
extraction definitions encoded in the header file were:

\begin{itemize}
\item If DEFAULT\_BLACK was ON, all combinations of features were
      automatically extracted.  For example, if features f1, f2 and f3
      represented the input variables, the following feature deficient
      networks were statistically compared to the full network (f1,
      f2, f3): (f1, f2), (f1, f3), (f2, f3), (f1), (f2), and (f3).
\item If INTERACT\_BLACK was ON, specific combination subsets of
      features could be feature extracted.  Using the above example,
      just (f1, f2), (f1, f3), (f2, f3) could be compared to the full
      network (f1, f2, f3).
\item MAX\_CLASS specified the maximum number of features to be
      extracted (not necessarily the number of input nodes, as more
      than one input nodes commonly encode a single feature.)
\item MAX\_CASE specified the maximum number of possible combinations
      of features, or $2^{MAX\_CLASS}$.
\item D\_NARY $= MAX\_CLASS + 1$.
\end{itemize}
\index{History!neUROn++|)}

% Index bookmark
\special{pdf: out 0 << /Title (Index) /Dest [ @nextpage /FitB
@ypos ] >>}

% Index see cross-references
\index{Cross-entropy|see{Error function, Cross-entropy}}
\index{Transfer function|see{Forward propagation}}
\index{On-line|see{Backpropagation, On-line}}
\index{Off-line|see{Backpropagation, Off-line}}
\index{Batch/epoch|see{Backpropagation, Off-line}}

\printindex

\end{document}


























































































